% Generated mechanically from frozen input inputs/p12/ja/examples.tex.
% Do not edit this generated file; edit the bound locale source instead.

% OK, start here.
%

\setcounter{chapter}{109}
\chapter{例}



\phantomsection
\label{examples-section-phantom}


% BEGIN EXPANDED MODULE ch110_01.tex
\section{序論}
\label{examples-section-introduction}

\noindent
本章では，理論を照らし出す例を収める．





\section{空の極限}
\label{examples-section-empty-limit}

\noindent
この例は Waterhouse による；\cite{Waterhouse} を参照せよ．
$S$ を非可算集合とする．各有限部分集合
$T \subset S$ に対し，単射 $T \to \mathbf{N}$ 全体の集合を $M_T$ とする．
$T \subset T' \subset S$ が有限ならば，制限写像 $M_{T'} \to M_T$
は全射である．したがって，$S$ の有限部分集合からなる
有向半順序集合上に，全射な遷移写像をもつ逆系が得られる．
しかし $\lim M_T = \emptyset$ である．実際，極限の元があれば，
単射 $S \to \mathbf{N}$ が定まるからである．



\section{零極限}
\label{examples-section-zero-limit}

\noindent
$(S_i)_{i \in I}$ を，非空集合からなる有向逆系で，
遷移写像が全射かつ
$\lim S_i = \emptyset$ であるものとする；Section
\ref{examples-section-empty-limit} を参照せよ．
$K$ を体とし，
$$
V_i = \bigoplus\nolimits_{s \in S_i} K
$$
と置く．このとき，$i \geq j$ に対する遷移写像 $V_i \to V_j$ は全射である．
しかし，$\lim V_i = 0$ である．実際，$v = (v_i)$ が極限の元ならば，
$v_i$ の台は有限部分集合 $T_i \subset S_i$ となり，
$\lim T_i \not = \emptyset$ である；Categories, Lemma
\ref{categories-lemma-nonempty-limit} を参照せよ．

\medskip\noindent
各 $i$ に対し，すべての基底ベクトル $s \in S_i$ を $1$ に送る
一意な $K$-線形写像 $V_i \to K$ を考える．その核を
$W_i \subset V_i$ とする．このとき
$$
0 \to (W_i) \to (V_i) \to (K) \to 0
$$
は，有向集合 $I$ 上のベクトル空間の逆系の，分裂しない短完全列である．
したがって，
% P12-ERR-0315: see ../control/ERRATA_CANDIDATES.jsonl
$W_i$ は，全射な遷移写像をもち，極限が零で，
$R^1\lim$ が零でない $K$-ベクトル空間の有向系である．
% END EXPANDED MODULE ch110_01.tex
% BEGIN EXPANDED MODULE ch110_02.tex
\section{準コンパクト空間の非準コンパクトな逆極限}
\label{examples-section-lim-not-quasi-compact}

\noindent
$\mathbf{N}$ で自然数の集合を表す．
各整数 $n$ に対し，$I_n$ を $> n$ である自然数全体の集合とする．
$T_n$ を，$\{1\}, \ldots , \{n\}, I_n$ を基とする
$\mathbf{N}$ 上の一意な位相と定める．位相空間
$(\mathbf{N}, T_n)$ を $X_n$ と書く．各 $m < n$ に対して恒等写像
$$
f_{n, m} : X_n \longrightarrow X_m
$$
は連続である．明らかに $m < n < p$ ならば，合成
$f_{p, n} \circ f_{n, m}$ は $f_{p, m}$ に等しい．したがって
$((X_n), (f_{n,m}))$ は準コンパクト位相空間からなる有向逆系である．

\medskip\noindent
$T$ を $\mathbf{N}$ 上の離散位相とし，$X$ を
$(\mathbf{N}, T)$ とする．このとき各整数 $n$ に対して恒等写像
$$
f_n : X \longrightarrow X_n
$$
は連続である．これが上の有向逆系の逆極限であることを示そう．
任意の位相空間 $(Y, S)$ をとる．各整数 $n$ に対し，
$$
g_n : (Y, S) \longrightarrow (\mathbf{N}, T_n)
$$
を連続写像とする．すべての $m < n$ に対して
$f_{n,m} \circ g_n = g_m$，すなわち系 $(g_n)$ が上の有向逆系と
両立すると仮定する．特に，集合の写像 $g_n$ はすべて共通の写像
$$
g : Y \longrightarrow \mathbf{N}.
$$
に等しい．さらに各整数 $n$ に対し，$\{n\}$ は $X_n$ の開集合なので，
$g^{-1}(\{n\}) = g_n^{-1}(\{n\})$ も $Y$ の開集合である．
したがって集合の写像 $g$ は，$Y$ 上の位相 $S$ と
$\mathbf{N}$ 上の位相 $T$ に関して連続である．ゆえに
$(X, (f_n))$ は上の有向逆系の逆極限である．

\medskip\noindent
しかし $X$ は明らかに準コンパクトでない．実際，単点集合からなる
無限開被覆は逆極限をもたない．
% P12-ERR-0316: see ../control/ERRATA_CANDIDATES.jsonl

\begin{lemma}
\label{examples-lemma-lim-not-quasi-compact}
$\mathbf{N}$ 上の準コンパクト位相空間の逆系で，その極限が
準コンパクトでないものが存在する．
\end{lemma}

\begin{proof}
上の議論を参照せよ．
\end{proof}
% END EXPANDED MODULE ch110_02.tex
% BEGIN EXPANDED MODULE ch110_03.tex
\section{ファイバー積上の構造層}
\label{examples-section-silly}

\noindent
$X, Y, S, a, b, p, q, f$ を Derived Categories of Schemes,
Section \ref{perfect-section-kunneth} の序論におけるものとする．図式は
$$
\xymatrix{
& X \times_S Y \ar[ld]^p \ar[rd]_q \ar[dd]^f \\
X \ar[rd]_a & & Y \ar[ld]^b \\
& S
}
$$
である．このとき標準写像
$$
can :
p^{-1}\mathcal{O}_X \otimes_{f^{-1}\mathcal{O}_S} q^{-1}\mathcal{O}_Y
\longrightarrow
\mathcal{O}_{X \times_S Y}
$$
があるが，これは一般には同型でない．

\medskip\noindent
例えば $S = \Spec(\mathbf{R})$，$X = \Spec(\mathbf{C})$，
$Y = \Spec(\mathbf{C})$ とする．このとき
$X \times_S Y = \Spec(\mathbf{C}) \amalg \Spec(\mathbf{C})$
は二点からなる離散空間であり，層 $p^{-1}\mathcal{O}_X$，
$q^{-1}\mathcal{O}_Y$，$f^{-1}\mathcal{O}_S$ はそれぞれ
$\mathbf{C}$，$\mathbf{C}$，$\mathbf{R}$ に値をとる定数層である．
したがって $can$ の始域は，二点からなる離散空間上の
$\mathbf{C} \otimes_\mathbf{R} \mathbf{C}$ に値をとる定数層である．
ゆえにその大域切断は $\mathbf{R}$-ベクトル空間として次元 $8$ をもつが，
$can$ の終域の大域切断をとると
$\mathbf{C} \times \mathbf{C}$ が得られ，これは
$\mathbf{R}$-ベクトル空間として次元 $4$ をもつ．

\medskip\noindent
別の例を挙げる．$k$ を代数閉体とし，
$S = \Spec(k)$，$X = \mathbf{A}^1_k$，$Y = \mathbf{A}^1_k$ とする．
非空開集合 $U \subset X \times_S Y = \mathbf{A}^2_k$ に対し，像
$p(U) \subset X = \mathbf{A}^1_k$ と $q(U) \subset \mathbf{A}^1_k$
は開集合であり，読者は
$$
\left(
p^{-1}\mathcal{O}_X \otimes_{f^{-1}\mathcal{O}_S} q^{-1}\mathcal{O}_Y
\right)(U) = \mathcal{O}_X(p(U)) \otimes_k \mathcal{O}_Y(q(U))
$$
を示すことができる．$U$ が $X \times_S Y = \mathbf{A}^2_k$ における
既約曲線 $C$ の補集合で，$p|_C$ と $q|_C$ がともに非定数ならば，
これは $\mathcal{O}_{X \times_S Y}(U)$ に等しくない．

\medskip\noindent
一般の場合に戻り，$z = (x, y, s, \mathfrak p)$ を
Schemes, Lemma \ref{schemes-lemma-points-fibre-product} における
$X \times_S Y$ の点とする．このとき $z$ における茎上で，写像 $can$ は
$$
can_z :
\mathcal{O}_{X, x} \otimes_{\mathcal{O}_{S, s}} \mathcal{O}_{Y, y}
\longrightarrow
\mathcal{O}_{X \times_S Y, z}
$$
を与える．これは平坦な環準同型である．実際，終域は始域の局所化である
（詳細は省略する；$X$，$Y$，$S$ がアフィンの場合に帰着せよ）．
始域は一般には局所環でないことに注意せよ．このことからも，$can$ が
一般には同型でないことが分かる．

\medskip\noindent
さらに一般に，$\mathcal{O}_X$-加群 $\mathcal{F}$ と
$\mathcal{O}_Y$-加群 $\mathcal{G}$ が与えられているとする．このとき標準写像
% P12-ERR-0317: see ../control/ERRATA_CANDIDATES.jsonl
\begin{align*}
& p^{-1}\mathcal{F} \otimes_{f^{-1}\mathcal{O}_S} q^{-1}\mathcal{G} \\
& =
p^{-1}(\mathcal{F} \otimes_{\mathcal{O}_X} \mathcal{O}_X)
\otimes_{f^{-1}\mathcal{O}_S}
q^{-1}(\mathcal{O}_Y \otimes_{\mathcal{O}_Y} \mathcal{G}) \\
& =
p^{-1}\mathcal{F} \otimes_{p^{-1}\mathcal{O}_X} p^{-1}\mathcal{O}_X
\otimes_{f^{-1}\mathcal{O}_S}
q^{-1}\mathcal{O}_Y \otimes_{q^{-1}\mathcal{O}_Y} q^{-1}\mathcal{G} \\
& \xrightarrow{can}
p^{-1}\mathcal{F} \otimes_{q^{-1}\mathcal{O}_X}
\mathcal{O}_{X \times_S Y}
\otimes_{q^{-1}\mathcal{O}_Y} q^{-1}\mathcal{G}  \\
& =
p^{-1}\mathcal{F} \otimes_{q^{-1}\mathcal{O}_X}
\mathcal{O}_{X \times_S Y}
\otimes_{\mathcal{O}_{X \times_S Y}}
\mathcal{O}_{X \times_S Y}
\otimes_{q^{-1}\mathcal{O}_Y} q^{-1}\mathcal{G}  \\
& =
p^*\mathcal{F} \otimes_{\mathcal{O}_{X \times_S Y}} q^*\mathcal{G}
\end{align*}
があり，これが同型になることは稀である．
% END EXPANDED MODULE ch110_03.tex
% BEGIN EXPANDED MODULE ch110_04.tex
\section{局所環がすべて整域である非整な連結スキーム}
\label{examples-section-connected-locally-integral-not-integral}

\noindent
連結で，すべての局所環が整域であるが，整でないアフィンスキーム
$X = \Spec(A)$ の例を与える．$X$ が連結であることは，$A$ が
非自明な冪等元をもたないことを意味する；Algebra, Lemma
\ref{algebra-lemma-disjoint-decomposition} を参照せよ．$A$ において
$fg = 0$ であるとき，$X$ の各点が，$f$ または $g$ のいずれか一方が
消える近傍をもつならば，$X$ の局所環は整域である．一方，$A$ が
整域でなければ $X$ は整でない（Properties, Definition
\ref{properties-definition-integral}）．

\medskip\noindent
大まかにいえば，構成は次の通りである．$X_0$ をアフィン平面上の
十字形（座標軸の和集合）とする．次に，平面のブローアップにおける
$X_0$ の被約な全逆像を $X_1$ とする（$X_1$ は鎖状に並ぶ三つの
有理成分をもつ）．得られた曲面を $X_1$ の二つの特異点で
ブローアップし，$X_1$ の被約逆像（五つの有理成分をもつ）を
$X_2$ とし，以下同様に続ける．$X$ をその逆極限とする．
この構成の唯一の問題は，ブローアップでは射影直線が貼り付けられるため，
$X_1$ がアフィンでないことである．そこで代わりにアフィン直線を
貼り付けることによってこれを修正する（したがって，ここで得るスキームは
上で述べたものの開部分スキームとなる）．

\medskip\noindent
完全に代数的な構成は次の通りである．各 $k \ge 0$ に対し，$A_k$ を
次の環とする．その元は，$i = 0, \ldots, 2^k$ で添字付けられた
多項式 $p_i \in \mathbf{C}[x]$ の族で，$p_i(1) = p_{i + 1}(0)$
を満たすものである．$X_k = \Spec(A_k)$ と置く．$X_k$ は，
鎖状に横断的に交わる $2^k + 1$ 本のアフィン直線の和である．
環準同型 $A_k \to A_{k + 1}$ を
$$
(p_0, \ldots, p_{2^k})
\longmapsto
(p_0, p_0(1), p_1, p_1(1), \ldots, p_{2^k}),
$$
によって定める．すなわち，一つおきの多項式は定数である．これにより
$A_k$ は $A_{k + 1}$ の部分環と同一視される．$A$ を $A_k$ の
帰納極限（本質的にはそれらの和集合）とし，$X = \Spec(A)$ と置く．
各 $k$ に対し自然な埋め込み $A_k \to A$，すなわち写像
$X\to X_k$ がある．各 $A_k$ は連結であるが整でない．したがって
$A$ も連結であるが整でない．残るのは，$A$ の局所環が整域であることを
示すことである．

\medskip\noindent
$fg = 0$ を満たす $f, g \in A$ と $x \in X$ をとる．$f$ または
$g$ のいずれか一方が消える $x$ の近傍を構成しよう．
$f, g \in A_{k - 1}$ となる $k$ を選ぶ（添字が $k - 1$ であることに
注意せよ）．$y$ を $X_k$ における $x$ の像とする．$\mathcal{O}_{X_k}$
の切断とみなした $f$ または $g$ のいずれか一方が消えるような
$y$ の近傍が存在することを示せば十分である．$y$ が $X_k$ の
滑らかな点，すなわち $2^k + 1$ 本の直線のうち一本のみに属するなら，
これは明らかである．したがって $y$ が $2^k$ 個の特異点の一つであり，
$X_k$ の二つの成分が $y$ を通ると仮定してよい．しかし，この二成分の
一方（添字が奇数のもの）では，$f$ と $g$ は $X_{k - 1}$ 上の函数の
引き戻しなので，ともに定数である．至る所で $fg = 0$ であるから，
他方の成分では $f$ または $g$ のいずれか一方（例えば $f$）が消える．
すると $f$ は両方の成分上で消え，所望の結論が従う．
% END EXPANDED MODULE ch110_04.tex
% BEGIN EXPANDED MODULE ch110_05.tex
\section{完備でない完備化}
\label{examples-section-noncomplete-completion}

\noindent
$R$ を環，$\mathfrak m$ を極大イデアルとする．完備化
$$
R^\wedge = \lim R/\mathfrak m^n.
$$
を考える．$R^\wedge$ は，極大イデアル
$\mathfrak m' = \Ker(R^\wedge \to R/\mathfrak m)$ をもつ局所環であることに
注意せよ．実際，$x = (x_n) \in R^\wedge$ が $\mathfrak m'$ に属さなければ，
$y = (x_n^{-1}) \in R^\wedge$ は $xy = 1$ を満たす．したがって Algebra,
Lemma \ref{algebra-lemma-characterize-local-ring} により $R^\wedge$ は局所環である．
また，$R^\wedge$ はその極限位相に関して常に完備である
% P12-ERR-0318: see ../control/ERRATA_CANDIDATES.jsonl
（More on Algebra, Section \ref{more-algebra-section-topological-ring} の議論を
参照せよ）．しかし，さらに次の問いがある．
\begin{enumerate}
\item $\mathfrak m R^\wedge = \mathfrak m'$ は成り立つか？
\item $R^\wedge$-加群とみなした $R^\wedge$ は $\mathfrak m'$-進完備か？
\item $R$-加群とみなした $R^\wedge$ は $\mathfrak m$-進完備か？
\end{enumerate}
これらの問いにはすべて否定的な答えが与えられる．以下の例は
Bart de Smit と Hendrik Lenstra の未公刊のノートから採った．
\cite[Exercise III.2.12]{Bourbaki-CA} および
\cite[Example 1.8]{Yekutieli} も参照せよ．

\medskip\noindent
$k$ を体とし，$R = k[x_1, x_2, x_3, \ldots]$ および
$\mathfrak m = (x_1, x_2, x_3, \ldots)$ とする．$R^\wedge$ の元 $f$ を，
（多重添字記法を用いた）無限でありうる和
$$
f = \sum a_I x^I
$$
で，各 $d \geq 0$ に対し $|I| = d$ となる非零の $a_I$ が有限個しか
ないものと考える．極大イデアル $\mathfrak m' \subset R^\wedge$ は
定数項が零である $f$ 全体からなる．特に元
$$
f = x_1 + x_2^2 + x_3^3 + \ldots
$$
は $\mathfrak m'$ に属するが $\mathfrak m R^\wedge$ には属さない．
したがって，この例では (1) は偽である．一方，(1) が偽ならば (3) も
必ず偽である．実際，$\mathfrak m' = \Ker(R^\wedge \to R/\mathfrak m)$
であり，$n = 1$ として Algebra, Lemma \ref{algebra-lemma-hathat} を
適用できるからである．

\medskip\noindent
最後に，$R^\wedge$ が $\mathfrak m'$-進完備でないことを証明する．
$n \geq 1$ に対して $K_n = \Ker(R^\wedge \to R/\mathfrak m^n)$ と置く．
このとき短完全列
$$
0 \to K_n/(\mathfrak m')^n \to R^\wedge/(\mathfrak m')^n \to
R/\mathfrak m^n \to 0
$$
がある．射影 $R^\wedge \to R/\mathfrak m^{n + 1}$ は
$(\mathfrak m')^n$ を $\mathfrak m^n/\mathfrak m^{n + 1}$ の上に写す．
したがって $K_{n + 1} \to K_n/(\mathfrak m')^n$ は全射である．
ゆえに逆系 $\left(K_n/(\mathfrak m')^n\right)$ の遷移写像は全射であり，
逆極限をとると，Algebra, Lemma \ref{algebra-lemma-Mittag-Leffler} により
完全列
$$
0 \to \lim K_n/(\mathfrak m')^n \to
\lim R^\wedge/(\mathfrak m')^n \to
\lim R/\mathfrak m^n \to 0
$$
を得る．したがって，$R^\wedge$ が $\mathfrak m'$ に関して完備であることと，
すべての $n \geq 1$ に対して $K_n = (\mathfrak m')^n$ であることは同値である．

\medskip\noindent
この例の $R^\wedge$ が $\mathfrak m'$-進完備でないことを示すため，
$K_2 = \Ker(R^\wedge \to R/\mathfrak m^2)$ が $(\mathfrak m')^2$ に
等しくないことを示す．$(\mathfrak m')^2$ の元は，定数項が零である
$f_i, g_i \in R^\wedge$ を用いて有限和
\begin{equation}
\label{examples-equation-sum}
\sum\nolimits_{i = 1, \ldots, t} f_i g_i
\end{equation}
と書けることに注意せよ．例を得るため，次の形の $z \in K_2$ を選ぶ．
% P12-ERR-0319: see ../control/ERRATA_CANDIDATES.jsonl
$$
z = z_1 + z_2 + z_3 + \ldots
$$
ここで次の条件を課す．
\begin{enumerate}
\item 数列 $1 < d_1 < d_2 < d_3 < \ldots $ と
$0 < n_1 < n_2 < n_3 < \ldots$ が存在し，
\par\noindent
$z_i \in k[x_{n_i}, x_{n_i + 1}, \ldots, x_{n_{i + 1} - 1}]$ は
次数 $d_i$ の斉次元である．
\item 環 $k[[x_{n_i}, x_{n_i + 1}, \ldots, x_{n_{i + 1} - 1}]]$ において，
$z_i$ は $t \leq i$ を満たす和 (\ref{examples-equation-sum}) として書けない．
\end{enumerate}
これは明らかに $z$ が $(\mathfrak m')^2$ に属さないことを意味する．実際，十分大きな
$i$ に対して関係 (\ref{examples-equation-sum}) の像を環
$k[[x_{n_i}, x_{n_i + 1}, \ldots, x_{n_{i + 1} - 1}]]$ でとれば
矛盾が生じる．したがって，任意の $t > 0$ に対して，ある $d \gg 0$ と
整数 $n$ が存在し，与えられた $t$ に対して
$k[[x_1, \ldots, x_n]]$ の中で和 (\ref{examples-equation-sum}) として書けない
次数 $d$ の斉次元 $z \in k[x_1, \ldots, x_n]$ が存在することを
証明すれば十分である．$n > 2t$ とし，$k$ の標数と互いに素な任意の
$d > 1$ をとって，$z = \sum_{i = 1, \ldots, n} x_i^d$ と置く．すると
イデアル
$$
(\frac{\partial z}{\partial x_1}, \ldots, \frac{\partial z}{\partial x_n})
=
(dx_1^{d - 1}, \ldots, dx_n^{d - 1})
$$
の零点集合は一点からなる．一方，Leibniz 則により
$$
\frac{\partial ( \sum\nolimits_{i = 1, \ldots, t} f_i g_i ) }{\partial x_j}
\in (f_1, \ldots, f_t, g_1, \ldots, g_t)
$$
であるから，これらの偏導函数の零点集合は少なくとも
$$
V(f_1, \ldots, f_t, g_1, \ldots, g_t) \subset
\Spec(k[[x_1, \ldots, x_n]]).
$$
を含む．しかし
$V(f_1, \ldots, f_t, g_1, \ldots, g_t)$ の次元は少なくとも
$n - 2t \geq 1$ なので，これは矛盾である．

\begin{lemma}
\label{examples-lemma-noncomplete-completion}
局所環 $R$ と極大イデアル $\mathfrak m$ であって，$R$ の
$\mathfrak m$ に関する完備化 $R^\wedge$ が次の性質をもつものが存在する．
\begin{enumerate}
\item $R^\wedge$ は局所環であるが，その極大イデアルは
$\mathfrak m R^\wedge$ に等しくない．
\item $R^\wedge$ は完備局所環でない．
\item $R^\wedge$ は $R$-加群として $\mathfrak m$-進完備でない．
\end{enumerate}
\end{lemma}

\begin{proof}
上の議論から従う．実際，$R = k[x_1, x_2, x_3, \ldots]$ とすると，
局所化 $R_{\mathfrak m}$ の完備化は $R$ の完備化に等しい．
\end{proof}
% END EXPANDED MODULE ch110_05.tex
% BEGIN EXPANDED MODULE ch110_06.tex
\section{完備でない商}
\label{examples-section-noncomplete-quotient}

\noindent
$k$ を体とし，
$$
R = k[t, z_1, z_2, z_3, \ldots, w_1, w_2, w_3, \ldots, x]/
(z_it - x^iw_i, z_i w_j)
$$
と置く．特に，この環では $z_iz_jt = 0$ であることに注意せよ．
$R$ の任意の元 $f$ は，一意的に有限和
$$
f = \sum\nolimits_{i = 0, \ldots, d} f_i x^i
$$
と書ける．ここで，各 $f_i \in k[t, z_i, w_j]$ は積 $z_it$ または
$z_iw_j$ を含む項をもたない．さらに，$f$ がこのように書かれているとき，
$f \in (x^n)$ であることと，$i < n$ に対して $f_i = 0$ であることは
同値である．したがって $x$ は非零因子であり，$\bigcap (x^n) = 0$ である．
$R^\wedge$ をイデアル $(x)$ に関する $R$ の完備化とする．
$R^\wedge$ は $(x)$-進完備であることに注意せよ；Algebra, Lemma
\ref{algebra-lemma-hathat-finitely-generated} を参照せよ．上のことから，
$R^\wedge$ の元は一意的に無限和
$$
f = \sum\nolimits_{i = 0}^\infty f_i x^i
$$
と書ける．ここでも，各 $f_i \in k[t, z_i, w_j]$ は積 $z_it$ または
$z_iw_j$ を含む項をもたない．元
$$
f = \sum\nolimits_{i = 1}^\infty x^i w_i =
xw_1 + x^2w_2 + x^3w_3 + \ldots
$$
を考える．すなわち $f_n = w_n$ である．すべての $n$ に対して
$f \in (t , x^n)$ であることに注意せよ．実際，すべての $m$ に対し
$x^mw_m \in (t)$ だからである．

$f \not \in (t)$ であると主張する．これを示すため，
$tg = f$ であると仮定し，$g = \sum g_lx^l$ を上と同じ標準形で書く．
$tz_iz_j = 0$ なので，どの $g_l$ も積 $z_iz_j$ を含む項をもたないと
仮定してよい．$tg$ を標準形にする過程を調べると，次が分かる．
$g_l$ の任意の項 $c m$ をとる．ここで $c \in k$ で，$m$ は
$t, z_i, w_j$ の単項式である．このとき次の置換を行う．
\begin{enumerate}
\item 単項式 $m$ がどの $z_i$ も含まなければ，$ctm$ は $f_l$ の項である．
\item 単項式 $m$ がある $z_i$ を含むなら，$m = z_i$ であり，
% P12-ERR-0320: see ../control/ERRATA_CANDIDATES.jsonl
$cw_i$ は $f_{l + i}$ の項となる．
\end{enumerate}
$g_0$ は多項式なので，その中に現れる変数 $z_i$ は有限個しかない．
$g_0$ に $z_n$ が現れないように $n$ を選ぶ．すると上の規則によれば
$f_n$ に $w_n$ は現れず，矛盾である．したがって
$R^\wedge/(t)$ は完備でない；Algebra, Lemma
\ref{algebra-lemma-quotient-complete} を参照せよ．

\begin{lemma}
\label{examples-lemma-noncomplete-quotient}
単項イデアル $I$ に関して完備な環 $R$ と単項イデアル $J$ であって，
$R/J$ が $I$-進完備でないものが存在する．
\end{lemma}

\begin{proof}
上の議論を参照せよ．
\end{proof}
% END EXPANDED MODULE ch110_06.tex
% BEGIN EXPANDED MODULE ch110_07.tex
\section{完備化は完全でない}
\label{examples-section-completion-not-exact}

\noindent
簡単な例は次の通りである．$R = k[t]$ と仮定する．$R$-加群の
$I = (t)$ に関する完備化を $M^\wedge$ と書く．
$P = K = \bigoplus_{n \in \mathbf{N}} R$ および
$M = \bigoplus_{n \in \mathbf{N}} R/(t^n)$ と置く．このとき短完全列
$0 \to K \to P \to M \to 0$ があり，最初の写像は第 $n$ 成分上で
$t^n$ を掛けることによって与えられる．
$0 \to K^\wedge \to P^\wedge \to M^\wedge \to 0$ は中央で完全でないと
主張する．実際，$\xi = (t^2, t^3, t^4, \ldots) \in P^\wedge$ は
$M^\wedge$ で零に写るが，$K^\wedge \to P^\wedge$ の像には属さない．
なぜなら，その原像は $(t, t, t, \ldots)$ でなければならないが，これは
$K^\wedge$ の元でないからである．

\medskip\noindent
より「小さい」例は次の通りである．Lemma
\ref{examples-lemma-noncomplete-quotient} の状況では，短完全列
$0 \to J \to R \to R/J \to 0$ は $I$ に関する完備化をとった後には
完全でなくなる．実際，$f \in J$ が生成元ならば，
$f : R \to J$ は全射であり，したがって $R \to J^\wedge$ は全射である．
ゆえに $J^\wedge \to R$ の像は $(f) = J$ である．しかし $R/J$ が
完備でないということは，全射 $R \to (R/J)^\wedge$ の核が $J$ より
真に大きいことを意味する；Algebra, Lemmas
\ref{algebra-lemma-completion-generalities} および
\ref{algebra-lemma-quotient-complete} を参照せよ．同じ理由により，列
$R \to R \to R/(f) \to 0$ も完備化後には完全でなくなる．

\begin{lemma}
\label{examples-lemma-completion-not-exact}
\begin{slogan}
完備化は一般に左完全でも右完全でもない．
\end{slogan}
完備化は一般には完全函手でなく，右完全ですらない．これは，$I$ が
有限生成で，有限表示加群の圏に制限した場合にも成り立つ．
\end{lemma}

\begin{proof}
上の議論を参照せよ．
\end{proof}
% END EXPANDED MODULE ch110_07.tex
% BEGIN EXPANDED MODULE ch110_08.tex
\section{完備加群の圏はアーベル圏でない}
\label{examples-section-non-abelian}

\noindent
$R$ を環，$I \subset R$ を有限生成イデアルとする．$I$-進完備な
$R$-加群の圏 $\mathcal{A}$ を考える；Algebra, Definition
\ref{algebra-definition-complete} を参照せよ．
$\varphi : M \to N$ を $\mathcal{A}$ の射とする．$\mathcal{A}$ における
$\varphi$ の余核は，通常の余核の $I$-進完備化
$(\Coker(\varphi))^\wedge$ である（$I$ は有限生成なので，この完備化は
完備である；Algebra, Lemma
\ref{algebra-lemma-hathat-finitely-generated} を参照せよ）．
$K = \Ker(\varphi)$ と置く．$\varphi$ は連続なので，これは $M$ の
閉部分加群である．下の Lemma \ref{examples-lemma-closed-is-complete} により，
$K$ は $I$-進完備であり，したがって $\mathcal{A}$ における
$\varphi$ の核である．

\begin{lemma}
\label{examples-lemma-closed-is-complete}
$I \subset R$ を環のイデアル，$M$ を $I$-進完備な $R$-加群，
$N \subset M$ を部分加群とする．次の条件は同値である．
\begin{enumerate}
\item $N$ は $M$ の閉部分加群である．
\item $N = \bigcap_{n \geq 1} (N + I^n M)$ である．
\item $M/N$ は $I$-進完備である．
\end{enumerate}
$I$ が有限生成ならば，これらの条件から $N$ が $I$-進完備であることが従う．
\end{lemma}

\begin{proof}
(1) と (2) の同値性は Formal Spaces, Lemma
\ref{formal-spaces-lemma-closed} から従う．(2) と (3) の同値性は
Algebra, Lemma \ref{algebra-lemma-quotient-complete} である．
$I$ が有限生成で，(1) が成り立つと仮定する．$N^\wedge$ を $N$ の
$I$-進完備化とする．$M$ は $I$-進完備なので，包含写像は
$$
N \to N^\wedge \to M
$$
と分解する．$N^\wedge \to M$ は連続で，$N$ は $N^\wedge$ で稠密であり，
$N$ は $M$ で閉であるから，$N^\wedge \to M$ は $N$ を経由する．したがって
$N^\wedge = N \oplus C$
であるような $A$-加群 $C$ が存在する．
% P12-ERR-0321: see ../control/ERRATA_CANDIDATES.jsonl
ゆえに $N$ は $I$-進完備加群 $N^\wedge$ の直和因子である；Algebra,
Lemma \ref{algebra-lemma-hathat-finitely-generated} を参照せよ（ここで
$I$ が有限生成であることを用いる）．したがって $N$ は $I$-進完備である．
\end{proof}

\noindent
一般には圏 $\mathcal{A}$ において $\Im \not = \Coim$ となることを示す例を
与える．$R = \mathbf{Z}_p = \lim_n \mathbf{Z}/p^n\mathbf{Z}$ を
$p$-進整数環とし，$I = (p)$ とする．写像
$$
\text{diag}(1, p, p^2, \ldots) :
\left(\bigoplus\nolimits_{n \geq 1} \mathbf{Z}_p\right)^\wedge
\longrightarrow
\prod\nolimits_{n \geq 1} \mathbf{Z}_p
$$
を考える．左辺は直和の $p$-進完備化である．したがって左辺の元は，
$x_i \in \mathbf{Z}_p$ であり，$i \to \infty$ のとき $p$-進付値が
$v_p(x_i) \to \infty$ となるベクトル $(x_1, x_2, x_3, \ldots)$ である．
これは $(x_1, px_2, p^2x_3, \ldots)$ に写る．ゆえに
$(1, p, p^2, \ldots)$ は像の閉包には属するが，像には属さない．
上で述べた核と余核の記述から，この写像について
$\Im \not = \Coim$ であることは明らかである．

\begin{lemma}
\label{examples-lemma-complete-modules-not-abelian}
$R$ を環，$I \subset R$ を有限生成イデアルとする．$I$-進完備な
$R$-加群の圏は核と余核をもつが，一般には $R$ が Noether 環であっても
アーベル圏ではない．
\end{lemma}

\begin{proof}
上を参照せよ．
\end{proof}
% END EXPANDED MODULE ch110_08.tex
% BEGIN EXPANDED MODULE ch110_09.tex
\section{導来的完備加群の圏}
\label{examples-section-derived-complete-modules}

\noindent
この節を読む前に More on Algebra, Section
\ref{more-algebra-section-derived-complete-modules} を読んでほしい．

\medskip\noindent
$A$ を環，$I$ を $A$ のイデアルとし，More on Algebra, Definition
\ref{more-algebra-definition-derived-complete} で定義された導来的完備加群の圏を
$\mathcal{C}$ と書く．

\medskip\noindent
$T$ を集合とし，$M_t$（$t \in T$）を導来的完備加群の族とする．
一般には $\bigoplus M_t$ は導来的完備加群でないと主張する．具体例として，
$A = \mathbf{Z}_p$，$I = (p)$ とし，
$\bigoplus_{n \in \mathbf{N}} \mathbf{Z}_p$ を考える．
$\bigoplus_{n \in \mathbf{N}} \mathbf{Z}_p$ からその $p$-進完備化への写像は
全射でない．したがって $\bigoplus_{n \in \mathbf{N}} \mathbf{Z}_p$ が
導来的完備であることはありえない．もしそうなら逆の結論が従うからである；
More on Algebra, Lemma
\ref{more-algebra-lemma-complete-derived-complete} を参照せよ．
したがって包含函手 $\mathcal{C} \to \text{Mod}_A$ は直和とも
（フィルター）余極限とも可換しない．

\medskip\noindent
$I$ が有限生成であると仮定する．More on Algebra, Section
\ref{more-algebra-section-derived-complete-modules} の議論により，圏
$\mathcal{C}$ は任意の余極限をもつ．しかし，$\mathcal{C}$ における
フィルター余極限は完全でないと主張する．実際，
$A = \mathbf{Z}_p$，$I = (p)$ とする．$p$-進完備な $A$-加群の包含写像
$f_n : \mathbf{Z}_p/p\mathbf{Z}_p \to \mathbf{Z}_p/p^n\mathbf{Z}_p$
を，$p^{n - 1}$ 倍によって定める．可換図式
$$
\xymatrix{
\mathbf{Z}_p/p\mathbf{Z}_p \ar[r]_{f_n} \ar[d]^1 &
\mathbf{Z}_p/p^n\mathbf{Z}_p \ar[d]_p \\
\mathbf{Z}_p/p\mathbf{Z}_p \ar[r]^{f_{n + 1}} &
\mathbf{Z}_p/p^{n + 1}\mathbf{Z}_p
}
$$
がある．主張は，圏 $\mathcal{C}$ におけるこれらの包含写像の余極限が
写像 $\mathbf{Z}_p/p\mathbf{Z}_p \to 0$ を与える，ということである．
実際，右側の系の $\text{Mod}_A$ における余極限は
$\mathbf{Q}_p/\mathbf{Z}_p$ である．したがって $\mathcal{C}$ における
余極限は
$$
H^0((\mathbf{Q}_p/\mathbf{Z}_p)^\wedge) =
H^0(\mathbf{Z}_p[1]) = 0
$$
である．これは More on Algebra, Section
\ref{more-algebra-section-derived-complete-modules} による．ここで
${}^\wedge$ は導来的完備化を表す．以上で主張が証明された．

\begin{lemma}
\label{examples-lemma-derived-complete-modules}
$A$ を環，$I \subset A$ をイデアルとする．導来的完備加群の圏
$\mathcal{C}$ はアーベル圏であり，包含函手
$F : \mathcal{C} \to \text{Mod}_A$ は完全で，任意の極限と可換する．
$I$ が有限生成ならば，$\mathcal{C}$ は任意の直和と余極限をもつが，
一般に $F$ はそれらと可換しない．最後に，フィルター余極限は一般に
$\mathcal{C}$ において完全でない．したがって $\mathcal{C}$ は
Grothendieck アーベル圏でない．
\end{lemma}

\begin{proof}
More on Algebra, Lemma
\ref{more-algebra-lemma-derived-complete-modules} および上の議論を参照せよ．
\end{proof}
% END EXPANDED MODULE ch110_09.tex
% BEGIN EXPANDED MODULE ch110_10.tex
\section{平坦でない完備化}
\label{examples-section-nonflat}

\noindent
Noether 環の場合とは異なり，環のイデアルに関する完備化は必ずしも
平坦でない．この現象の二つの例を More on Algebra, Example
\ref{more-algebra-example-not-glueing-pair} で見た．この節では，さらに
二つの例を与える．

\begin{lemma}
\label{examples-lemma-countable-fg-tensor}
$R$ を環とし，$M$ を可算な $R$-加群とする．このとき $M$ が有限
$R$-加群であることと，$M \otimes_R R^\mathbf{N} \to M^\mathbf{N}$ が
全射であることは同値である．
\end{lemma}

\begin{proof}
$M$ が有限加群ならば，Algebra, Proposition
\ref{algebra-proposition-fg-tensor} により写像は全射である．逆に写像が
全射であると仮定する．$m_1, m_2, m_3, \ldots$ を $M$ の元の列挙とする．
テンソル積の元 $\sum_{j = 1, \ldots, m} x_j \otimes a_j$ が
元 $(m_n) \in M^\mathbf{N}$ に写るとする．Algebra, Proposition
\ref{algebra-proposition-fg-tensor} の証明と同様に，
$x_1, \ldots, x_m$ は $R$ 上 $M$ を生成する．
\end{proof}

\begin{lemma}
\label{examples-lemma-countable-fp-tensor}
$R$ を可算環，$M$ を可算な $R$-加群とする．このとき $M$ が
有限表示であることと，標準写像
$M \otimes_R R^\mathbf{N} \to M^\mathbf{N}$ が同型であることは同値である．
\end{lemma}

\begin{proof}
$M$ が有限表示加群ならば，Algebra, Proposition
\ref{algebra-proposition-fp-tensor} により写像は同型である．逆に写像が
同型であると仮定する．Lemma \ref{examples-lemma-countable-fg-tensor} により，
加群 $M$ は有限である．核を $K$ とする全射
$R^{\oplus m} \to M$ を選ぶ．$K$ は $R^{\oplus m}$ の部分加群として
可算である．Algebra, Proposition \ref{algebra-proposition-fp-tensor} の
証明と同様に論じると，$K \otimes_R R^\mathbf{N} \to K^\mathbf{N}$ は
全射である．したがって Lemma \ref{examples-lemma-countable-fg-tensor} により
$K$ は有限 $R$-加群である．ゆえに $M$ は有限表示である．
\end{proof}

\begin{lemma}
\label{examples-lemma-countable-coherent}
$R$ を可算環とする．このとき，$R$ がコヒーレントであることと
$R^\mathbf{N}$ が平坦な $R$-加群であることは同値である．
\end{lemma}

\begin{proof}
$R$ がコヒーレントならば，Algebra, Proposition
\ref{algebra-proposition-characterize-coherent} により
$R^\mathbf{N}$ は平坦加群である．$R^\mathbf{N}$ が平坦であると仮定する．
$I \subset R$ を有限生成イデアルとする．補題を証明するため，$I$ が
$R$-加群として有限表示であることを示す．写像
$I \otimes_R R^\mathbf{N} \to R^\mathbf{N}$ は，$R^\mathbf{N}$ が
平坦であるため単射であり，その像は Lemma
\ref{examples-lemma-countable-fg-tensor} により $I^\mathbf{N}$ である．したがって
Lemma \ref{examples-lemma-countable-fp-tensor} から結論が従う．
\end{proof}

\noindent
$R$ を可算環とする．$R[[x]]$ は $R$-加群として $R^\mathbf{N}$ と
同型であることに注意せよ．Lemma \ref{examples-lemma-countable-coherent} により，
$R \to R[[x]]$ が平坦であることと $R$ がコヒーレントであることは
同値である．可算な非コヒーレント環は多数存在する．例えば
$$
R = k[y, z, a_1, b_1, a_2, b_2, a_3, b_3, \ldots]/
(a_1 y + b_1 z, a_2 y + b_2 z, a_3 y + b_3 z, \ldots)
$$
がそうである．ここで $k$ は可算体である．この環がコヒーレントでないのは，
$R$ のイデアル $(y, z)$ が有限表示 $R$-加群でないからである．
$R[[x]]$ は単項イデアル $(x)$ による $R[x]$ の完備化であることに注意せよ．

\begin{lemma}
\label{examples-lemma-completion-polynomial-ring-not-flat}
$R[x]$ の $(x)$ における完備化 $R[[x]]$ が $R$ 上平坦でなく，したがって
$R[x]$ 上でも平坦でないような環が存在する．
\end{lemma}

\begin{proof}
上の議論を参照せよ．
\end{proof}

\noindent
$R$ 上では $R[[x]]$ が平坦であるが，$R[[x]]$ は $R[x]$ 上では
平坦でないような環 $R$ が存在することが分かっている．Badam Baplan による
\href{https://math.stackexchange.com/users/164860/badam-baplan}{この投稿}
を参照せよ．実際，$R$ を付値環とする．このとき $R$ はコヒーレントであり
（Algebra, Example \ref{algebra-example-valuation-ring-coherent}），したがって
Algebra, Proposition \ref{algebra-proposition-characterize-coherent} により
$R[[x]]$ は $R$ 上平坦である．一方，次の補題が成り立つ．

\begin{lemma}
\label{examples-lemma-almost-integral-when-powerseries-flat}
$R$ を分数体 $K$ をもつ整域とする．$R[[x]]$ が $R[x]$ 上平坦ならば，
$R$ が正規であることと $R$ が完整閉であることは同値である
（Algebra, Definition \ref{algebra-definition-almost-integral}）．
\end{lemma}

\begin{proof}
$\alpha \in K$ と非零元 $r \in R$ があり，すべての $n \geq 1$ に対して
$r \alpha^n \in R$ であると仮定する．$R[[x]]$ の元
$f = \sum r \alpha^{n - 1} x^n$ を考える．
$a, b \in R$，$b$ は非零として $\alpha = a/b$ と書く．すると
$(a x - b)f = -rb$ である．したがって $rb$ はイデアル
$(ax - b)R[[x]]$ に属する．
$S = \{h \in R[x] : h(0) = 1\}$ と置く．これは乗法的部分集合であり，
$R[x] \to R[[x]]$ の平坦性から $S^{-1}R[x] \to R[[x]]$ が忠実平坦で
あることが従う（詳細は省略する；Algebra, Lemma
\ref{algebra-lemma-ff-rings} を用いよ）．したがって
$$
S^{-1}R[x]/(ax - b)S^{-1}R[x] \to R[[x]]/(ax - b)R[[x]]
$$
は単射である．よって，ある $h \in S$ と $g \in R[x]$ に対して
$h rb = (ax - b) g$ である．$h = 1 + h_1 x + \ldots + h_d x^d$ と書けば
$$
1 + h_1 x + \ldots + h_d x^d = (1/r)(\alpha x - 1)g
$$
を得る．$K[x]$ におけるこの因数分解から $K[x^{-1}]$ における対応する
因数分解が得られ，$\alpha$ が $R$ に係数をもつモニック多項式の根で
あることが分かる．
\end{proof}

\begin{lemma}
\label{examples-lemma-completion-polynomial-ring-not-flat-bis}
$R$ が次元 $> 1$ の付値環ならば，$R[[x]]$ は $R$ 上平坦であるが
$R[x]$ 上平坦でない．
\end{lemma}

\begin{proof}
上の議論により，$R$ が完整閉でないことを示せば十分である
（付値環は正規である；Algebra, Lemma
\ref{algebra-lemma-valuation-ring-normal} を参照せよ）．素イデアルの鎖
$\mathfrak p \subset \mathfrak m \subset R$ をとる．非零元
$x \in \mathfrak p$ と $y \in \mathfrak m \setminus \mathfrak p$ を選ぶ．
すべての $n \geq 1$ に対して $x y^{-n} \in R$ である（そうでなければ
$y^n/x \in R$ となるが，$y \not \in \mathfrak p$ なので不合理である）．
したがって $1/y$ は $R$ 上概整であるが，$R$ には属さない．
\end{proof}

\noindent
次に，局所化の完備化が平坦でない例を構成する．そのため環
$$
R = k[y, z, a_1, a_2, a_3, \ldots]/(ya_i, a_i a_j)
$$
を考える．$f \in R$ を $z$ の剰余類とする．環準同型
\begin{equation}
\label{examples-equation-nonflat}
R[[x]] \longrightarrow R_f[[x]]
\end{equation}
は平坦でないと主張する．$I$ を写像
$y : R[[x]] \to R[[x]]$ の核とする．$I$ の典型的な元 $g$ は
$g = \sum g_{n, m} a_mx^n$ の形をしている．ここで
$g_{n, m} \in k[z]$ であり，固定した $n$ に対して非零の
$g_{n, m}$ は有限個しかない．$J$ を写像
$y : R_f[[x]] \to R_f[[x]]$ の核とする．
$J \not = I R_f[[x]]$ と主張する．もし等しいならば，
$$
\sum z^{-n} a_n x^n = \sum\nolimits_{i = 1, \ldots, m} h_i g_i
$$
が，ある $m \geq 1$，$g_i \in I$，$h_i \in R_f[[x]]$ に対して
成り立つはずである．
\par\noindent
$h_i = \bar h_i \bmod (y, a_1, a_2, a_3, \ldots)$ とし，
$\bar h_i \in k[z, 1/z][[x]]$ とする．$a_n$ の係数を比較し，上で述べた
$g_i$ の元の記述を用いると，
$$
z^{-n} x^n = \sum \bar h_i \bar g_{i, n}
$$
を，ある $\bar g_{i, n} \in k[z][[x]]$ に対して得る．これはすべての
$z^{-n}x^n$ が，元 $\bar h_i$ で生成される有限
$k[z][[x]]$-加群に含まれることを意味する．$k[z][[x]]$ は Noether 環なので，
% P12-ERR-0322: see ../control/ERRATA_CANDIDATES.jsonl
$1, z^{-1}x, z^{-2}x^2, \ldots$ で生成される $k[z, 1/z][[x]]$ の
$R[z][[x]]$-部分加群が有限であることが従う．Algebra, Lemma
\ref{algebra-lemma-characterize-integral-element} によれば，これは
$z^{-1}x$ が $k[z][[x]]$ 上整であることを意味するが，不合理である．
一方，(\ref{examples-equation-nonflat}) が平坦ならば，完全列
$0 \to I \to R[[x]] \xrightarrow{y} R[[x]]$ を $R_f[[x]]$ と
テンソルすることにより $J = IR_f[[x]]$ を得るはずである．

\begin{lemma}
\label{examples-lemma-nonflat-completion-localization}
単項イデアル $I$ に関して完備な環 $A$ と元 $f \in A$ であって，
$A_f$ の $I$-進完備化 $A_f^\wedge$ が $A$ 上平坦でないものが存在する．
\end{lemma}

\begin{proof}
$A = R[[x]]$，$I = (x)$ と置き，$R_f[[x]]$ が $R[[x]]_f$ の
完備化であることに注意せよ．
\end{proof}
% END EXPANDED MODULE ch110_10.tex
% BEGIN EXPANDED MODULE ch110_11.tex
\section{準連接加群の非アーベル圏}
\label{examples-section-nonabelian-QCoh}

\noindent
Sheaves on Stacks, Section \ref{stacks-sheaves-section-quasi-coherent} で，
$\Sch$ 上の亜群にファイバー化された圏上の準連接加群の圏を定義した．
Sheaves on Stacks, Section
\ref{stacks-sheaves-section-quasi-coherent-algebraic-stacks} では，代数スタックに
対してこの圏がアーベル圏であることを示すが，この節では形式代数空間に
対してはそうでないことを示す．

\medskip\noindent
実際，$\mathbf{Z}_p$ を $p$-進位相を備えた位相環とみなし，
$X = \text{Spf}(\mathbf{Z}_p)$ と置く；Formal Spaces, Definition
\ref{formal-spaces-definition-affine-formal-spectrum} を参照せよ．このとき
$X$ は $(\Sch/\mathbf{Z})_{fppf}$ 上の集合の層であり，集合亜群のスタック
$\mathcal{X}$ を与える；Stacks, Lemma
\ref{stacks-lemma-when-stack-in-sets} を参照せよ．したがって Sheaves on Stacks,
Section \ref{stacks-sheaves-section-quasi-coherent-algebraic-stacks} の議論を
適用できる．

\medskip\noindent
$\mathcal{F}$ を $\mathcal{X}$ 上の準連接加群とする．
$X = \colim \Spec(\mathbf{Z}/p^n\mathbf{Z})$ なので，Sheaves on Stacks,
Lemma \ref{stacks-sheaves-lemma-quasi-coherent} から，$\mathcal{F}$ は
次を満たす列 $(\mathcal{F}_n)$ によって与えられることが分かる．
\begin{enumerate}
\item $\mathcal{F}_n$ は $\Spec(\mathbf{Z}/p^n\mathbf{Z})$ 上の準連接加群である．
\item 遷移写像は同型
$\mathcal{F}_n = \mathcal{F}_{n + 1}/p^n\mathcal{F}_{n + 1}$ を与える．
\end{enumerate}
加群に移せば，$\mathcal{F}$ は系 $(M_n)$ に対応する．ここで各 $M_n$ は
$p^n$ で零化されるアーベル群であり，遷移写像は同型
$M_n = M_{n + 1}/p^n M_{n + 1}$ を誘導する．この状況では加群
$M = \lim M_n$ は $p$-進完備であり，$M_n = M/p^n M$ である；Algebra,
Lemma \ref{algebra-lemma-limit-complete} を参照せよ．したがって $X$ 上の
準連接加群の圏は $p$-進完備アーベル群の圏と同値である．この圏は
アーベル圏でない；Section \ref{examples-section-non-abelian} を参照せよ．

% P12-ERR-0323: see ../control/ERRATA_CANDIDATES.jsonl
\begin{lemma}
\label{examples-lemma-quasi-coherent-not-abelian}
形式代数空間 $X$ 上の準連接\footnote{ここでいう準連接加群は上で定義したものである．
Stacks project における構成の仕方から，これは実際に正しい定義である．
上で見た通り，この定義は $\text{Spf}(A)$ 上の準連接加群として素朴に
定義されるであろうもの，すなわち完備 $A$-加群と一致する．}
加群の圏は，一般にはアーベル圏でない．これは $X$ が Noether アフィン
形式代数空間である場合にも成り立つ．
\end{lemma}

\begin{proof}
上の議論を参照せよ．
\end{proof}
% END EXPANDED MODULE ch110_11.tex
% BEGIN EXPANDED MODULE ch110_12.tex
\section{局所的には分裂する非分裂列}
\label{examples-section-nonsplit-locally-split}

\noindent
短完全列
$$
0 \to M \to
\bigoplus\nolimits_{p\text{ prime}} \mathbf{Z}_{(p)} \to \mathbf{Q} \to 0
$$
を考える．ここで $M$ は右側の写像の核であり，この写像は各直和成分上で
包含 $\mathbf{Z}_{(p)} \to \mathbf{Q}$ を用いる．この $\mathbf{Z}$-加群の
列は分裂しない．実際，非零準同型 $\mathbf{Q} \to \mathbf{Z}_{(p)}$ は
存在しないからである．一方，任意の素数 $p$ で局所化すると，列は分裂する．

\begin{lemma}
\label{examples-lemma-nonsplit-locally-split}
環 $R$ と，Zariski 局所的には分裂するが分裂しない加群の列が存在する．
\end{lemma}

\begin{proof}
上の議論を参照せよ．
\end{proof}
% END EXPANDED MODULE ch110_12.tex
% BEGIN EXPANDED MODULE ch110_13.tex
\section{正則列と基底変換}
\label{examples-section-regular-base-change}

\noindent
正則列 $(x, y, z)$ をもつ環 $R$ であって，
$x\delta = y\delta = 0$ を満たす非零元 $\delta \in R/zR$ が存在するものを
構成する．

\medskip\noindent
例を構成するため，まず任意の体 $k$ に対して環 $k[x, y, z]$ 上の
特異な加群 $E$ を構成する．すなわち，$E$ を次の図式における押し出しとする．
$$
\xymatrix{
\frac{xk[x, y, z, y^{-1}]}{xyk[x, y, z]} \ar[r] \ar[d]^{z/x} &
\frac{k[x, y, z, x^{-1}, y^{-1}]}{yk[x, y, z, x^{-1}]} \ar[r] \ar[d] &
\frac{k[x, y, z, x^{-1}, y^{-1}]}{yk[x, y, z, x^{-1}] + xk[x, y, z, y^{-1}]}
\ar[d] \\
\frac{k[x, y, z, y^{-1}]}{yzk[x, y, z]} \ar[r] &
E \ar[r] &
\frac{k[x, y, z, x^{-1}, y^{-1}]}{yk[x, y, z, x^{-1}] + xk[x, y, z, y^{-1}]}
}
$$
ここで各行は短完全列である（組版上の問題のため，両端の零を省略した）．
$E$ は別の仕方では
$$
E = \{(f, g) \mid f \in k[x, y, z, x^{-1}, y^{-1}],
g \in k[x, y, z, y^{-1}] \}/\sim
$$
と記述できる．ここで $(f, g) \sim (f', g')$ であるための必要十分条件は，
ある $h \in k[x, y, z, y^{-1}]$ が存在して
$$
f = f' + xh \bmod yk[x, y, z, x^{-1}], \quad
g = g' - zh \bmod yzk[x, y, z]
$$
となることである．次を主張する：(a) $x : E \to E$ は単射である，
(b) $y : E/xE \to E/xE$ は単射である，(c) $E/(x, y)E = 0$ である，
(d) $x\delta = y\delta = 0$ を満たす非零元 $\delta \in E/zE$ が存在する．

\medskip\noindent
(a) を示すため，$E$ の元を与える対 $(f, g)$ で
$(xf, xg) \sim 0$ となるものをとる．このとき
$xf + xh \in yk[x, y, z, x^{-1}]$ および
$xg - zh \in yzk[x, y, z]$ を満たす
$h \in k[x, y, z, y^{-1}]$ が存在する．この元は，$j \leq 0$ のみが
現れる単項式の和 $h = \sum a_{i, j, k}x^iy^jz^k$ であると仮定してよい．
すると $xg - zh \in yzk[x, y, z]$ から $i > 0$ のみが現れることが従う．
すなわち，ある $h' \in k[x, y, z, y^{-1}]$ に対して $h = xh'$ である．
このとき $(f, g) \sim (f + xh', g - zh')$ であるから，$g = 0$ かつ
$h = 0$ と仮定してよい．この場合，$xf \in yk[x, y, z, x^{-1}]$ から
$f \in yk[x, y, z, x^{-1}]$ が従い，$(f, g) \sim 0$ であることが分かる．
したがって $x : E \to E$ は単射である．

\medskip\noindent
乗法 $x$ は
$\frac{k[x, y, z, x^{-1}, y^{-1}]}{yk[x, y, z, x^{-1}]}$ 上の同型なので，
$E/xE$ は
$$
\frac{k[x, y, z, y^{-1}]}{
yzk[x, y, z] + xk[x, y, z, y^{-1}] + zk[x, y, z, y^{-1}]}
=
\frac{k[x, y, z, y^{-1}]}{xk[x, y, z, y^{-1}] + zk[x, y, z, y^{-1}]}
$$
と同型である．したがって $y$ による乗法は $E/xE$ 上の同型である．
これは明らかに (b) と (c) を意味する．

\medskip\noindent
$e \in E$ を $(1, 0)$ の同値類とする．$e \in zE$ と仮定する．すると
$f \in k[x, y, z, x^{-1}, y^{-1}]$，$g \in k[x, y, z, y^{-1}]$，および
$h \in k[x, y, z, y^{-1}]$ が存在して
$$
1 + zf + xh \in yk[x, y, z, x^{-1}], \quad
0 + zg - zh \in yzk[x, y, z].
$$
となる．しかしこれは不可能である．単項式 $1$ は $zf$ にも $xh$ にも
現れえないからである．一方，$ye = 0$ であり，
$xe = (x, 0) \sim (0, -z) = z(0, -1)$ である．したがって $\delta$ を
$E/zE$ における $e$ の合同類とすれば，(d) が得られる．

\begin{lemma}
\label{examples-lemma-strange-regular-sequence}
局所環 $R$ と（極大イデアル内の）正則列 $x, y, z$ であって，
$x\delta = y\delta = 0$ を満たす非零元 $\delta \in R/zR$ が存在するものが
存在する．
\end{lemma}

\begin{proof}
$R = k[x, y, z] \oplus E$ と置く．ここで上の加群 $E$ を平方零イデアルと
みなす．このとき $x, y, z$ が $R$ の正則列であり，
$\delta \in E/zE \subset R/zR$ が所望の性質をもつ元を与えることは明らかで
ある．局所的な例を得るには，$R$ を極大イデアル
$\mathfrak m = (x, y, z, E)$ で局所化すればよい．局所化は完全なので，
列 $x, y, z$ は正則列のままである．また $\delta$ は $\mathfrak m$ に
台をもつので，非零のままである．
\end{proof}

\begin{lemma}
\label{examples-lemma-base-change-regular-sequence}
局所環の局所準同型 $A \to B$ と，$B$ の極大イデアル内の正則列 $x, y$ で
あって，$B/(x, y)$ は $A$ 上平坦であるが，$B/\mathfrak m_AB$ における
$x, y$ の像 $\overline{x}, \overline{y}$ は正則列をなさず，Koszul 正則列
ですらないものが存在する．
\end{lemma}

\begin{proof}
$A = k[z]_{(z)}$ とし，$B = (k[x, y, z] \oplus E)_{(x, y, z, E)}$ と置く．
Lemma \ref{examples-lemma-strange-regular-sequence} の証明で見たように，
$x, y, z$ は $B$ の正則列である．したがって $x, y$ は $B$ の正則列で，
$B/(x, y)$ は捩れのない $A$-加群，ゆえに平坦である．一方，
$\overline{x}, \overline{y}$ で零化される非零元
$\delta \in B/\mathfrak m_AB = B/zB$ が存在する．したがって
$H_2(K_\bullet(B/\mathfrak m_AB, \overline{x}, \overline{y})) \not = 0$ である．
ゆえに $\overline{x}, \overline{y}$ は Koszul 正則でなく，特に正則列でない；
More on Algebra, Lemma
\ref{more-algebra-lemma-regular-koszul-regular} を参照せよ．
\end{proof}
% END EXPANDED MODULE ch110_13.tex
% BEGIN EXPANDED MODULE ch110_14.tex
\section{無限次元 Noether 環}
\label{examples-section-Noetherian-infinite-dimension}

\noindent
Algebra, Proposition \ref{algebra-proposition-dimension} で見たように，
Noether 局所環の次元は有限である．しかし無限次元の Noether 環も存在する．
\cite[Appendix, Example 1]{Nagata} を参照せよ．

\medskip\noindent
実際，$k$ を体とし，環
$$
R = k[x_1, x_2, x_3, \ldots ].
$$
を考える．$i = 1, 2, \ldots$ に対して
$\mathfrak p_i = (x_{2^{i - 1}}, x_{2^{i - 1} + 1}, \ldots, x_{2^i - 1})$
と置く．これらは $R$ の素イデアルである．$S$ を乗法的部分集合
$$
S = \bigcap\nolimits_{i \geq 1} (R \setminus \mathfrak p_i).
$$
とし，環 $A = S^{-1}R$ を考える．次を主張する．
\begin{enumerate}
\item 環 $A$ の極大イデアルは $\mathfrak m_i = \mathfrak p_iA$ である．
\item $A_{\mathfrak m_i} = R_{\mathfrak p_i}$ であり，これは次元 $2^i$ の
% P12-ERR-0324: see ../control/ERRATA_CANDIDATES.jsonl
Noether 局所環である．
\item 環 $A$ は Noether 環である．
\end{enumerate}
したがって，これが求める例であることは明らかである．詳細は省略する．
% END EXPANDED MODULE ch110_14.tex
% BEGIN EXPANDED MODULE ch110_15.tex
\section{完備化が非被約となる局所環}
\label{examples-section-local-completion-nonreduced}

\noindent
Algebra, Example \ref{algebra-example-bad-dvr-char-p} では，標数 $p$ の
次元 $1$ の Noether 局所整域 $R$ で，その完備化が非被約となる例を
与えた．この節では，標数零で同様の環を与える
\cite[Proposition 3.1]{Ferrand-Raynaud} の例を提示する．

\medskip\noindent
$\mathbf{C}\{x\}$ を複素数体 $\mathbf{C}$ 上の収束冪級数環とする．
すべての冪級数からなる環 $\mathbf{C}[[x]]$ はその完備化である．
$K = \mathbf{C}\{x\}[1/x]$ を収束 Laurent 級数体とする．
$\mathbf{C}$ 上の $K$ の代数的微分からなる $K$-加群
$\Omega_{K/\mathbf{C}}$ は無限次元 $K$-ベクトル空間である
（証明は省略する）．$n \geq 1$ に対して
$f_n \in x\mathbf{C}\{x\}$ を，
$
\text{d}x, \text{d}f_1, \text{d}f_2, \ldots
$
が $\Omega_{K/\mathbf{C}}$ の基底の一部となるように選べる．したがって
$\mathbf{C}$-導分
$$
D : \mathbf{C}\{x\} \longrightarrow \mathbf{C}((x))
$$
で，$D(x) = 0$ および
% P12-ERR-0325: see ../control/ERRATA_CANDIDATES.jsonl
$D(f_i) = x^{-n}$ を満たすものを見つけられる．
$$
A = \{f \in \mathbf{C}\{x\} \mid D(f) \in \mathbf{C}[[x]]\}
$$
と置く．次を主張する．
\begin{enumerate}
\item $\mathbf{C}\{x\}$ は $A$ 上整である．
\item $A$ は局所整域である．
\item $\dim(A) = 1$ である．
\item $A$ の極大イデアルは $x$ と $xf_1$ で生成される．
\item $A$ は Noether 環である．
\item $A$ の完備化は $\mathbf{C}[[x]]$ 上の双対数環に等しい．
\end{enumerate}
双対数環は非被約なので，環 $A$ が求める例を与える．

\medskip\noindent
$0 \not = f \in x\mathbf{C}\{x\}$ ならば，ある $n \geq 0$ と
$h \in \mathbf{C}[[x]]$ に対して $D(f) = h/f^n$ と書けることに注意せよ．
したがって $D(f^{n + 1}/(n + 1)) \in \mathbf{C}[[x]]$ および
$D(f^{n + 2}/(n + 2)) \in \mathbf{C}[[x]]$ である．ゆえに
$f^{n + 1}, f^{n + 2} \in A$ である！ 特に (1) が成り立つことが分かる．
また，$A$ の分数体は $\mathbf{C}\{x\}$ の分数体に等しいことも従う．
さらに $A \cap x\mathbf{C}\{x\}$ は $A$ の非単元全体であるから，
$A$ は次元 $1$ の局所整域である．(4) を示せば，$A$ が Noether 環である
ことが従う（証明は省略する）．$f \in A \cap x\mathbf{C}\{x\}$ とする．
$D(f) = h$，$h \in \mathbf{C}[[x]]$ と書き，さらに
$c \in \mathbf{C}$，$h' \in \mathbf{C}[[x]]$ として $h = c + xh'$ と書く．
すると $D(f - cxf_1) = c + xh' - c = xh'$ である．一方，
$g \in \mathbf{C}\{x\}$ として $f - cxf_1 = xg$ と書けるが，上の計算により
$D(g) = h' \in \mathbf{C}[[x]]$ である．したがって $g \in A$ であり，
所望の通り $f = cxf_1 + xg \in (x, xf_1)$ となる．

\medskip\noindent
最後に，なぜ $A$ の完備化は非被約なのだろうか．$A$ の完備化を
$\hat A$ と書く．もちろん，$x \in A$ なので，これは $\mathbf{C}\{x\}$ の
完備化 $\mathbf{C}[[x]]$ に全射的に写る．この写像を
$\psi : \hat A \to \mathbf{C}[[x]]$ と書く．上で
$\mathfrak m_A = (x, xf_1)$ であることを見た．したがって簡単な計算により
$D(\mathfrak m_A^n) \subset (x^{n - 1})$ である．ゆえに
$D : A \to \mathbf{C}[[x]]$ は連続であり，$\psi$ 上の連続導分
$\hat D : \hat A \to \mathbf{C}[[x]]$ を誘導する．したがって環準同型
$$
\psi + \epsilon \hat D :
\hat A
\longrightarrow
\mathbf{C}[[x]][\epsilon].
$$
を得る．$\hat A$ は一次元 Noether 完備局所環なので，この射が全射である
ことを示せば，$\hat A$ が非被約であることが従う．実際にはこの写像は
同型であるが，その確認は省略する．部分環
$\mathbf{C}[x]_{(x)} \subset A$ は完備化上の写像
$i : \mathbf{C}[[x]] \to \hat A$ を誘導し，
% P12-ERR-0326: see ../control/ERRATA_CANDIDATES.jsonl
$i \circ \psi = \text{id}$ および $D \circ i = 0$ を満たす
（構成により $D(x) = 0$ だからである）．元 $x^nf_n \in A$ を考える．
すべての $n \geq 1$ に対して
$$
(\psi + \epsilon D)(x^nf_n) = x^n f_n + \epsilon
$$
である．以上のことから全射性は容易に従う．
% END EXPANDED MODULE ch110_15.tex
% BEGIN EXPANDED MODULE ch110_16.tex
\section{完備化が非被約となるもう一つの局所環}
\label{examples-section-another-local-completion-nonreduced}

\noindent
この節では，単項イデアルに関して完備な次元 $2$ の Noether 局所整域で，
ある局所化を再完備化すると非被約になる例を構成する．

\medskip\noindent
$p$ を素数とする．$k$ を標数 $p$ の体で，$k$ は $p$ 乗元からなる部分体
$k^p$ 上の次数が無限であるものとする．例えば
$k = \mathbf{F}_p(t_1, t_2, t_3, \ldots)$ とすればよい．環
$$
A =
\left\{
\begin{matrix}
\sum a_{i, j} x^iy^j \in k[[x, y]] \text{ such that for all }n \geq 0
\text{ we have } \\
[k^p(a_{n, n}, a_{n, n + 1}, a_{n + 1, n},
a_{n, n + 2}, a_{n + 2, n}, \ldots) : k^p] < \infty
\end{matrix}
\right\}
$$
を考える．集合として
$$
k^p[[x, y]] \subset A  \subset k[[x, y]]
$$
である．$A$ の各元 $f$ は，一意的に級数
$$
f = f_0 + f_1 xy + f_2 (xy)^2 + f_3 (xy)^3 + \ldots
$$
と書ける．ここで
$$
f_n = a_{n, n} + a_{n, n + 1} y + a_{n + 1, n} x +
a_{n, n + 2} y^2 + a_{n + 2, n} x^2 + \ldots
$$
であり，$A$ を定義する式の条件は，$f_n$ の係数が $k^p$ の有限次拡大を
生成することを意味する．この表示から，$A$ は $k[[x, y]]$ の
$k^p[[x, y]]$-部分代数で，イデアル $xy$ に関して完備であることが明らかで
ある．さらに，明らかに
$$
A/xy A = C \times_k D
$$
である．ここで $k^p[[x]] \subset C \subset k[[x]]$ および
$k^p[[y]] \subset D \subset k[[y]]$ は，Algebra, Example
\ref{algebra-example-bad-dvr-char-p} の冪級数からなる部分環である．
したがって $C$ と $D$ は離散付値環であり，$A/ xy A$ は Noether 環である．
Algebra, Lemma \ref{algebra-lemma-completion-Noetherian} により，$A$ は
Noether 環である．$\dim(k[[x, y]]) = 2$ なので，Algebra, Lemma
\ref{algebra-lemma-integral-sub-dim-equal} を用いると
$\dim(A) = 2$ が従う．

\medskip\noindent
$f = \sum a_i x^i$ を，$k^p(a_0, a_1, a_2, \ldots)$ の $k^p$ 上の次数が
無限であるような冪級数とする．このとき $f \not \in A$ だが
$f^p \in A$ である．
$$
B = A[f] \subset k[[x, y]]
$$
と置く．$B$ は $A$ 上有限なので，$B$ は Noether 環である．また，
$B$ は $xy$ で生成されるイデアルに関して完備である；Algebra, Lemma
\ref{algebra-lemma-completion-tensor} を参照せよ．実際，$B$ は $A$ 上
基底 $1, f, f^2, \ldots, f^{p - 1}$ をもつ自由加群であるが，証明は省略する．

\medskip\noindent
環
$$
(B_y)^\wedge = (B[1/y])^\wedge = \lim B[1/y]/(xy)^n B[1/y] =
\lim B[1/y] / x^n B[1/y]
$$
は非被約であると主張する．実際，この環は
$$
(A_y)^\wedge = (A[1/y])^\wedge = \lim A[1/y]/(xy)^n A[1/y] =
\lim A[1/y] / x^n A[1/y]
$$
上，基底 $1, f, \ldots, f^{p - 1}$ をもつ自由加群である．しかし
$f^p = g^p$ となる元 $g \in (A_y)^\wedge$ が存在する．実際，$f$ に用いた
ものと同じ式 $g = \sum a_i x^i$ をとればよく，これは
$(A_y)^\wedge$ の元として意味をもつ．したがって
$$
(B_y)^\wedge = (A_y)^\wedge[f]/(f^p - g^p) \cong
(A_y)^\wedge[\tau]/(\tau^p)
$$
は非被約である．実際，この例はもう少し強いことを示す．
すなわち，$(A_y)^\wedge$ は一様化元 $x$ をもち，剰余体が上の離散付値環
$D$ の分数体である離散付値環であることに注意せよ．したがって
$$
(B_y)^\wedge[1/(xy)] = ((B_y)^\wedge)_{xy}
$$
ですら非被約である．これは次の種類の例を与える．

\begin{lemma}
\label{examples-lemma-nonreduced-recompletion}
単項イデアル $I = (b)$ に関して完備な局所 Noether $2$ 次元整域
$(B, \mathfrak m)$ と，$f \not \in I$ を満たす元
$f \in \mathfrak m$ であって，単項局所化 $B_f$ の $I$-進完備化
$C = (B_f)^\wedge$ が非被約であり，さらに
$C_b = C[1/b] = (B_f)^\wedge[1/b]$ も非被約となるものが存在する．
\end{lemma}

\begin{proof}
上の議論を参照せよ．
\end{proof}
% END EXPANDED MODULE ch110_16.tex
% BEGIN EXPANDED MODULE ch110_17.tex
\section{鎖状でない Noether 局所環}
\label{examples-section-non-catenary-Noetherian-local}

\noindent
Noether 局所環にはよく機能する次元論があるにもかかわらず，鎖状でない
Noether 局所環も存在する．一つの例は
\cite[Appendix, Example 2]{Nagata} にある．実際，ここではこの例の
最も簡単な場合を示す．すなわち，普遍鎖状でない次元 $2$ の局所 Noether
整域 $A$ を構成する．（$A$ は自動的に鎖状であることに注意せよ；Exercises,
Exercise \ref{exercises-exercise-Noetherian-local-domain-dim-2-catenary} を
参照せよ．）普遍鎖状でない Noether 局所環が存在すれば，鎖状でない
Noether 局所環も存在することになる．以下では，ここで扱う具体例について
節の最後にこのことを詳述する．

\medskip\noindent
$k$ を体とし，$k$ 上の一変数形式冪級数環 $k[[x]]$ を考える．
$$
z = \sum\nolimits_{i = 1}^\infty a_i x^i
$$
を形式冪級数とする．Laurent 級数体 $k((x)) = k[[x]][1/x]$ の元として
$z$ は $k(x)$ 上超越的であると仮定する．
$$
z_j
=
x^{-j}(z - \sum\nolimits_{i = 1, \ldots, j - 1} a_i x^i)
=
\sum\nolimits_{i = j}^\infty a_i x^{i - j}
\in k[[x]].
$$
と置く．$z = xz_1$ であることに注意せよ．$R$ を $x$，$z$，および
すべての $z_j$ で生成される $k[[x]]$ の部分環，すなわち
$$
R = k[x, z_1, z_2, z_3, \ldots ] \subset k[[x]].
$$
とする．$R$ のイデアル $\mathfrak m = (x)$ および
$\mathfrak n = (x - 1, z_1, z_2 + a_1, z_3 + a_1 + a_2, \ldots)$ を
考える．

\medskip\noindent
$xz_{j + 1} + a_j = z_j$ である．したがって $R/\mathfrak m = k$ であり，
$\mathfrak m$ は極大イデアルである．さらに，$R$ の
$\mathfrak m$ に属さない任意の元は $k[[x]]$ の単元に写るので，
$R_{\mathfrak m} \subset k[[x]]$ である．実際，
$R_{\mathfrak m}$ は離散付値環で，剰余体 $k$ をもつことが容易に従う．
% P12-ERR-0327: see ../control/ERRATA_CANDIDATES.jsonl

\medskip\noindent
次を主張する．
$$
R/(x - 1) =
k[x, z_1, z_2, z_3, \ldots ]/(x - 1)
\cong
k[z].
$$
実際，上の関係から
$z_{j + 1} = z_j - a_j - (x - 1)z_{j + 1}$ が従うので，商
$R/(x - 1)$ における $z_{j + 1}$ の類を $z_j$ を用いて表せる．構成により
$R$ の分数体は $k$ 上超越次数 $2$ をもつから，$R/(x - 1)$ において
$z$ は $k$ 上超越的であり，所望の同型が得られる．したがって
$\mathfrak n = (x - 1, z)$ は極大イデアルである．実際，写像
$$
k[x, x^{-1}, z]_{(x - 1, z)} \longrightarrow R_{\mathfrak n}
$$
は同型である（$R_{\mathfrak n}$ では $x^{-1}$ が可逆であり，かつ
$z_{j + 1} = x^{-1}z_j - a_j = \ldots = f_j(x, x^{-1}, z)$ だからである）．
これにより，$R_{\mathfrak n}$ は剰余体 $k$ をもつ次元 $2$ の
正則局所環であることが分かる．

\medskip\noindent
$S$ を乗法的部分集合
$$
S =
(R \setminus \mathfrak m) \cap (R \setminus \mathfrak n) =
R \setminus (\mathfrak m \cup \mathfrak n)
$$
とし，$B = S^{-1}R$ と置く．次を主張する．
\begin{enumerate}
\item 環 $B$ は $k$-代数である．
\item 環 $B$ の極大イデアルは $\mathfrak mB$ と $\mathfrak nB$ の二つである．
\item これらの極大イデアルにおける剰余体は $k$ である．
\item $B_{\mathfrak mB} = R_{\mathfrak m}$ および
$B_{\mathfrak nB} = R_{\mathfrak n}$ であり，これらはそれぞれ次元
$1$ および $2$ の Noether 正則局所環である．
\item 環 $B$ は Noether 環である．
\end{enumerate}
確認の詳細は省略する．

\medskip\noindent
上記の性質をもつ $k$-代数 $B$ が与えられれば，次のように例を得る．
$\text{rad}(B) = \mathfrak mB \cap \mathfrak nB$ を Jacobson 根基として，
$A = k + \text{rad}(B) \subset B$ と置く．$B$ は $A$ 上有限であることが
容易に分かるので，Eakin の定理により $A$ は Noether 環である
（\cite{Eakin}，\cite[Appendix A1]{Nagata}，または，将来の参考文献を
% P12-ERR-0328: see ../control/ERRATA_CANDIDATES.jsonl
ここに挿入せよ，を参照せよ）．さらに $A$ は $B$ と同じ分数体と剰余体
$k$ をもつ局所整域である．$B$ の次元は $2$ なので，Algebra, Lemma
\ref{algebra-lemma-integral-sub-dim-equal} により $A$ の次元も $2$ である．

\medskip\noindent
$A$ が普遍鎖状ならば，次元公式 Algebra, Lemma
\ref{algebra-lemma-dimension-formula} により
$\dim(B_{\mathfrak mB}) = 2$ となるが，これは矛盾である．

\medskip\noindent
$B$ は $A$ 上一元で生成されることに注意せよ．したがって，$A[x]$ の
ある素イデアル $\mathfrak p$ に対して $B = A[x]/\mathfrak p$ である．
$\mathfrak m' \subset A[x]$ を $\mathfrak mB$ に対応する極大イデアルとする．
一方では $\dim(A[x]_{\mathfrak m'}) = 3$ であり，他方では
$$
(0)
\subset \mathfrak pA[x]_{\mathfrak m'}
\subset \mathfrak m'A[x]_{\mathfrak m'}
$$
が素イデアルの極大鎖である．したがって $A[x]_{\mathfrak m'}$ は
鎖状でない Noether 局所環の例である．
% END EXPANDED MODULE ch110_17.tex
% BEGIN EXPANDED MODULE ch110_18.tex
\section{悪い局所 Noether 環の存在}
\label{examples-section-bad}

\noindent
$(A, \mathfrak m, \kappa)$ を Noether 完備局所環とする．\cite{Lech} では，この $A$ について，
$\text{depth}(A) \geq 1$ であり，$A$ が部分環として $\mathbf{Q}$ または
$\mathbf{F}_p$ を含むか，あるいは部分環として $\mathbf{Z}$ を含み，
$\mathbf{Z}$-加群として捩れをもたないならば，$A$ は Noether 局所整域の
完備化であることが示された．これは「奇妙な」性質をもつ Noether
局所整域の多くの例を生み出す．

\medskip\noindent
例えばこれを $A = \mathbf{C}[[x, y]]/(y^2)$ に適用すると，完備化が
非被約となる Noether 局所整域が得られる．Section
\ref{examples-section-local-completion-nonreduced} と比較せよ．

\medskip\noindent
\cite{LLPY} では，$A$ が被約な局所 Noether 環の完備化となるための
必要十分条件が見いだされた．

\medskip\noindent
\cite{Heitmann-completion-UFD} では，この $A$ について，$\text{depth}(A) \geq 2$ であり，
$A$ が部分環として $\mathbf{Q}$ または $\mathbf{F}_p$ を含むか，あるいは
部分環として $\mathbf{Z}$ を含み，$\mathbf{Z}$-加群として捩れをもたない
ならば，$A$ は局所 Noether UFD $R$ の完備化であることが示された．
特に $R$ は正規であり（Algebra, Lemma \ref{algebra-lemma-UFD-normal}），
したがって $R$ の Hensel 化も正規整域である
（More on Algebra, Lemma \ref{more-algebra-lemma-henselization-normal}）．
ゆえに上の $A$ は Hensel 的 Noether 局所正規整域の完備化である
（$R$ とその Hensel 化の完備化は一致するからである；More on Algebra,
Lemma \ref{more-algebra-lemma-henselization-noetherian} を参照せよ）．

\medskip\noindent
これを適用して，
$R^\wedge \cong \mathbf{C}[[x, y, z, w]]/(wx, wy)$ を満たす Noether
局所 UFD $R$ を得る．$\Spec(R^\wedge)$ は，正則な $2$ 次元成分と
正則な $3$ 次元成分の和であることに注意せよ．環 $R$ は普遍鎖状で
ありえない．極大イデアルのブローアップ
$$
X \longrightarrow \Spec(R)
$$
を考える．このとき $X$ は整スキームである．
$\dim(\mathcal{O}_{X, x}) = 2$ となる閉点 $x \in X$ が存在する．実際，
完備局所環のレベルでは，$x$ を $2$ 次元成分の狭義変換上にあり，
$3$ 次元成分の狭義変換上にはないように選ぶ．Morphisms, Lemma
\ref{morphisms-lemma-dimension-formula} により，$R$ が普遍鎖状でないことが
分かる．Section \ref{examples-section-non-catenary-Noetherian-local} と比較せよ．

\medskip\noindent
上の環 $R$ は鎖状である（$3$ 次元局所 Noether UFD だからである）．しかし
\cite{Ogoma-example} では，完備化
$R^\wedge \cong \mathbf{C}[[x, y, z, w]]/(wx, wy)$ をもち，鎖状でない
正規局所 Noether 整域 $R$ が構成されている．\cite{Heitmann-Ogoma} および
\cite{Lech-YAPO} も参照せよ．

\medskip\noindent
\cite{Heitmann-isolated} では，$A$ が部分環として $\mathbf{Q}$ または
$\mathbf{F}_p$ を含むか，あるいは $A$ の剰余標数が $p > 0$ で，$p$ が
$A$ のいかなる真の局所化においても非零零因子に写らないならば，$A$ は
孤立特異点をもつ局所 Noether 環 $R$ の完備化であることが示された．
ここで Noether 局所環 $R$ が孤立特異点をもつとは，すべての非極大
素イデアル $\mathfrak p \subset R$ に対して $R_\mathfrak p$ が正則局所環で
あることをいう．

\medskip\noindent
論文 \cite{Nishimura-few} と \cite{Nishimura-few-II} には，指定された完備化を
もつ「悪い」Noether 局所環の長い一覧が含まれている．特に，完備化が
% P12-ERR-0329: see ../control/ERRATA_CANDIDATES.jsonl
$\mathbf{C}[[x, y, z]]/(yz)$ である $2$ 次元 Nagata 局所正規整域の例と，
完備化が $\mathbf{C}[[x, y, z]]/(y^2 - z^3)$ である例が構成されている．

\medskip\noindent
なお，\cite{Loepp} では，$A$ が被約かつ等次元で，$A$ のどの整数も
零因子でないならば，$A$ は優秀 Noether 局所整域の完備化であることが
示された．しかし，優秀環が得られるので，これは「悪い」Noether 局所環を
与えない！
% END EXPANDED MODULE ch110_18.tex
% BEGIN EXPANDED MODULE ch110_19.tex
\section{Noether Jacobson 環における次元}
\label{examples-section-noetherian-jacobson}

\noindent
$k$ を有限体の代数閉包とする．$A = k[x, y]$ と置き，$X = \Spec(A)$ とする．
$C = V(x)$ を $y$-軸とする（$X$ の他の任意の $1$ 次元整閉部分スキームを
選んでもよい）．$C_1, C_2, C_3, \ldots$ を，$X$ の残りの $1$ 次元
整閉部分スキームの一つの列挙とする．$p_1, p_2, p_3, \ldots$ を $C$ の
閉点の一つの列挙とする．

\medskip\noindent
主張：任意の $n$ に対し，集合論的に
$$
\{p_n\} = Z_n \cap (C \cup C_1 \cup \ldots \cup C_n)
$$
となる次元 $1$ の既約閉集合 $Z_n \subset X$ が存在する．これを示すため，
$Y = C \cup C_1 \cup C_2 \cup \ldots \cup C_n$ と置く．これは $k$ 上の
次元 $1$ の被約アフィン代数スキームである．集合論的に
$V(f) \cap  Y = \{p_n\}$ となる $f \in k[x, y]$ を見つければ十分である．
実際，そのとき $Z_n$ として $V(f)$ の適当な既約成分を取れる．制限写像
$$
k[x, y] \longrightarrow \Gamma(Y, \mathcal{O}_Y)
$$
は全射なので，零点集合が集合論的に $\{p_n\}$ となる $Y$ 上の正則関数
$g$ を見つければ十分である．これが可能であることを見るため，台が $p_n$
である有効 Cartier 因子 $D \subset Y$ を選ぶ（これは Varieties, Lemma
\ref{varieties-lemma-complement-codim-1-closed-points} により可能である）．
したがって，ある $N > 0$ に対して
$\mathcal{O}_X(ND) \cong \mathcal{O}_X$ であることを示せば十分である．
% P12-ERR-0330: see ../control/ERRATA_CANDIDATES.jsonl
しかし，有限体の代数閉包上のアフィン $1$ 次元代数スキームの Picard 群は
捩れ群であり（将来の参考文献をここに挿入せよ），主張が従う．
% P12-ERR-0331: see ../control/ERRATA_CANDIDATES.jsonl

\medskip\noindent
すべての $n$ に対して上のような $Z_n$ を選ぶ．$k[x, y]$ は UFD なので，
ある既約元 $f_n \in A$ により $Z_n = V(f_n)$ と書ける．$S \subset k[x, y]$
を $f_1, f_2, f_3, \ldots$ で生成される乗法的部分集合とする．Noether 環
$B = S^{-1}A$ を考える．

\medskip\noindent
明らかに，環準同型 $A \to B$ は局所環を同一視し，単射
$\Spec(B) \to \Spec(A)$ を誘導する．さらに曲線 $C_1$ を見ると，
$\Spec(A)$ から $\Spec(B)$ へ移る際に除かれるのは $C \cap C_1$ の点だけで
あることが分かる．特に，$\Spec(B)$ には $A$ の極大イデアルに対応する
極大イデアルが無限個ある．一方，$C = V(x)$ のすべての閉点（ただし
生成点は除く）を取り除いたので，$B/xB$ のスペクトルはただ一つの
素イデアルからなり，したがって $xB$ は極大イデアルである．最後に，
$i \geq 1$ に対して曲線 $C_i$ を考える．ある既約元 $g_i \in A$ により
$C_i = V(g_i)$ と書く．ある $n$ に対して $C_i = Z_n$ ならば，$g_iB$ は
単位イデアルである．そうでなければ，$A$ から $B$ へ移っても $C_i$ の
有限個を除くすべての閉点が残る．実際，$C_i$ から除かれるのは
$(Z_1 \cup \ldots \cup Z_{i - 1} \cup C) \cap C_i$ の点だけである．

\medskip\noindent
上で記述した $B$ の素スペクトルの構造と Algebra, Lemma
\ref{algebra-lemma-noetherian-dim-1-Jacobson} により，$B$ は Jacobson 環で
ある．その極大イデアルは，$\Spec(B)$ に属する $A$ の極大イデアル
（そのようなものは無限個ある）と，極大イデアル $xB$ である．
% P12-ERR-0332: see ../control/ERRATA_CANDIDATES.jsonl
したがって，次元 $1$ および $2$ の局所環が存在することが分かる．

\begin{lemma}
\label{examples-lemma-Noetherian-Jacobson}
極大イデアル $\mathfrak m_1, \mathfrak m_2$ をもち，
$\dim(B_{\mathfrak m_1}) = 1$ および $\dim(B_{\mathfrak m_2}) = 2$ を満たす
Jacobson 普遍鎖状 Noether 整域 $B$ が存在する．
\end{lemma}

\begin{proof}
$B$ の構成は上で与えた．Algebra, Lemma
\ref{algebra-lemma-localization-catenary} および Morphisms, Lemma
\ref{morphisms-lemma-ubiquity-uc} により $B$ が普遍鎖状であることだけを
指摘しておく．
\end{proof}
% END EXPANDED MODULE ch110_19.tex
% BEGIN EXPANDED MODULE ch110_20.tex
\section{基礎空間は Noether だが Noether でないスキーム}
\label{examples-section-noetherian-not-noetherian}

\noindent
基礎位相空間が Noether であっても Noether スキームではないことがあると
示す二つの例を与える．

\begin{example}
\label{examples-example-many-variables}
$k$ を体とし，$A = k[x_1, x_2, x_3, \dots] /
(x_1^2, x_2^2, x_3^2, \dots)$ と置く．$A$ の任意の素イデアルは冪零元
$x_1, x_2, x_3, \dots$ を含むので，$\mathfrak p = (x_1, x_2, x_3,
\dots)$ が $A$ の唯一の素イデアルである．したがって
$\operatorname{Spec} A$ の基礎位相空間は一点からなり，特に Noether である．
しかし，$\mathfrak p$ が有限生成でないことは明らかである．
\end{example}

\begin{example}
\label{examples-example-many-mononials}
$k$ を体とし，$A \subseteq k[x, y]$ を $k$ と単項式
$\{xy^i\}_{i \ge 0}$ で生成される部分環とする．$x$ を含まない $A$ の
素イデアルは，$A_x \cong k[x, x^{-1}, y]$ の素イデアルと一対一に対応する．
一方，素イデアル $\mathfrak p$ が $x$ を含むならば，任意の $i \ge 0$ に
対して $(xy^i)^2 = x(xy^{2i}) \in \mathfrak p$ であり，$\mathfrak p$ は
根基的なので，すべての $xy^i$ を含む．したがって
$\mathfrak p = (\{xy^i\}_{i \ge 0})$ である．ゆえに
$\operatorname{Spec} A$ の基礎位相空間は，Noether スキーム
$\Spec(A[x, x^{-1}, y])$ の点と素イデアル $\mathfrak p$ からなるので
Noether である．しかし環 $A$ は Noether でない．実際，$\mathfrak p$ は
有限生成でない．この例では，$A$ は整域でもあることに注意せよ．
\end{example}
% END EXPANDED MODULE ch110_20.tex
% BEGIN EXPANDED MODULE ch110_21.tex
\section{準アフィンな正規化をもつ非準アフィン多様体}
\label{examples-section-nonquasi-affine}

\noindent
この種の例が存在することは \cite[II Remark 6.6.13]{EGA} で述べられている．
そこではそのような例について EGA 第五巻を参照しているが，第五巻は刊行
されなかった．

\medskip\noindent
$k$ を体とする．$Y = \mathbf{A}^2_k \setminus \{(0, 0)\}$ と置く．$k$ 上の
多様体 $X$ で準アフィンでないものと，有限全射双有理射
$\pi : Y \longrightarrow X$ を構成する．そのため，$Y$ 内の次の曲線を考える．
$$
\begin{matrix}
C_1 & : & x = 0 \\
C_2 & : & y = 0
\end{matrix}
$$
$C_1 \cap C_2 = \emptyset$ であることに注意せよ．同型
$\varphi : C_1 \to C_2$，$(0, y) \mapsto (y^{-1}, 0)$ を選ぶ．次の図式が
多様体の圏（さらにはスキームの圏）における余等化子図式となるような上記の
射 $\pi : Y \to X$ が一意に存在すると主張する．
$$
\xymatrix{
C_1
\ar@<1ex>[rr]^{\text{id}} \ar@<-1ex>[rr]_{\varphi}
& &
Y \ar[r]^\pi & X
}
$$
ひとまずこれを認め，そのような $X$ が準アフィンではありえないことを示す．
実際，明らかに
$$
\Gamma(X, \mathcal{O}_X) =
\{ f \in k[x, y] \mid f(0, y) = f(y^{-1}, 0)\} =
k \oplus (xy) \subset k[x, y].
$$
となる．特に，これらの関数は点 $(1, 0)$ と $(-1, 0)$ を分離しないが，
$X$ におけるそれらの像は（下で示すように）異なる（$k$ の標数が $2$ で
なければ）．

\medskip\noindent
$X$ の存在を示すため，$Y$ の Zariski 開集合 $D(x + y) \subset Y$ を考える．
これは環 $k[x, y, 1/(x + y)]$ のスペクトルであり，曲線 $C_1$，$C_2$ は
$D(x + y)$ に完全に含まれる．さらに射
$$
C_1 \amalg C_2
\longrightarrow
D(x + y) \cap Y = \Spec(k[x, y, 1/(x + y)])
$$
は閉埋入である．More on Algebra, Lemma
\ref{more-algebra-lemma-fibre-product-finite-type} により，環
$$
A =
\{f \in k[x, y, 1/(x + y)] \mid f(0, y) = f(y^{-1}, 0)\}
$$
は $k$ 上有限型である．一方，曲線 $C_1$，$C_2$ と交わらない $Y$ の開集合
$D(xy) \subset Y$ がある．これは環
$$
B = k[x, y, 1/xy].
$$
のスペクトルである．$A_{xy} \cong B_{x + y}$ であることに注意せよ
（$A$ は任意の多項式 $P$ に対する元 $xyP(x, y)$ と，元
$xy/(x + y)$ を明らかに含む）．
% P12-ERR-0333: see ../control/ERRATA_CANDIDATES.jsonl
アフィンスキーム $\Spec(A)$ と $\Spec(B)$ を同型
$A_{xy} \cong B_{x + y}$ により貼り合わせて $X$ が得られるので，$X$ は
明らかに $k$ 上有限型である．これが分離的であることを見るには，環準同型
$A \otimes_k B \to B_{x + y}$ が全射であることを示さなければならない．
これには，$A \otimes_k B$ が元 $xy/(x + y) \otimes 1/xy$ を含み，それが
$1/(x + y)$ に写ることを用いる．射 $Y \to X$ は自然な写像
$D(x + y) \to \Spec(A)$ および $D(xy) \to \Spec(B)$ によって与えられる．
これらはいずれも有限なので，所望のとおり $Y \to X$ は有限である．
$X$ が上に表示した図式の余等化子であることの確認は省略するが，
（スキームの圏における押し出しの将来の参考文献をここに挿入せよ）を
参照せよ．
% P12-ERR-0334: see ../control/ERRATA_CANDIDATES.jsonl
射 $\pi : Y \to X$ は点 $(1, 0)$ と $(-1, 0)$ を $X$ の異なる点へ写すことに
注意せよ．実際，関数 $(x + y^3)/(x + y)^2 \in A$ のそれらにおける値は
それぞれ $1/1$ と\ $-1/(-1)^2 = -1$ であり，常に異なる（標数が $2$ の
場合を除く；標数 $2$ については各自で点を見つけよ）．以上を補題として
まとめる．

\begin{lemma}
\label{examples-lemma-quasi-affine-normalization-not-quasi-affine}
$k$ を体とする．正規化が準アフィンであるが，自身は準アフィンでない
多様体 $X$ が存在する．
\end{lemma}

\begin{proof}
上の議論および（正規化に関する将来の参考文献をここに挿入せよ）を参照せよ．
% P12-ERR-0335: see ../control/ERRATA_CANDIDATES.jsonl
\end{proof}
% END EXPANDED MODULE ch110_21.tex
% BEGIN EXPANDED MODULE ch110_22.tex
\section{スキーム論的像を取ること}
\label{examples-section-scheme-theoretic-image}

\noindent
$k$ を体とし，$t$ を変数とする．$Y = \Spec(k[t])$ および
$X = \coprod_{n \geq 1} \Spec(k[t]/(t^n))$ と置く．各 $n \geq 1$ に対する
閉埋入 $\Spec(k[t]/(t^n)) \to \Spec(k[t])$ を用いて定まる射を
$f : X \to Y$ と書く．この場合，次が成り立つ．
\begin{enumerate}
\item $f$ のスキーム論的像（Morphisms, Definition
\ref{morphisms-definition-scheme-theoretic-image}）は $Y$ である．一方，
$f$ の像は $Y$ の閉点 $t = 0$ である．したがって，$f$ のスキーム論的像の
基礎閉集合は $f$ の像の閉包と等しくない．
\item スキーム論的像を形成する操作は，開部分スキーム
$V = \Spec(k[t, 1/t]) \subset Y$ への制限と可換でない．実際，$X$ における
$V$ の逆像は空であるから，$f|_{f^{-1}(V)} : f^{-1}(V) \to V$ の
スキーム論的像は空スキームである．これは $Y \cap V$ と等しくない．
\end{enumerate}
% END EXPANDED MODULE ch110_22.tex
% BEGIN EXPANDED MODULE ch110_23.tex
\section{局所閉部分集合の像}
\label{examples-section-non-chevalley}

\noindent
Chevalley の定理によれば，アフィンスキームの有限表示射による構成可能集合の
像は構成可能である．Algebra, Theorem \ref{algebra-theorem-chevalley} および
Morphisms, Section \ref{morphisms-section-constructible} を参照せよ．
局所閉部分集合の像については同じことが成り立たないことを見る．

\medskip\noindent
$k$ を標数 $0$ の体とする．射影射
$$
f :
X = \Spec(k[t, x_1, x_2, \ldots, y_1, y_2, \ldots])
\longrightarrow
\Spec(k[x_1, x_2, \ldots, y_1, y_2, \ldots]) = Y
$$
を考える．これは有限表示射である．$Z$ を，次の方程式で定義される $X$ の
閉部分集合とする．
$$
x_1(t - 1) = 0,\quad
x_2(t - 1)(t - 2) = 0,\quad
x_3(t - 1)(t - 2)(t - 3) = 0,\quad \ldots
$$
$U = \bigcup_{j \geq 1} U_j$ を，次により定義される $X$ の開集合とする．
$$
U_j = \text{次を満たす点： }y_j(t - 1)(t - 2) ... (t - j)\text{ は零でない}
$$
すると
$$
f(Z \cap U_j) =
\text{次を満たす点： }x_1, \ldots, x_j\text{ は零で， }y_j\text{ は零でない}
$$
である．$B = f(Z \cap U) = \bigcup_{j \geq 1} f(Z \cap U_j)$ は $Y$ の
局所閉部分集合の有限和ではないと主張する．

\medskip\noindent
主張の証明．$B = A_1 \cup \cdots \cup A_m$ が $Y$ の局所閉部分集合による
$B$ の有限被覆であると仮定する．$n$ に関する帰納法で $m \geq n$ を示す．
$B$ は空でないから，基底の場合 $n = 1$ はよい．$n > 1$ とし，$n - 1$ に
対する帰納法の仮定が成り立つとする．$B$ の閉包は $(x_1 = 0)$ なので，
ある $A_i$ は $(x_1 = 0)$ の空でない開部分集合を含まなければならない．
すると $A_i$ は $(x_1 = 0)$ において開でなければならない．しかし，
そのような開部分集合は $y_1 = 0$ を満たす点を含みえない．実際，$B$ の
点について $y_1 = 0$ ならば $x_2 = 0$ であり，これは $B$ が
$(x_1 = 0)$ 内で $(x, y)$ の近傍を含まないことを示す．したがって，残る
$m - 1$ 個の集合は $B \cap (y_1 = 0)$ の構成可能被覆を与える．ところが，
右シフト写像 $x_i \mapsto x_{i + 1}$，$y_i \mapsto y_{i + 1}$ は $B$ と
$B \cap (y_1 = 0)$ を同一視することに注意せよ．ゆえに帰納法の仮定から
$m - 1 \geq n - 1$ であり，$m \geq n$ を得る．
% P12-ERR-0336: see ../control/ERRATA_CANDIDATES.jsonl
これで帰納段階の証明，したがって主張の証明が終わる．

\begin{lemma}
\label{examples-lemma-no-chevalley}
アフィンスキーム間の有限表示射 $f : X \to Y$ と，$X$ の局所閉部分集合
$T$ で，$f(T)$ が $Y$ の局所閉部分集合の有限和でないものが存在する．
\end{lemma}

\begin{proof}
上の議論を参照せよ．
\end{proof}
% END EXPANDED MODULE ch110_23.tex
% BEGIN EXPANDED MODULE ch110_24.tex
\section{閉部分スキームの中で開でない局所閉部分スキーム}
\label{examples-section-strange-immersion}

\noindent
これは Morphisms, Example \ref{morphisms-example-thibaut} の複製である．
開埋入に続く閉埋入として合成表示できない埋入の例を与える．$k$ を体とする．
$X = \Spec(k[x_1, x_2, x_3, \ldots])$ とし，
$U = \bigcup_{n = 1}^{\infty} D(x_n)$ と置く．すると $U \to X$ は開埋入である．
イデアル
$$
I_n =
(x_1^n, x_2^n, \ldots, x_{n - 1}^n, x_n - 1, x_{n + 1}, x_{n + 2}, \ldots)
\subset
k[x_1, x_2, x_3, \ldots][1/x_n].
$$
を考える．任意の $m \not = n$ に対して
$I_n k[x_1, x_2, x_3, \ldots][1/x_nx_m] = (1)$ であることに注意せよ．
したがって，$D(x_n)$ 上の準連接イデアル $\widetilde I_n$ は $D(x_nx_m)$ 上で
一致する．すなわち，$n \not = m$ ならば
$\widetilde I_n|_{D(x_nx_m)} = \mathcal{O}_{D(x_n x_m)}$ である．ゆえに，
これらのイデアルは準連接イデアル層 $\mathcal{I} \subset \mathcal{O}_U$ に
貼り合わさる．$\mathcal{I}$ に対応する閉部分スキームを $Z \subset U$ と
する．このとき $Z \to X$ は埋入である．

\medskip\noindent
$Z \to X$ を $Z \to \overline{Z} \to X$ と分解することはできないと主張
する．ここで $\overline{Z} \to X$ は閉で，$Z \to \overline{Z}$ は開である．
実際，$\overline{Z}$ は，
$I_n = I k[x_1, x_2, x_3, \ldots][1/x_n]$ を満たすイデアル
$I \subset k[x_1, x_2, x_3, \ldots]$ により定義されなければならない．
しかし，すべての $I_n$ に入る $f \in k[x_1, x_2, x_3, \ldots]$ は $0$ だけ
である．したがって，そのような $I$ は存在しない．

\medskip\noindent
射 $Z \to X$ は，埋入のスキーム論的像が悪い振る舞いをする例も与える．
実際，上の議論から，埋入 $Z \to X$ のスキーム論的像は $X$ である．一方，
次が分かる．
\begin{enumerate}
\item $Z$ は $X$ に位相的に稠密でない．
\item $Z = Z \cap U \to U$ のスキーム論的像は $Z$ 自身である．これは
$U \cap X = U$ と等しくない．したがって，この場合スキーム論的像を形成する
操作は開集合への制限と可換でない．
\end{enumerate}
% END EXPANDED MODULE ch110_24.tex
% BEGIN EXPANDED MODULE ch110_25.tex
\section{適切な開集合の非存在}
\label{examples-section-nonexistence-opens}

\noindent
本節は Properties, Section
\ref{properties-section-finding-affine-opens} の結果を補足する．

\medskip\noindent
$k$ を体とし，$A = k[z_1, z_2, z_3, \ldots]/I$ とする．ここで $I$ は，
すべての相異なる二元の積 $z_iz_j$，$i \not = j$，$i, j \in \mathbf{N}$ で
生成されるイデアルである．$S = \Spec(A)$ と置く．$s \in S$ を極大イデアル
$(z_i)$ に対応する閉点とする．$S \setminus \{s\}$ において稠密な準コンパクト
開集合 $V \subset S \setminus \{s\}$ は存在しないと主張する．
$S \setminus \{s\} = \bigcup D(z_i)$ であることに注意せよ．各 $D(z_i)$ は，
生成点 $\eta_i$ をもつ開既約集合である．したがって，すべての $i$ に対して
$\eta_i \in V$ でなければならない．しかし，$S \setminus \{s\}$ の主アフィン
開集合は，ある $f \in (z_1, z_2, \ldots)$ に対する $D(f)$ の形である．ある
$n$ に対して $f \in (z_1, \ldots, z_n)$ だから，$D(f)$ は点 $\eta_i$ を
有限個しか含まない．ゆえに，$V$ は準コンパクトではありえない．

\medskip\noindent
$k$ を体とし，$B = k[x, z_1, z_2, z_3, \ldots]/J$ とする．ここで $J$ は，
積 $xz_i$，$i \in \mathbf{N}$，およびすべての相異なる二元の積 $z_iz_j$，
$i \not = j$，$i, j \in \mathbf{N}$ で生成されるイデアルである．
$T = \Spec(B)$ と置き，主開集合 $U = D(x)$ を考える．$V \cap U = \emptyset$
かつ $V \cup U$ が $T$ に稠密となる準コンパクト開集合 $V \subset T$ は
存在しないと主張する．$t \in T$ を極大イデアル $(x, z_i)$ に対応する閉点と
する．$T$ における $U$ の閉包は $\overline{U} = U \cup \{t\}$ である．
したがって，$V \subset \bigcup_i D(z_i)$ は準コンパクト開集合である．前段落の
議論により，$V$ は $\bigcup D(z_i)$ に稠密ではありえない．

\begin{lemma}
\label{examples-lemma-complement-of-affine-does-not-contain-qc-dense-open}
アフィンスキームにおける準コンパクト開集合の非存在：
% P12-ERR-0337: see ../control/ERRATA_CANDIDATES.jsonl
\begin{enumerate}
\item アフィンスキーム $S$ とアフィン開集合 $U \subset S$ で，
$U \cap V = \emptyset$ かつ $U \cup V$ が $S$ に稠密となる準コンパクト開集合
$V \subset S$ が存在しないものがある．
\item アフィンスキーム $S$ と閉点 $s \in S$ で，$S \setminus \{s\}$ が
準コンパクト稠密開集合を含まないものがある．
\end{enumerate}
\end{lemma}

\begin{proof}
上の議論を参照せよ．
\end{proof}

\noindent
% END EXPANDED MODULE ch110_25.tex
% BEGIN EXPANDED MODULE ch110_26.tex
$X$ を，アフィン開集合 $U$ に沿ってアフィンスキーム $T$（上を参照）の二つの
コピーを貼り合わせたものとする．したがって，射 $\pi : X \to T$ と
$X = U_1 \cup U_2$ があり，$\pi$ は $U_i$ を $T$ へ同型に写し，
$U_1 \cap U_2$ を $U$ へ同型に写す．$X$ は準分離的であり（Schemes, Lemma
\ref{schemes-lemma-characterize-quasi-separated}），準コンパクトであることに
注意せよ．分離的で稠密な準コンパクト開集合 $W \subset X$ は存在しないと
主張する．実際，上で導入した閉点 $t \in T$ に写る二つの閉点
$x_1 \in U_1$，$x_2 \in U_2$ を考える．$U$ の（一意な）生成点 $\eta$ に
写る生成点を $\tilde \eta \in U_1 \cap U_2$ とする．特殊化
$\eta \leadsto t$ の上で，$\tilde\eta \leadsto x_1$ および
$\tilde\eta \leadsto x_2$ であることに注意せよ．$\pi|_W : W \to T$ は
分離的なので，$x_1$ がその開集合に属し，かつ $x_2 \in W$ であることはない
（分離性の付値判定法 Schemes, Lemma
\ref{schemes-lemma-valuative-criterion-separatedness} による）．
$x_1 \not \in W$ としてよい．このとき $W \cap U_1$ は，$x_1$ を含まない
$U_1$ の準コンパクト（$X$ が準分離的だから）稠密開集合である．ここで，
$x$ を $z_1$ へ，$z_i$ を $z_{i + 1}$ へ写すことにより，スキームの同型
$(T, t) \cong (S, s)$ が存在することに注意せよ．したがって，本節の
第一段落により矛盾に至る．

\begin{lemma}
\label{examples-lemma-no-dense-separated-quasi-compact-open-in-qcqs}
分離的な準コンパクト稠密開集合を含まない準コンパクト準分離的スキーム
$X$ が存在する．
\end{lemma}

\begin{proof}
上の議論を参照せよ．
\end{proof}
% END EXPANDED MODULE ch110_26.tex
% BEGIN EXPANDED MODULE ch110_27.tex
\section{準コンパクト稠密開部分スキームの非存在}
\label{examples-section-nonexistence-qc-dense-open-subscheme}

\noindent
$X$ を体 $k$ 上の準コンパクト準分離的代数空間とする．スキーム軌跡
$X' \subset X$ が稠密開部分空間であることは知られている．Properties of
Spaces, Proposition
\ref{spaces-properties-proposition-locally-quasi-separated-open-dense-scheme}
を参照せよ．実際，この結果は $X$ が reasonable である場合にも成り立つ．
Decent Spaces, Proposition
\ref{decent-spaces-proposition-reasonable-open-dense-scheme} を参照せよ．
$X$ の準コンパクト稠密開部分スキームを見つけられるか，という自然な問いが
生じる．一般には不可能であることが分かる．

\medskip\noindent
$k$ の標数は 2 でないと仮定する．
$B = k[x, z_1, z_2, z_3, \ldots]/J$ とする．ここで $J$ は積 $xz_i$，
$i \in \mathbf{N}$，およびすべての相異なる二元の積 $z_iz_j$，
$i \not = j$，$i, j \in \mathbf{N}$ で生成されるイデアルである．
$U = \Spec(B)$ と置く．すべての座標が零である閉点を $0 \in U$ と書く．
$$
j : R = \Delta \amalg \Gamma \longrightarrow U \times_k U
$$
と置く．ここで $\Delta$ は $k$ 上の $U$ の対角射の像であり，
$$
\Gamma = \{((x, 0, 0, 0, \ldots), (-x, 0, 0, 0, \ldots))
\mid x \in \mathbf{A}^1_k, x \not = 0\}.
$$
明らかに $s, t : R \to U$ はエタールであり，したがって $j$ はエタール
同値関係である．商 $X = U/R$ は代数空間である（Spaces, Theorem
\ref{spaces-theorem-presentation}）．$(0, 0) \in \Delta$ は $\Gamma$ の
閉包に属するので，$j$ は埋入でないことに注意せよ．したがって $X$ は
スキームでない．一方，$R$ は準コンパクトなので $X$ は準分離的である．
点 $0 \in U$ の像を $0_X$ と書く．$X \setminus \{0_X\}$ はスキームであり，
実際，次の表示をもつと主張する．
$$
X \setminus \{0_X\} =
\Spec\left(k[x^2, x^{-2}]\right) \amalg
\Spec\left(k[z_1, z_2, z_3, \ldots]/(z_iz_j)\right) \setminus \{0\}
$$
はスキームである（詳細は省略する）．他方，Section
\ref{examples-section-nonexistence-opens} で，右辺のスキームは準コンパクト稠密開集合を
含まないことを見た．

\begin{lemma}
\label{examples-lemma-nonexistence-qc-dense-open-subscheme}
準コンパクト稠密開部分スキームを含まない準コンパクト準分離的代数空間が
存在する．
\end{lemma}

\begin{proof}
上の議論を参照せよ．
\end{proof}

\noindent
Spaces, Example \ref{spaces-example-non-representable-descent} の構成を，
上で Spaces, Example \ref{spaces-example-affine-line-involution} の構成を用いた
のと同じ仕方で用いると，準コンパクト稠密開部分スキームを含まない，
準コンパクト，準分離的かつ局所分離的な代数空間の例が得られる．
% END EXPANDED MODULE ch110_27.tex
% BEGIN EXPANDED MODULE ch110_28.tex
\section{代数空間上のアフィンスキーム}
\label{examples-section-embedding-affines}

\medskip\noindent
$f : Y \to X$ をスキームの射とし，$f$ は局所有限型で，$Y$ はアフィンで
あると仮定する．このとき，$X$ 上のアフィン $n$-空間への $Y$ の埋入
$Y \to \mathbf{A}^n_X$ が存在する．少し一般的な Morphisms, Lemma
\ref{morphisms-lemma-quasi-affine-finite-type-over-S} を参照せよ．

\medskip\noindent
次に，$f : Y \to X$ を代数空間の射とし，$f$ は局所有限型で，$Y$ は
アフィンスキームであると仮定する．このとき，一般には $X$ 上のアフィン
$n$-空間への $Y$ の埋入を見つけることはできない．

\medskip\noindent
第一の（厄介な）反例は $Y = \Spec(k)$ および
$X = [\mathbf{A}^1_k/\mathbf{Z}]$ である．ここで $k$ は標数零の体で，
$\mathbf{Z}$ は $\mathbf{A}^1_k$ に平行移動 $(n, t) \mapsto t + n$ で作用
する．実際，$X$ 上の任意の射 $Y \to \mathbf{A}^n_X$ を $X$ の被覆
$\mathbf{A}^1_k$ へ引き戻すと，$\mathbf{A}^{n + 1}_k$ へ写る
$\mathbf{A}^1_k$ の無限非交和が得られるが，これは埋入でない．
% P12-ERR-0338: see ../control/ERRATA_CANDIDATES.jsonl

\medskip\noindent
第二の反例は $Y = \mathbf{A}^1_k \to X = \mathbf{A}^1_k/R$ である．ここで
$R = \{(t, t)\} \amalg \{(t, -t), t \not = 0\}$ とする．実際，この場合
射 $Y \to \mathbf{A}^n_X$ は $Y$ 上の正則関数 $f_1, \ldots, f_n$ によって
与えられるので，被覆 $\mathbf{A}^{n + 1}_k \to \mathbf{A}^n_X$ と $Y$ の
ファイバー積はスキーム
$$
\{(f_1(t), \ldots, f_n(t), t)\} \amalg
\{(f_1(t), \ldots, f_n(t), -t), t \not = 0\}
$$
となり，$\mathbf{A}^{n + 1}_k$ への明らかな射をもつが，これは埋入でない．
これは $X$ が準分離的である反例を与えることに注意せよ．

\begin{lemma}
\label{examples-lemma-cannot-embed-into-affine}
$Y$ はアフィン，$X$ は準分離的であるような代数空間の有限型射
$Y \to X$ で，$X$ 上の埋入 $Y \to \mathbf{A}^n_X$ が存在しないものがある．
\end{lemma}

\begin{proof}
上の議論を参照せよ．
\end{proof}
% END EXPANDED MODULE ch110_28.tex
% BEGIN EXPANDED MODULE ch110_29.tex
\section{準連接加群の順像}
\label{examples-section-push-quasi-coherent}

\noindent
Schemes, Lemma \ref{schemes-lemma-push-forward-quasi-coherent} で，$f$ が
準コンパクトかつ準分離的ならば，$f_*$ は準連接加群を準連接加群へ写すことを
証明した．これら二つの条件がともに必要であることを示す例を与える．

\medskip\noindent
$Y = \Spec(A)$ をアフィンスキームとし，$X = \coprod_{n \in \mathbf{N}} Y$
と仮定する．明らかな射を $f : X \to Y$ とすると，
$f_*\mathcal{O}_X$ は準連接でないと主張する．実際，$a \in A$ に対して
$$
f_*\mathcal{O}_X(D(a)) = \prod\nolimits_{n \in \mathbf{N}} A_a
$$
である．したがって，$f_*\mathcal{O}_X$ が準連接であるためには，すべての
$a \in A$ に対して
$$
\prod\nolimits_{n \in \mathbf{N}} A_a
=
\left(\prod\nolimits_{n \in \mathbf{N}} A\right)_a
$$
でなければならない．これは一般には成り立たない．例えば，
$A = \mathbf{Z}$ かつ $a = 2$ ならば，$(1, 1/2, 1/4, 1/8, \ldots)$ は左辺の
元であるが右辺には属さない．$f$ は準コンパクトでない分離的射であることに
注意せよ．

\medskip\noindent
$k$ を体とする．
$$
A = k[t, z, x_1, x_2, x_3, \ldots]/(tx_1z, t^2x_2^2z, t^3x_3^3z, \ldots)
$$
と置く．$Y = \Spec(A)$ とし，$V \subset Y$ を開部分スキーム
$V = D(x_1) \cup D(x_2) \cup \ldots$ とする．$V$ に沿って $Y$ の二つの
コピーを貼り合わせたものを $X$ とし，明らかな射を $f : X \to Y$ とする．
このとき，完全列
$$
0 \to f_*\mathcal{O}_X \to
\mathcal{O}_Y \oplus \mathcal{O}_Y \xrightarrow{(1, -1)} j_*\mathcal{O}_V
$$
がある．ここで $j : V \to Y$ は包含射である．写像
$$
A \longrightarrow \prod A_{x_n}
$$
は単射なので（詳細は省略する），$\Gamma(Y, f_*\mathcal{O}_X) = A$ である．
一方，写像
$$
A_t \longrightarrow \prod A_{tx_n}
$$
の核は元 $z$ を含むので零でない．したがって，
$\Gamma(D(t), f_*\mathcal{O}_X)$ は $(z, 0)$ を含むから $A_t$ より真に大きい．
ゆえに，$f_*\mathcal{O}_X$ は準連接でない．$f$ は準コンパクトであるが
準分離的でないことに注意せよ．

\begin{lemma}
\label{examples-lemma-pushforward-quasi-coherent}
Schemes, Lemma \ref{schemes-lemma-push-forward-quasi-coherent} は，
準コンパクト性の仮定も準分離性の仮定も取り除けないという意味で鋭い．
\end{lemma}

\begin{proof}
上の議論を参照せよ．
\end{proof}
% END EXPANDED MODULE ch110_29.tex
% BEGIN EXPANDED MODULE ch110_30.tex
\section{すべての茎が階数 1 の有限自由加群である非有限加群}
\label{examples-section-nonfree}

\noindent
$R = \mathbf{Q}[x]$ とする．$R$ の分数体の部分加群として
$M = \sum_{n \in \mathbf{N}} \frac{1}{x - n}R$ と置く．このとき $M$ は
有限生成でないが，$R$ の任意の素イデアル $\mathfrak p$ に対して，
$R_{\mathfrak p}$-加群として $M_{\mathfrak p} \cong R_{\mathfrak p}$ である．

\medskip\noindent
同様な趣の例として $R = \mathbf{Z}$ および
$M = \sum_{p \text{ は素数}} \frac{1}{p} \mathbf{Z} \subset \mathbf{Q}$ がある．
これは $b$ が零でない平方因子をもたない整数である分数 $\frac{a}{b}$ の
全体に等しい．
% END EXPANDED MODULE ch110_30.tex
% BEGIN EXPANDED MODULE ch110_31.tex
\section{各茎では可逆であるが可逆でないイデアル}
\label{examples-section-locally-invertible-not-invertible}

\noindent
$A$ を整域とし，$I \subset A$ を非零イデアルとする．$I$ が可逆であるという
とき，$I$ が $A$-加群として可逆であることを意味することを思い出そう．
$I$ は可逆でないが，任意の素イデアル $\mathfrak q \subset A$ に対して
$I_\mathfrak q = (f) \subset A_\mathfrak q$ が（非零）単項イデアルとなる例を
構成する．文献では，すべての素イデアル $\mathfrak q$ に対して
$I_\mathfrak q$ が単項であるという性質を，「$I$ は局所単項イデアルである」
と表すことがある．しかし，本書で「局所的」は常に「Zariski 位相に関して
局所的」（またはその時点で扱う位相に関して局所的）を意味するので，この
用語法は採用できない．

\medskip\noindent
$R = \mathbf{Q}[x]$ とし，$M = \sum \frac{1}{x - n}R$ を Section
\ref{examples-section-nonfree} で構成した加群とする．環\footnote{環 $A$ は，
局所環が Noether である非 Noether 整域の例である．}
$$
A = \text{Sym}^*_R(M)
$$
と，イデアル $I = M A = \bigoplus_{d \geq 1} \text{Sym}^d_R(M)$ を考える．
$M$ は $R$-加群として有限生成でないので，$I$ は $A$ のイデアルとして
有限個の元で生成されえない．可逆加群は有限生成なので，$I$ は可逆でない．
一方，$\mathfrak p \subset R$ を素イデアルとする．構成により
$M_\mathfrak p \cong R_\mathfrak p$ である．したがって，次数付き
$R_\mathfrak p$-代数として
$$
A_\mathfrak p = \text{Sym}^*_{R_\mathfrak p}(M_\mathfrak p) \cong
\text{Sym}^*_{R_\mathfrak p}(R_\mathfrak p) =
R_\mathfrak p[T]
$$
である．ゆえに，$I_\mathfrak p \subset A_\mathfrak p$ は非零因子 $T$ で
生成される．したがって確かに，任意の素イデアル $\mathfrak q \subset A$ に
対し $I_\mathfrak q$ は一元で生成される．

\begin{lemma}
\label{examples-lemma-locally-principal-not-invertible}
整域 $A$ と非零イデアル $I \subset A$ で，すべての素イデアル
$\mathfrak q \subset A$ に対して $I_\mathfrak q \subset A_\mathfrak q$ は
単項イデアルであるが，$I$ は可逆 $A$-加群でないものが存在する．
\end{lemma}

\begin{proof}
上の議論を参照せよ．
\end{proof}
% END EXPANDED MODULE ch110_31.tex
% BEGIN EXPANDED MODULE ch110_32.tex
\section{射影的でない有限平坦加群}
\label{examples-section-finite-flat-not-projective}

\noindent
これは Algebra, Remark \ref{algebra-remark-warning} の複製である．有限
$R$-加群が $R$-平坦であっても，自動的に射影的であるとは限らない．反例は，
$R = \mathcal{C}^\infty(\mathbf{R})$ を $\mathbf{R}$ 上の無限回微分可能関数の
環とし，$M = R_{\mathfrak m} = R/I$ とする場合である．ここで
$\mathfrak m = \{f \in R \mid f(0) = 0\}$ および
$I = \{f \in R \mid \exists \epsilon, \epsilon > 0 :
f(x) = 0\ \forall x, |x| < \epsilon\}$ とする．

\medskip\noindent
射 $\Spec(R/I) \to \Spec(R)$ は，開でない平坦閉埋入の例でもある．

\begin{lemma}
\label{examples-lemma-finite-flat-non-projective}
奇妙な平坦加群．
\begin{enumerate}
\item 射影的でない有限平坦 $R$-加群 $M$ をもつ環 $R$ が存在する．
\item 平坦であるが開でない閉埋入が存在する．
\end{enumerate}
\end{lemma}

\begin{proof}
上の議論を参照せよ．
\end{proof}
% END EXPANDED MODULE ch110_32.tex
% BEGIN EXPANDED MODULE ch110_33.tex
\section{局所自由でない射影加群}
\label{examples-section-projective-not-locally-free}

\noindent
二つの例を与える．一つは階数が $0$ と $1$ の間にある例で，もう一つは
階数が $\aleph_0$ である例である．

\begin{lemma}
\label{examples-lemma-ideal-generated-by-idempotents-projective}
$R$ を環とし，$I \subset R$ を可算個の冪等元で生成されるイデアルとする．
このとき，$I$ は射影 $R$-加群である．
\end{lemma}

\begin{proof}
$I = (e_1, e_2, e_3, \ldots)$ とし，$e_n$ は $R$ の冪等元であるとする．
$e_{n + 1}$ を帰納的に $e_n + (1 - e_n)e_{n + 1}$ で置き換えることにより，
$(e_1) \subset (e_2) \subset (e_3) \subset \ldots$ と仮定してよい．したがって，
$I = \bigcup_{n \geq 1} (e_n) = \colim_n e_nR$ である．この場合
$$
\Hom_R(I, M) = \Hom_R(\colim_n e_nR, M)
= \lim_n \Hom_R(e_nR, M) = \lim_n e_nM
$$
である．遷移写像 $e_{n + 1}M \to e_nM$ は $e_n$ による乗法で与えられ，
全射であることに注意せよ．したがって Algebra, Lemma
\ref{algebra-lemma-ML-exact-sequence} により，関手 $\Hom_R(I, M)$ は完全で
あり，すなわち $I$ は射影 $R$-加群である．
\end{proof}

\noindent
$P \subset Q$ を $R$-加群の包含とし，$Q$ は有限 $R$-加群，$P$ は局所自由
であると仮定する．Algebra, Definition
\ref{algebra-definition-locally-free} を参照せよ．$Q$ は $R$-加群として
$N$ 個の元で生成できると仮定する．Algebra, Lemma
\ref{algebra-lemma-map-cannot-be-injective} から，$P$ は有限局所自由である
（自由部分の階数は高々 $N$）．この場合，$P$ は有限 $R$-加群である．
Algebra, Lemma \ref{algebra-lemma-finite-projective} を参照せよ．

\medskip\noindent
これと上の結果を合わせると，可算個の冪等元で生成される有限生成でない
イデアルは射影的であるが局所自由でないことが分かる．具体例として
$R = \prod_{n \in \mathbf{N}} \mathbf{F}_2$ と，冪等元
$$
e_n = (1, 1, \ldots, 1, 0, \ldots )
$$
で生成されるイデアル $I$ を取れる．ここで $1$ の列の長さは $n$ である．

\begin{lemma}
\label{examples-lemma-ideal-projective-not-locally-free}
環 $R$ とイデアル $I$ で，$I$ は $R$-加群として射影的であるが
$R$-加群として局所自由でないものが存在する．
\end{lemma}

\begin{proof}
上を参照せよ．
\end{proof}

\begin{remark}
\label{examples-remark-projective-with-dual-finite-not-finite}
上の例 $R = \prod_{n \in \mathbf{N}} \mathbf{F}_2$ と，冪等元
$e_n = (1, 1, \ldots, 1, 0, \ldots )$ で生成されるイデアル $I$ は，双対が
有限生成で，さらには有限自由である有限生成でない射影加群 $M = I$ の例も
与える．実際，Lemma \ref{examples-lemma-ideal-generated-by-idempotents-projective} の
証明の方針を用いて，読者は $\Hom_R(I, R) \cong R$ を示せる．
\end{remark}

\begin{lemma}
\label{examples-lemma-chow-group-product}
$K$ を体とする．$C_i$，$i = 1, \ldots, n$ を $K$ 上の滑らかで射影的な
幾何学的既約曲線とする．$P_i \in C_i(K)$ を有理点とし，
$Q_i \in C_i$ を $[\kappa(Q_i) : K] = 2$ を満たす点とする．このとき，
$U_i = C_i \setminus \{Q_i\}$ と置けば，
$[P_1 \times \ldots \times P_n]$ は
$\CH_0(U_1 \times_K \ldots \times_K U_n)$ において非零である．
\end{lemma}

\begin{proof}
次数写像
$\deg : \CH_0(C_1 \times_K \ldots \times_K C_n) \to \mathbf{Z}$ がある．
各 $Q_i$ は $K$ 上次数 $2$ をもつので，「境界」
$$
C_1 \times_K \ldots \times_K C_n
\setminus
U_1 \times_K \ldots \times_K U_n
$$
に台をもつ任意の零サイクルの次数は $2$ で割り切れる．
\end{proof}

\noindent
上の補題を用いて，射影的だが局所自由でない加群の別の例を次のように
構成できる．$C_n$，$n = 1, 2, 3, \ldots$ を，それぞれ点の対
$P_n, Q_n \in C_n$ をもつ $\mathbf{Q}$ 上の滑らかで射影的な幾何学的既約
曲線とする．ここで $\kappa(P_n) = \mathbf{Q}$ であり，$\kappa(Q_n)$ は
$\mathbf{Q}$ の二次拡大である．$U_n = C_n \setminus \{Q_n\}$ と置く．これは
アフィン曲線である．$\mathcal{L}_n$ を，$U_n$ 上の $P_n$ のイデアル層の
逆とする．零サイクル群 $\CH_0(U_n)$ において
$c_1(\mathcal{L}_n) = [P_n]$ であることに注意せよ．
$A_n = \Gamma(U_n, \mathcal{O}_{U_n})$ と置く．
$L_n = \Gamma(U_n, \mathcal{L}_n)$ とすれば，これは $A_n$ 上階数 $1$ の
局所自由加群である．
$$
B_n = A_1 \otimes_{\mathbf{Q}} A_2 \otimes_{\mathbf{Q}} \ldots
\otimes_{\mathbf{Q}} A_n
$$
と置く．すると $\Spec(B_n) = U_1 \times \ldots \times U_n$ であり，すべての
積は $\Spec(\mathbf{Q})$ 上で取る．$i \leq n$ に対して
$$
L_{n, i} =
A_1 \otimes_{\mathbf{Q}} \ldots \otimes_{\mathbf{Q}} L_i
\otimes_{\mathbf{Q}} \ldots \otimes_{\mathbf{Q}} A_n
$$
と置く．これは階数 $1$ の局所自由 $B_n$-加群である．これはまた
$\text{pr}_i^*\mathcal{L}_n$ の大域切断でもあることに注意せよ．
% P12-ERR-0339: see ../control/ERRATA_CANDIDATES.jsonl
$$
B_\infty = \colim_n B_n
\quad\text{and}\quad
L_{\infty, i} = \colim_n L_{n, i}
$$
と置く．最後に
$$
M = \bigoplus\nolimits_{i \geq 1} L_{\infty, i}.
$$
と置く．これは有限局所自由加群の直和なので射影的である．$M$ は局所自由で
ないと主張する．実際，$M_f$ が $(B_\infty)_f$ 上自由となる非零関数
$f \in B_\infty$ が存在すると仮定する．$e_1, e_2, \ldots$ を基底とする．
$f \in B_n$ となる $n \geq 1$ を選ぶ．さらに，
$e_1, \ldots, e_{n + 1}$ が
$$
\bigoplus\nolimits_{1 \leq i \leq m} L_{m, i}.
$$
に属するように $m \geq n + 1$ を選ぶ．元 $e_1, \ldots, e_{n + 1}$ は忠実
平坦な基底変換後に基底の一部なので，Chern 類
$$
c_i(\text{pr}_1^*\mathcal{L}_1 \oplus \ldots \oplus
\text{pr}_m^*\mathcal{L}_m), \quad i = m, m - 1, \ldots, m - n
$$
は
$$
D(f) \subset U_1 \times \ldots \times U_m
$$
の Chow 群において零である．$f$ は $U_1 \times \ldots \times U_n$ 上の
関数の引き戻しなので，特に
$$
c_{m - n}(\mathcal{O}_W^{\oplus n} \oplus
\text{pr}_1^*\mathcal{L}_{n + 1} \oplus \ldots
\oplus \text{pr}_{m - n}^*\mathcal{L}_m) = 0.
$$
が，体 $K = \mathbf{Q}(C_1 \times \ldots \times C_n)$ 上の多様体
$$
W = (C_{n + 1} \times \ldots \times C_m)_K
$$
上で成り立つ．言い換えると，サイクル
$$
[(P_{n + 1} \times \ldots \times P_m)_K]
$$
は $W$ 上の零サイクルの Chow 群において零である．これは上の Lemma
\ref{examples-lemma-chow-group-product} に矛盾する．実際，点 $Q_i$，
$n + 1 \leq i \leq m$ は $(C_n)_K$ 上の対応する点 $Q_i'$ を誘導し，
% P12-ERR-0340: see ../control/ERRATA_CANDIDATES.jsonl
$K/\mathbf{Q}$ は幾何学的既約なので $[\kappa(Q_i') : K] = 2$ である．
% P12-ERR-0341: see ../control/ERRATA_CANDIDATES.jsonl

\begin{lemma}
\label{examples-lemma-projective-not-locally-free}
可算環 $R$ と射影加群 $M$ で，$M$ は階数 $1$ の局所自由加群の可算直和で
あるが，局所自由でないものが存在する．
\end{lemma}

\begin{proof}
上を参照せよ．
\end{proof}
% END EXPANDED MODULE ch110_33.tex
% BEGIN EXPANDED MODULE ch110_34.tex
\section{非零平坦イデアルをもつ零次元局所環}
\label{examples-section-zero-dimensional-flat-ideal}

\noindent
\cite{Lazard} および \cite{Autour} には，非零平坦イデアルをもつ零次元局所環の
例がある．その構成を述べる．$k$ を体とし，$X_i, Y_i$，$i \geq 1$ を変数と
する．$R = k[X_i, Y_i]/(X_i - Y_i X_{i + 1}, Y_i^2)$ とする．この環における
$X_i$ の像，および $Y_i$ の像をそれぞれ $x_i$，$y_i$ と書く．この環では
$$
x_i = y_i x_{i + 1} = y_i y_{i +1} x_{i + 2} =
y_i y_{i + 1} y_{i + 2} x_{i + 3} = \ldots
$$
であることに注意せよ．環 $R$ はただ一つの素イデアル
$\mathfrak m = (x_i, y_i)$ をもつ．イデアル $I = (x_i)$ は $R$-加群として
平坦であると主張する．

\medskip\noindent
$R$ における $x_i$ の零化イデアルは
$(x_1, x_2, x_3, \ldots, y_i, y_{i + 1}, y_{i + 2}, \ldots)$ であることに
注意せよ．元 $e_i$，$i \geq 1$ と関係 $e_i = y_i e_{i + 1}$ で生成される
$R$-加群 $M$ を考える．$M$ は，遷移写像
$$
R \xrightarrow{y_1} R \xrightarrow{y_2} R \xrightarrow{y_3} \ldots
$$
をもつ $R$ のコピーの余極限 $\colim_i R$ なので平坦である．$M$ における
$e_i$ の零化イデアルは
$(x_1, x_2, x_3, \ldots, y_i, y_{i + 1}, y_{i + 2}, \ldots)$ であることに
注意せよ．$M$ の任意の元，および $I$ の任意の元は，ある $f, h \in R$ に
よりそれぞれ $f e_i$，$h x_i$ と書けるので，写像 $M \to I$，
$e_i \to x_i$ は同型であり，$I$ は平坦である．

\begin{lemma}
\label{examples-lemma-zero-dimensional-flat-ideal}
\begin{slogan}
平坦イデアルをもつ零次元環．
\end{slogan}
ただ一つの素イデアルをもつ局所環 $R$ と，平坦 $R$-加群である非零イデアル
$I \subset R$ が存在する．
\end{lemma}

\begin{proof}
上の議論を参照せよ．
\end{proof}
% END EXPANDED MODULE ch110_34.tex
% BEGIN EXPANDED MODULE ch110_35.tex
\section{全射でない零次元環のエピ射}
\label{examples-section-epimorphism-not-surjective}

\noindent
\cite{Lazard-deux} および \cite{Autour} に次の例がある．$k$ を体とし，環準同型
$$
k[x_1, x_2, \ldots, z_1, z_2, \ldots]/
(x_i^{4^i}, z_i^{4^i})
\longrightarrow
k[x_1, x_2, \ldots, y_1, y_2, \ldots]/
(x_i^{4^i}, y_i - x_{i + 1}y_{i + 1}^2)
$$
で，$x_i$ を $x_i$ へ，$z_i$ を $x_iy_i$ へ写すものを考える．右辺では
$y_i^{4^{i + 1}}$ は零であるが，$y_1$ は零でないことに注意せよ
（詳細は省略する）．この写像は全射でない．実際，
$\deg(x_i) = -1$，$\deg(z_i) = 0$，$\deg(y_i) = 1$ と置いて，上を
$\mathbf{Z}$-次数付き代数の写像とみなせば，$y_1$ が像に属さないことは
明らかである．最後に，この写像はエピ射である．なぜなら
$$
y_{i - 1} \otimes 1 = x_i y_i^2 \otimes 1 = y_i \otimes x_i y_i =
x_i y_i \otimes y_i = 1 \otimes x_i y_i^2.
$$
したがって，始域上で取った終域のテンソル積は終域と同型である．
% P12-ERR-0342: see ../control/ERRATA_CANDIDATES.jsonl

\begin{lemma}
\label{examples-lemma-epi-not-surjective}
全射でない次元 $0$ の局所環のエピ射が存在する．
\end{lemma}

\begin{proof}
上の議論を参照せよ．
\end{proof}
% END EXPANDED MODULE ch110_35.tex
% BEGIN EXPANDED MODULE ch110_36.tex
\section{有限型で有限表示でなく，素点で平坦な例}
\label{examples-section-ft-not-fp-flat-at-prime}

\noindent
$k$ を体とし，局所環 $A_0 = k[x, y]_{(x, y)}$ を考える．
$\mathfrak p_{0, n} = (y + x^n + x^{2n + 1})$ と置く．これは素イデアルである．
$$
A = A_0[z_1, z_2, z_3, \ldots]/(z_n z_m, z_n(y + x^n + x^{2n + 1}))
$$
と置く．$A \to A_0$ は，核が平方零イデアルである全射であることに注意せよ．
したがって $A$ も局所環であり，$A$ の素イデアルは $A_0$ の素イデアルと
一対一に対応する．$\mathfrak p_{0, n}$ に対応する $A$ の素イデアルを
$\mathfrak p_n$ と書く．$\mathfrak p_n$ は $A$ における $z_n$ の零化イデアル
であることに注意せよ．
$$
C = A[z]/(xz^2 + z + y)[\frac{1}{2zx + 1}]
$$
と置く．$A \to C$ はエタール環準同型であることに注意せよ．Algebra,
Example \ref{algebra-example-make-standard-smooth} を参照せよ．
$\mathfrak q \subset C$ を $x$，$y$，$z$ およびすべての $z_n$ で生成される
極大イデアルとする．$A \to C$ は平坦なので，$C$ における $z_n$ の
零化イデアルは $\mathfrak p_nC$ である．次を計算する．
\begin{align*}
C/\mathfrak p_n C
& =
A_0[z]/(xz^2 + z + y, y + x^n + x^{2n + 1})[1/(2zx + 1)] \\
& =
k[x]_{(x)}[z]/(xz^2 + z - x^n - x^{2n + 1})[1/(2zx + 1)] \\
& =
k[x]_{(x)}[z]/(z - x^n) \times
k[x]_{(x)}[z]/(xz + x^{n + 1} + 1)[1/(2zx + 1)] \\
& =
k[x]_{(x)} \times k(x)
\end{align*}
これは
$(z - x^n)(xz + x^{n + 1} + 1) = xz^2 + z - x^n - x^{2n + 1}$ による．
したがって，$\mathfrak r_n = \mathfrak p_nC + (z - x^n)C$ および
$\mathfrak q_n = \mathfrak p_nC + (xz + x^{n + 1} + 1)C$ と置けば，
$\mathfrak p_nC = \mathfrak r_n \cap \mathfrak q_n$ である．
$\mathfrak q_n + \mathfrak r_n = C$ なので，さらに
$\mathfrak p_nC = \mathfrak r_n \mathfrak q_n$ である．したがって，
$\mathfrak q_n$ は $\xi_n = (z - x^n)z_n$ の零化イデアルである．一方で
$\mathfrak r_n \subset \mathfrak q$ であり，他方で
$\mathfrak q_n + \mathfrak q = C$ であることに注意せよ．例えば，これは
$\mathfrak q_n$ が $\mathfrak q$ と異なる $C$ の極大イデアルであることから
従う．同様に，$n \not = m$ に対して
$\mathfrak q_n + \mathfrak q_m = C$ である．ここで
$$
B = \Im(C \longrightarrow C_{\mathfrak q})
$$
と置く．$xz + x^{n + 1} + 1$ は $\mathfrak q$ に属さないので，元 $\xi_n$ は
$B$ において零に写ることに注意せよ．$\mathfrak q$ の像を
$\mathfrak q' \subset B$ と書く．構成により $B$ は有限型 $A$-代数で，
$B_{\mathfrak q'} \cong C_{\mathfrak q}$ である．特に，
$B_{\mathfrak q'}$ は $A$ 上平坦である．

\medskip\noindent
$B_{g'}$ が $A$ 上有限表示となる元 $g' \in B$，
$g' \not \in \mathfrak q'$ は存在しないと主張する．この主張の証明を概説する．
$g' \in B$ に写る元 $g \in C$ を選ぶ．写像 $C_g \to B_{g'}$ を考える．
Algebra, Lemma \ref{algebra-lemma-finite-presentation-independent} により，
$B_g$ が $A$ 上有限表示であることと，$C_g \to B_{g'}$ の核が有限生成で
あることは同値である．
% P12-ERR-0343: see ../control/ERRATA_CANDIDATES.jsonl
しかし元 $g \in C$ は $\mathfrak q$ に属さないので，
$A_0[z]/(xz^2 + z + y)$ の非零元に写る．したがって，$g$ が属しうる
素イデアル $\mathfrak q_n$ は有限個だけである．実際，素イデアル
$(y + x^n + x^{2n + 1}, xz + x^{n + 1} + 1)$ は二次元既約 Noether 空間
$\Spec(k[x, y, z]/(xz^2 + z + y))$ の余次元 1 の点の無限族である．写像
$$
\bigoplus\nolimits_{g \not \in \mathfrak q_n} C/\mathfrak q_n
\longrightarrow
C_g, \quad
(c_n) \longrightarrow \sum c_n \xi_n
$$
は単射で，その像は $C_g \to B_{g'}$ の核である．この主張の証明は省略する．
（ヒント：$I$ を $A \to A_0$ の核として，$A_0$-加群として
$A = A_0 \oplus I$ と書く．同様に $C = C_0 \oplus IC$ と書く．
$IC = \bigoplus Cz_n \cong \bigoplus (C/\mathfrak r_n \oplus C/\mathfrak q_n)$
と書き，各直和因子への $g$ による乗法の作用を調べよ．）
以上で主張の証明の概説が終わる．またこれは，上のような任意の $g'$ に対して
$B_{g'}$ が $A$ 上平坦でないことも示す．実際，もし平坦なら，$B_{g'}$ に
おける $z_n$ の像の零化イデアルは $\mathfrak p_nB_{g'}$ であり，
$z - x^n$ を含まないはずである．

\medskip\noindent
帰結として，\cite[Part I, Remarques (3.4.7) (\romannumeral 5)]{GruRay} で
提起された問いに否定的に答えられる．正確な主張は次のとおりである．

\begin{lemma}
\label{examples-lemma-example-raynaud-gruson}
局所環 $A$，有限型環準同型 $A \to B$，および $\mathfrak m_A$ の上にある
素イデアル $\mathfrak q$ で，$B_{\mathfrak q}$ は $A$ 上平坦であるが，
任意の元 $g \in B$，$g \not \in \mathfrak q$ に対して環 $B_g$ は
$A$ 上有限表示でなく，また $A$ 上平坦でもないものが存在する．
\end{lemma}

\begin{proof}
上の議論を参照せよ．
\end{proof}
% END EXPANDED MODULE ch110_36.tex
% BEGIN EXPANDED MODULE ch110_37.tex
\section{有限型かつ平坦で有限表示でない例}
\label{examples-section-finite-type-flat-not-finite-presentation}

\noindent
本節では，有限型かつ平坦であるが有限表示でない環準同型と射の例を与える．

\medskip\noindent
$R$ を，$R/I$ が有限平坦だが射影的でない加群となるイデアル $I$ をもつ環と
する．具体例については Section \ref{examples-section-finite-flat-not-projective} を
参照せよ．これは $I$ が有限生成でないことを意味することに注意せよ．
Algebra, Lemma \ref{algebra-lemma-finitely-generated-pure-ideal} を参照せよ．
また，$I = I^2$ であることに注意せよ．Algebra, Lemma
\ref{algebra-lemma-pure} を参照せよ．例の基礎環は $R$ であり，対応する
基礎スキームは $S = \Spec(R)$ とする．

\medskip\noindent
慣例により $\epsilon^2 = 0$ として，環準同型 $R \to R \oplus R/I\epsilon$ を
考える．これは有限かつ平坦であるが有限表示でない環準同型である．すべての
ファイバー環は完全交叉かつ幾何学的既約である．

\medskip\noindent
$A = R[x, y]/(xy, ay; a \in I)$ とする．$R$-加群として
$A = \bigoplus_{i \geq 0} Ry^i \oplus \bigoplus_{j > 0} R/Ix^j$ であることに
注意せよ．
% P12-ERR-0344: see ../control/ERRATA_CANDIDATES.jsonl
したがって，$R \to A$ は有限表示でない平坦有限型環準同型である．各
ファイバー環は $\kappa(\mathfrak p)[x, y]/(xy)$ または
$\kappa(\mathfrak p)[x]$ のいずれかと同型である．

\medskip\noindent
前の例から，$B = R[X_0, X_1, X_2]/(X_1X_2, aX_2; a \in I)$ と置くことに
より射影射の例を作れる．この場合，$X = \text{Proj}(B) \to S$ は有限表示で
ない固有平坦射で，各 $s \in S$ に対してファイバー $X_s$ は
$\mathbf{P}^1_s$ または $X_1X_2$ の消滅で定義される
$\mathbf{P}^2_s$ の閉部分スキームのいずれかと同型である（後者は算術種数
$0$ の射影節点曲線である）．

\medskip\noindent
$M = R \oplus R \oplus R/I$ とする．$B = \text{Sym}_R(M)$ を $M$ の対称代数と
置き，$X = \text{Proj}(B)$ とする．このとき $X \to S$ は有限表示でない
固有平坦射で，$s \in S$ に対して幾何学的ファイバーは
$\mathbf{P}^1_s$ または $\mathbf{P}^2_s$ のいずれかと同型である．特に，
これらのファイバーは滑らかで幾何学的既約である．

\begin{lemma}
\label{examples-lemma-finite-type-flat-not-finitely-presented}
次の例が存在する．
\begin{enumerate}
\item 幾何学的既約な完全交叉ファイバー環をもつ，有限表示でない平坦有限型
環準同型．
\item 幾何学的連結，幾何学的被約，次元 1 の完全交叉ファイバー環をもつ，
有限表示でない平坦有限型環準同型．
\item 各ファイバーが $\mathbf{P}^1_s$ または $\mathbf{P}^2_s$ 内の
$X_1X_2$ の消滅跡のいずれかと同型である，有限表示でないスキームの
固有平坦射 $X \to S$．
\item 各ファイバーが $\mathbf{P}^1_s$ または $\mathbf{P}^2_s$ のいずれかと
同型である，有限表示でないスキームの固有平坦射 $X \to S$．
\end{enumerate}
\end{lemma}

\begin{proof}
上の議論を参照せよ．
\end{proof}
% END EXPANDED MODULE ch110_37.tex
% BEGIN EXPANDED MODULE ch110_38.tex
\section{有限型環準同型の位相}
\label{examples-section-topology-finite-type}

\noindent
$A \to B$ を局所整域の局所準同型とする．$A$ が Noether，$A \to B$ が
本質的有限型，かつ $A/\mathfrak m_A \subset B/\mathfrak m_B$ が有限ならば，
$\mathfrak q \not = \mathfrak m_B$ を満たす素イデアル $\mathfrak q \subset B$
で，$A \to B/\mathfrak q$ が準有限環準同型の局所化となるものが存在する．
More on Morphisms, Lemma
\ref{more-morphisms-lemma-quasi-finite-quasi-section-meeting-nearby-open} を
参照せよ．

\medskip\noindent
本節では，$A$ がもはや Noether でないとこの結果が偽になることを示す例を
与える．
% P12-ERR-0345: see ../control/ERRATA_CANDIDATES.jsonl
すなわち，$k$ を体として
$$
A = \{a_0 + a_1 x + a_2 x^2 + \ldots \mid a_0 \in k, a_i \in k((y))
\text{，ただし }i\geq 1\}
$$
および
$$
C = \{a_0 + a_1 x + a_2 x^2 + \ldots \mid a_0 \in k[y], a_i \in k((y))
\text{，ただし }i\geq 1\}.
$$
と置く．$C$ は $A$ 上 $y$ で生成されるので，包含 $A \to C$ は有限型である．
$A$ は極大イデアル
$\mathfrak m = \{a_1 x + a_2 x^2 + \ldots \in A\}$ をもつ局所環で，
$(0)$ と $\mathfrak m$ 以外に素イデアルをもたないと主張する．実際，
$A$ の元 $f = a_0 + a_1 x + a_2 x^2 + \ldots$ は $a_0 \not = 0$ ならば
可逆である．$\mathfrak q \subset A$ を非零素イデアルとし，
$f = a_i x^i + \ldots \in \mathfrak q$ とする．冪級数の性質を用いると，
任意の $g \in k((y))$ に対し $g^{i + 1} x^{i + 1} \in \mathfrak q$，
すなわち $gx \in \mathfrak q$ であることが分かる．これにより
$\mathfrak q = \mathfrak m$ が示された．

\medskip\noindent
環 $C$ のスペクトルについても，上と同様に議論すると，任意の非零
素イデアルは $\mathfrak m$ の上にある素イデアル
$\mathfrak p = \{a_1 x + a_2 x^2 + \ldots \in C\}$ を含むことが分かる．
したがって，$(0)$ の上にある $C$ の素イデアルは $(0)$ だけである．
$\mathfrak m_C = yC + \mathfrak p$ および $B = C_{\mathfrak m_C}$ と置く．
このとき $A \to B$ が所望の例である．

\begin{lemma}
\label{examples-lemma-topology-finite-type}
局所整域の局所準同型 $A \to B$ で，本質的有限型かつ
$A/\mathfrak m_A \to B/\mathfrak m_B$ が有限であるが，$B$ の任意の
素イデアル $\mathfrak q \not = \mathfrak m_B$ に対して環準同型
$A \to B/\mathfrak q$ が準有限環準同型の局所化でないものが存在する．
\end{lemma}

\begin{proof}
上の議論を参照せよ．
\end{proof}
% END EXPANDED MODULE ch110_38.tex
% BEGIN EXPANDED MODULE ch110_39.tex
\section{純粋だが普遍的には純粋でない例}
\label{examples-section-pure-not-universally}

\noindent
$k$ を体とする．
$$
R = k[[x, xy, xy^2, \ldots]] \subset k[[x, y]].
$$
言い換えると，形式的冪級数 $f \in k[[x, y]]$ が $R$ に属するための必要十分条件は，
$f(0, y)$ が定数であることである．特に $R[1/x] = k[[x, y]][1/x]$ であり，
$R/xR$ は，その平方が零となる極大イデアルをもつ局所環である．
$R[y] \subset k[[x, y]]$ を，$f(0, y)$ が $y$ の多項式となる形式的冪級数
$f \in k[[x, y]]$ 全体とする．このとき $R \to R[y]$ は有限型だが有限表示でない
環準同型で，$x$ を可逆化すると同型を誘導する．また，全射
$R[y]/xR[y] \to k[y]$ があり，その核の平方は零である．有限表示な環準同型
$R \to S = R[t]/(xt - xy)$ を考える．ここでも
$R[1/x] \to S[1/x]$ は同型であり，この場合
$S/xS \cong (R/xR)[t]/(xy)$ は冪零な核をもつ全射として $k[t]$ に写る．
$x$ の可逆化後に同型を誘導し，$x$ を法として冪零核をもつ全射を誘導する全射
$S \to R[y]$，$t \longmapsto y$ がある．したがって $S \to R[y]$ の核は局所冪零である．
特に $S \to R[y]$ はスペクトル上に普遍的同相写像を誘導する．

\medskip\noindent
第一に，$S$ は $R$ 上相対的に純粋な $S$-加群であると主張する．$x$ を可逆化すると
同型を得るので，極大イデアル $\mathfrak m \subset R$ においてのみ確認すればよい．
$R$ はその極大イデアルに関して完備であるから Hensel 環であり，したがって，各素イデアル
$\mathfrak p \subset R$，$\mathfrak p \not = \mathfrak m$ に対し，$\mathfrak p$ の上にある
$S$ の一意な素イデアル $\mathfrak q$ が
$\mathfrak mS + \mathfrak q \not = S$ を満たすことだけを確認すればよい．
$\mathfrak p \not = \mathfrak m$ なので，これは $k[[x, y]][1/x]$ の一意な素イデアルに対応する．
したがって，$\mathfrak p = (0)$ であるか，または $x$ と同伴でないある既約元
$f \in k[[x, y]]$ に対して $\mathfrak p = (f)$ である
（ここでは $k[[x, y]]$ が UFD であることを用いる -- 将来の参照をここに挿入）．
% P12-ERR-0346: frozen unresolved editorial placeholder; see ERRATA_CANDIDATES.jsonl.
前者では $\mathfrak q = (0)$ であり，結果は明らかである．後者では，$f$ に単元を掛けることで
$f \in R[y]$ としてよい（Weierstrass の予備定理；詳細は省略する）．すると
$R[y]/fR[y] \cong k[[x, y]]/(f)$ であることが容易に分かるので，$f$ は $R[y]$ の
素イデアルを定め，$\mathfrak mR[y] + fR[y] \not = R[y]$ である．
$S \to R[y]$ はスペクトル上に普遍的同相写像を誘導するため，$S$ に対する所望の結果も従う．

\medskip\noindent
第二に，$S$ は $R$ 上普遍的に相対純粋ではないと主張する．実際，これを見るには，
付値環 $\mathcal{O}$ と局所環準同型 $R \to \mathcal{O}$ であって，
$\Spec(R[y] \otimes_R \mathcal{O}) \to \Spec(\mathcal{O})$ が
$\Spec(\mathcal{O})$ の閉点に到達しないものを見つければ十分である．同値な言い方をすれば，
$\mathcal{O}$ の付値を $v$ とするとき，$\varphi(x) \not = 0$ かつ
$v(\varphi(x)) > v(\varphi(xy))$ を満たす $\varphi : R \to \mathcal{O}$ を
見つけなければならない（これは $y$ の付値が負であることを意味する）．
そのため，明らかな仕方で定義される環準同型
$$
R
\longrightarrow
\{a_0 + a_1 x + a_2 x^2 + \ldots \mid a_0 \in k[y^{-1}], a_i \in k((y))\}
$$
を考える．右辺を極大イデアル $(y^{-1}, x)$ で局所化した環を支配する付値環
$\mathcal{O}$ を取れば，目的が達せられる．

\begin{lemma}
\label{examples-lemma-pure-not-universally-pure}
有限表示なアフィンスキームの射 $X \to S$ と有限表示な $\mathcal{O}_X$-加群
$\mathcal{F}$ で，$\mathcal{F}$ は $S$ に相対的に純粋だが，$S$ に相対的に普遍純粋ではない
ものが存在する．
\end{lemma}

\begin{proof}
上の議論を参照せよ．
\end{proof}
% END EXPANDED MODULE ch110_39.tex
% BEGIN EXPANDED MODULE ch110_40.tex
\section{平坦でない形式的に滑らかな環準同型}
\label{examples-section-formally-smooth-nonflat}

\noindent
$k$ を体とする．$k$-代数 $k[\mathbf{Q}]$ を考える．これは基底
$x_\alpha, \alpha \in \mathbf{Q}$ をもち，積が
$x_\alpha x_\beta = x_{\alpha + \beta}$ によって定まる $k$-代数である
（特に $x_0 = 1$ である）．$k$-代数準同型
$$
k[\mathbf{Q}] \longrightarrow k, \quad
x_\alpha \longmapsto 1.
$$
を考える．これは全射であり，その核 $J$ は元 $x_\alpha - 1$ で生成される．
$J/J^2$ を計算しよう．$J/J^2$ 上の $x_\alpha$ による乗法は恒等写像であることに注意する．
$J^2$ を法とする $x_\alpha - 1$ の類を $z_\alpha$ と書く．これらの類は $J/J^2$ を生成する．
$$
(x_\alpha - 1)(x_\beta - 1) = x_{\alpha + \beta} - x_\alpha - x_\beta + 1 =
(x_{\alpha + \beta} - 1) - (x_\alpha - 1) - (x_\beta - 1)
$$
なので，$J/J^2$ において $z_{\alpha + \beta} = z_\alpha + z_\beta$ である．
$J/J^2$ の一般の元は，$\lambda_\alpha \in k$（非零なものは有限個）を用いて
$\sum \lambda_\alpha z_\alpha$ の形に書ける．$k$ の標数が $p > 0$ ならば
$$
0 = pz_{\alpha/p} = z_{\alpha/p} + \ldots + z_{\alpha/p} = z_\alpha
$$
であり，$J/J^2 = 0$ と分かる．$k$ の標数が零ならば
$$
J/J^2 = \mathbf{Q} \otimes_{\mathbf{Z}} k \cong k
$$
（詳細は省略する）であり，零ではない．

\medskip\noindent
$k$ の標数が正ならば，$k[\mathbf{Q}] \to k$ は形式的に滑らかな環準同型であると主張する．
実際，実線部分が可換な図式
$$
\xymatrix{
k \ar[r] \ar@{..>}[rd] & A \\
k[\mathbf{Q}] \ar[u] \ar[r]^\varphi & A' \ar[u]
}
$$
が与えられ，$A' \to A$ は平方零な核 $I$ をもつ全射であるとする．
$k[\mathbf{Q}] \to k$ が形式的に滑らかであることを示すには，$\varphi$ が $k$ を経由することを
証明すればよい．$\varphi(x_\alpha - 1)$ は $A$ で零に写るので，$\varphi$ は写像
$\overline{\varphi} : J/J^2 \to I$ を誘導し，これが零であることが所望の分解の障害である．
標数が $p > 0$ ならば $J/J^2 = 0$ なので，所望の結果，すなわちこの場合
$k[\mathbf{Q}] \to k$ が形式的に滑らかであることを得る．最後に，この環準同型は平坦でない．
例えば，非零因子 $x_2 - 1$ が零に写るからである．

\begin{lemma}
\label{examples-lemma-formally-smooth-nonflat}
平坦でない形式的に滑らかな環準同型が存在する．
\end{lemma}

\begin{proof}
上の議論を参照せよ．
\end{proof}
% END EXPANDED MODULE ch110_40.tex
% BEGIN EXPANDED MODULE ch110_41.tex
\section{平坦でない形式的エタール環準同型}
\label{examples-section-formally-etale-nonflat}

\noindent
この節では，\cite[0, Example 19.10.3(i)]{EGA} の最後の文に対する反例を与える
（これは後の正誤表で捉えられた項目の一つではない）．局所環 $A$ と，
$J^2 = J$ を満たす非零の真のイデアル $J$（したがって $J$ は有限生成でない）に対し，
$A \to A/J$ を考える．$J$ を極大イデアルとする代数閉非 Archimedes 体の付値環は，
そのような $(A, J)$ の一つの供給源である．これらの平坦でない商写像は形式的エタールである．
実際，可換図式
$$
\xymatrix{
A/J \ar[r] & R/I \\
A \ar[u] \ar[r]^\varphi & R \ar[u]
}
$$
が与えられ，$I$ は $I^2 = 0$ を満たす環 $R$ のイデアルであるとする．すると
% P12-ERR-0347: frozen phi(J)A codomain/ideal expression; see ERRATA_CANDIDATES.jsonl.
$$
\varphi(J) = \varphi(J^2) \subset (\varphi(J)A)^2 \subset I^2 = 0.
$$
なので，$A \to R$ は $A/J$ を一意に経由する．したがって，これは
\cite[IV, Theorem 18.4.6(i)]{EGA} にある，局所有限型かつ形式的エタールな射に対する
「構造定理」の形式的エタールの場合にも反例を与える
（ただし，実際に通常用いられる部分 (ii) に対する反例ではない）．後者の証明の誤りは，
証明の最後の段階で誤った \cite[0, Example 19.3.10(i)]{EGA} を援用していることであり，
ここから今述べた反例が入り込む．
% P12-ERR-0348: frozen internally inconsistent EGA example number; see ERRATA_CANDIDATES.jsonl.

\begin{lemma}
\label{examples-lemma-formally-etale-not-presented}
形式的エタールだが平坦でない環準同型が存在する．
\end{lemma}

\begin{proof}
上の議論を参照せよ．
\end{proof}
% END EXPANDED MODULE ch110_41.tex
% BEGIN EXPANDED MODULE ch110_42.tex
\section{非自明な余接複体をもつ形式的エタール環準同型}
\label{examples-section-formally-etale-nontrivial-cotangent-complex}

\noindent
$k$ を体とする．環
$$
R = k[\{x_n\}_{n \geq 1}, \{y_n\}_{n \geq 1}]/(
x_1y_1, x_{nm}^m - x_n, y_{nm}^m - y_n)
$$
を考える．すべての $x_n, y_n$ で生成される極大イデアルにおける局所化を $A$ とし，
その極大イデアルを $J \subset A$ と書く．$B = A/J$ とおく．構成により
$J^2 = J$ であり，したがって $A \to B$ は形式的エタールである
（Section \ref{examples-section-formally-etale-nonflat} を参照）．元 $x_1 \otimes y_1$ は
$$
J \otimes_A J \longrightarrow J.
$$
の核の非零元であると主張する．実際，$(A, J)$ は，環
$$
R_n = k[x_n, y_n]/(x_n^n y_n^n)
$$
を対応する極大イデアルで局所化して得られる $(A_n, J_n)$ の余極限である．このとき
$x_1 \otimes y_1$ は $x_n^n \otimes y_n^n \in J_n \otimes_{A_n} J_n$ に対応し，
これは非零である（省略する直接計算による）．$\otimes$ は余極限と可換なので，主張が従う．
\cite[III Section 3.3]{cotangent} により，$J$ は弱正則でない．したがって
\cite[III Proposition 3.3.3]{cotangent} により，余接複体 $L_{B/A}$ は零でない．
実際，より精密に述べられる．$J/J^2 = 0$ なので，
$H_0(L_{B/A}) = \Omega_{B/A}$ かつ $H_1(L_{B/A}) = 0$ である．しかし Quillen の
基本スペクトル系列の五項完全系列
（Cotangent, Remark \ref{cotangent-remark-elucidate-ss} または
\cite[Corollary 8.2.6]{Reinhard} を参照）と
$\text{Tor}_2^A(B, B) = \Ker(J \otimes_A J \to J)$ の非消滅から，
$H_2(L_{B/A})$ は非零であると結論する．

\begin{lemma}
\label{examples-lemma-formally-etale-nontrivial-cotangent-complex}
$L_{B/A}$ が零でない形式的エタールな全射環準同型 $A \to B$ が存在する．
\end{lemma}

\begin{proof}
上の議論を参照せよ．
\end{proof}
% END EXPANDED MODULE ch110_42.tex
% BEGIN EXPANDED MODULE ch110_43.tex
\section{平坦かつ形式的不分岐でも形式的エタールとは限らない}
\label{examples-section-flat-formally-unramified-not-formally-etale}

\noindent
More on Morphisms, Lemma
\ref{more-morphisms-lemma-unramified-flat-formally-etale} では，スキームの不分岐平坦射
$X \to S$ が形式的エタールであることを示した．この節の目的は，ここで「不分岐」を
「形式的不分岐」に置き換えられないことを示す二つの例を与えることである．第一の例は
完全環の特殊な性質を利用し，第二の例は Noether 正則環の準同型についてさえ結果が
成り立たないことを示す．

\begin{lemma}
\label{examples-lemma-perfect-closure-polynomial-ring}
$A = \mathbb{F}_p[T]$ を $\mathbb{F}_p$ 上一変数の多項式環とし，$A_{perf}$ を
$A$ の完全閉包とする．このとき $A \rightarrow A_{perf}$ は平坦かつ形式的不分岐だが，
形式的エタールではない．
\end{lemma}

\begin{proof}
Frobenius 写像 $F_A : A \to A$ のもとで，終域側の $A$ は，始域上，基底
$\{1, T, \dots, T^{p - 1}\}$ をもつ自由加群であることに注意する．したがって
$F_A$ は忠実平坦であり，ゆえに忠実平坦写像の余極限である $A \to A_{perf}$ も
忠実平坦である．$A_{perf}$ は完全環なので，相対 Frobenius
$F_{A_{perf}/A}$ は全射である．言い換えると $A_{perf} = A[A_{perf}^p]$ であり，
これから直ちに $\Omega_{A_{perf}/A} = 0$ が従う．したがって More on Morphisms, Lemma
\ref{more-morphisms-lemma-formally-unramified-differentials} により
$A \rightarrow A_{perf}$ は形式的不分岐である．

\medskip\noindent
$A \rightarrow A_{perf}$ が形式的に滑らかでないことを示せば十分である．$A$ は滑らかな
$\mathbb{F}_p$-代数なので，余接複体
% P12-ERR-0349: frozen uppercase P in the prime-field subscript; see ERRATA_CANDIDATES.jsonl.
$L_{A/\mathbb{F}_P} \simeq \Omega_{A/\mathbb{F}_p}[0]$ は次数 $0$ に集中することに注意する．
COTANGENT, Lemma \ref{cotangent-lemma-when-projective} を参照せよ．さらに，
Cotangent, Lemma \ref{cotangent-lemma-perfect-zero} により，$D(A_{perf})$ において
$L_{A_{perf}/\mathbb{F}_p} = 0$ である．$D(A_{perf})$ における余接複体の識別三角形
$$
L_{A/\mathbb{F}_p} \otimes_A A_{perf} \to
L_{A_{perf}/\mathbb{F}_p} \to
L_{A_{perf}/A} \to
(L_{A/\mathbb{F}_p} \otimes_A A_{perf})[1]
$$
を考える．Cotangent, Section \ref{cotangent-section-triangle} を参照せよ．すると
$L_{A_{perf}/A} = \Omega_{A/\mathbb{F}_p} \otimes_A A_{perf}[1]$ を得る．すなわち
$L_{A_{perf}/A}$ は次数 $-1$ に置かれた階数 $1$ の自由 $A_{perf}$-加群に等しい．
したがって More on Morphisms, Lemma
\ref{more-morphisms-lemma-NL-formally-smooth} および Cotangent, Lemma
\ref{cotangent-lemma-relation-with-naive-cotangent-complex} により，
$A \rightarrow A_{perf}$ は形式的に滑らかでない．
\end{proof}

\noindent
次の例も素標数の環を用いるが，おそらく少し意外である．欠点は，前の例よりも標数
$p$ の現象について多くの知識を要することである．素標数の環 $A$ の Frobenius 写像が
有限であるとき，$A$ は $F$-有限であるということを思い出そう．

\begin{lemma}
\label{examples-lemma-completion-etale}
$(A, \mathfrak m, \kappa)$ を，$[\kappa : \kappa^p] < \infty$ を満たす素標数
$p > 0$ の Noether 局所環とする．このとき，$A$ からその完備化への標準写像
$A \to A^\wedge$ は平坦かつ形式的不分岐である．しかし $A$ が正則だが優秀でないなら，
この写像は形式的エタールではない．
\end{lemma}

\begin{proof}
完備化の平坦性は Algebra, Lemma \ref{algebra-lemma-completion-flat} である．写像が形式的不分岐
であることを示すには，Algebra, Lemma
\ref{algebra-lemma-characterize-formally-unramified} により，
$\Omega_{A^\wedge/A} = 0$ を示せば十分である．

\medskip\noindent
証明を概説する．$\kappa$ の $p$-基底，すなわち $\kappa$ が $\kappa^p$ 上
$\overline{x}_i$ により極小生成されるような $\overline{x}_1, \ldots, \overline{x}_r$ に写る
$x_1, \ldots, x_r \in A$ を選ぶ．$\mathfrak m$ の極小生成系
$y_1, \ldots, y_s$ を選ぶ．各 $n$ に対し，元
$x_1, \ldots, x_r, y_1, \ldots, y_s$ は Frobenius によって
$(A/\mathfrak m^n)^p$ 上 $A/\mathfrak m^n$ を生成する．詳細の一部は省略する．
$F : A^\wedge \to A^\wedge$ は有限であると結論する．したがって
$\Omega_{A^\wedge/A}$ は有限 $A^\wedge$-加群である．一方，任意の
$a \in A^\wedge$ と $n$ に対して，$a - a_0 \in \mathfrak m^nA^\wedge$ となる
$a_0 \in A$ を取れる．ゆえに $\text{d}(a) \in \bigcap \mathfrak m^n \Omega_{A^\wedge/A}$ であり，
Algebra, Lemma \ref{algebra-lemma-intersect-powers-ideal-module-zero} によって
$\text{d}(a)$ は零である．したがって $\Omega_{A^\wedge/A} = 0$ である．

\medskip\noindent
$A$ が正則であると仮定する．すると Cohen の構造定理を用いて，
$x_1, \ldots, x_r, y_1, \ldots, y_s$ は環 $A^\wedge$ の $p$-基底である．すなわち
$$
A^\wedge = \bigoplus\nolimits_{I, J} (A^\wedge)^p
x_1^{i_1} \ldots x_r^{i_r} y_1^{j_1} \ldots y_s^{j_s}
$$
である．ここで $I = (i_1, \ldots, i_r)$，$J = (j_1, \ldots, j_s)$ であり，
$0 \leq i_a, j_b \leq p - 1$ である．詳細は省略する．特に
$\Omega_{A^\wedge}$ は，基底
$\text{d}(x_1), \ldots, \text{d}(x_r), \text{d}(y_1), \ldots, \text{d}(y_s)$ をもつ
自由 $A^\wedge$-加群であることが分かる．

\medskip\noindent
いま $A \to A^\wedge$ が形式的エタール，あるいは単に形式的に滑らかであるとすると，
Algebra, Proposition \ref{algebra-proposition-characterize-formally-smooth} により
$\NL_{A^\wedge/A}$ の次数 $-1, 0$ のコホモロジーは消滅する．環準同型
$\mathbf{F}_p \to A \to A^\wedge$ に対する Jacobi--Zariski 系列
（Algebra, Lemma \ref{algebra-lemma-exact-sequence-NL}）から，同型
$\Omega_A \otimes_A A^\wedge \cong \Omega_{A^\wedge}$ を得る．したがって
$\Omega_A$ は $\text{d}(x_1), \ldots, \text{d}(x_r), \text{d}(y_1), \ldots, \text{d}(y_s)$
を基底とする自由加群である．分数体を見て $A$ が正規であることを用いると，
$a \in A$ が $p$ 乗であるための必要十分条件は，$A^\wedge$ におけるその像が $p$ 乗である
ことだと分かる（詳細は省略する；Algebra, Lemma
\ref{algebra-lemma-derivative-zero-pth-power} を用いよ）．第二の帰結として，作用素
$\partial/\partial x_a$ および $\partial/\partial y_b$ が $A$ 上で定義される．

\medskip\noindent
% P12-ERR-0350: frozen grammatically incomplete transition; see ERRATA_CANDIDATES.jsonl.
以上から，$A$ は $p$-基底 $x_1, \ldots, x_r, y_1, \ldots, y_s$ をもち，$A^p$ 上有限であると
結論できることを示す．これは Kunz の結果 \cite[Corollary 2.6]{Kun76} により，$A$ が
優秀でないことに矛盾する．実際，$a \in A$ とし，
$$
a =  \sum\nolimits_{I, J} (a_{I, J})^p
x_1^{i_1} \ldots x_r^{i_r} y_1^{j_1} \ldots y_s^{j_s}
$$
と書く．ここで $a_{I, J} \in A^\wedge$ である．証明を終えるには
$a_{I, J} \in A$ を示せば十分である．両辺に作用素
$$
(\partial/\partial x_1)^{p - 1} \ldots
(\partial/\partial x_r)^{p - 1}
(\partial/\partial y_1)^{p - 1} \ldots
(\partial/\partial y_s)^{p - 1}
$$
を施すと，$I = (p - 1, \ldots, p - 1)$ かつ $J = (p - 1, \ldots, p - 1)$ に対して
$a_{I, J}^p \in A$ であると分かる．上の注意により，これは $a_{I, J} \in A$ も意味する．
$a$ を $a' = a - a_{I, J}^p x^I y^J$ で置き換えた後，$1$ 階低い微分作用素を用いれば，
別の係数 $a_{I, J}$ が $A$ に属することを得る．以下同様である．
% P12-ERR-0351: frozen malformed English differential-operator phrase; see ERRATA_CANDIDATES.jsonl.
\end{proof}

\begin{remark}
\label{examples-remark-reference-existence-regular-nonexcellent-rings}
剰余体が有限 $p$-基底をもつ優秀でない正則環は，標数 $p$ の体
$k = \overline{k}$ 上の $\mathbb{P}^2_k$ の関数体の中でさえ構成できる．
\cite[$\mathsection 4.1$]{DS18} を参照せよ．
\end{remark}

\noindent
Lemma \ref{examples-lemma-completion-etale} の証明は，実際にはもう少し多くを示している．

\begin{lemma}
\label{examples-lemma-excellent-regular-local-rings}
$(A, \mathfrak m, \kappa)$ を標数 $p > 0$ の正則局所環とする．
$[\kappa : \kappa^p] < \infty$ と仮定する．このとき，$A$ が優秀であるための必要十分条件は，
$A \to A^\wedge$ が形式的エタールであることである．
\end{lemma}

\begin{proof}
逆向きの含意は Lemma \ref{examples-lemma-completion-etale} から従う．順向きについては，
Lemma \ref{examples-lemma-completion-etale} により，$A \to A^\wedge$ が形式的不分岐であること，
同値な言い方をすれば $\Omega_{A^\wedge/A}$ が零であることは既に分かっている．
したがって，$A$ が優秀なとき完備化写像が形式的に滑らかであることを示せば十分である．
N\'eron--Popescu 非特異化により，$A \to A^\wedge$ は滑らかな $A$-代数の
フィルター余極限として書ける
（Smoothing Ring Maps, Theorem \ref{smoothing-theorem-popescu}）．したがって
$\NL_{A^\wedge/A}$ の次数 $-1$ のコホモロジーは消滅する．ゆえに Algebra, Proposition
\ref{algebra-proposition-characterize-formally-smooth} により，$A \to A^\wedge$ は形式的に滑らかである．
\end{proof}
% END EXPANDED MODULE ch110_43.tex
% BEGIN EXPANDED MODULE ch110_44.tex
\section{冪等元の集合で生成されるイデアルと局所化}
\label{examples-section-ideal-locally-idempotents}

\noindent
$R$ を環とする．環
$$
B(R) = R[x_n; n \in \mathbf{Z}]/(x_n(x_n - 1), x_nx_m; n \not = m)
$$
を考える．各素イデアル $\mathfrak q \subset B(R)$ は，
$$
\mathfrak q = \mathfrak pB(R) + (x_n; n \in \mathbf{Z})
$$
という形か，または
$$
\mathfrak q =
\mathfrak pB(R) + (x_n - 1) + (x_m; n \not = m, m \in \mathbf{Z}).
$$
という形であることは容易に示せる．したがって
$$
\Spec(B(R)) =
\Spec(R) \amalg \coprod\nolimits_{n \in \mathbf{Z}} \Spec(R)
$$
であるが，その位相は単なる非交和位相ではない．これは次の性質をもつ．
$n \in \mathbf{Z}$ で添字付けられた各コピーは開部分スキームであり，具体的には標準開
$D(x_n)$ である．「中央」の $\Spec(R)$ のコピーは，他の $\Spec(R)$ のコピーのうち
任意の無限個の合併の閉包に含まれる．この最後の $\Spec(R)$ のコピーは，冪等元 $x_n$ で
生成されるイデアル $(x_n, n \in \mathbf{Z})$ により切り出されることに注意する．したがって，
$\Spec(R)$ が連結ならば，上の分解はちょうど $\Spec(B(R))$ の連結成分への分解である．

\medskip\noindent
次に $A = \mathbf{C}[x, y]/((y - x^2 + 1)(y + x^2 - 1))$ とする．$A$ のスペクトルは
二つの既約成分 $C_1 = \Spec(A_1)$，$C_2 = \Spec(A_2)$ からなり，ここで
$A_1 = \mathbf{C}[x, y]/(y - x^2 + 1)$ および
$A_2 = \mathbf{C}[x, y]/(y + x^2 - 1)$ である．これらは
$(x, y) = (t, t^2 - 1)$ および $(x, y) = (t, -t^2 + 1)$ によりパラメータ付けられ，
$P = (-1, 0)$ と $Q = (1, 0)$ で交わることに注意する．次のように $B(A_1)$ と $B(A_2)$ を
貼り合わせて，$B(A)$ のねじれた版を作れる．$P$ の上では
$x_n \in B(A_1) \otimes \kappa(P)$ を $x_n \in B(A_2) \otimes \kappa(P)$ に対応させるが，
$Q$ の上では $x_n \in B(A_1) \otimes \kappa(Q)$ を
$x_{n + 1} \in B(A_2) \otimes \kappa(Q)$ に対応させる．得られる $A$-代数を
$B^{twist}(A)$ と書く．詳細は省略する．構成により $B^{twist}(A)$ は $A$ 上 Zariski 局所的に
ねじれていない版と同型である．実際，これは主開集合
$\Spec(A) \setminus \{P\}$ と主開集合 $\Spec(A) \setminus \{Q\}$ のいずれの上でも成り立つ．
しかし，この貼り合わせの選択は十分な「モノドロミー」を生じさせ，
$\Spec(B^{twist}(A))$ は連結になる（詳細は省略する）．最後に，中央のコピー
$\Spec(A) \to \Spec(B^{twist}(A))$ があり，これは $B^{twist}(A)$ 上 Zariski 局所的には
冪等元で生成されるイデアルにより切り出される閉部分スキームを与えるが，大域的には
そうではない（$B^{twist}(A)$ は非自明な冪等元をもたないからである）．
% P12-ERR-0357: frozen ill-attached English "cut out" phrase; see ERRATA_CANDIDATES.jsonl.

\begin{lemma}
\label{examples-lemma-not-generated-idempotents}
アフィンスキーム $X = \Spec(A)$ と閉部分スキーム $T \subset X$ であって，$T$ は $X$ 上
Zariski 局所的には冪等元で生成されるイデアルにより切り出されるが，$T$ 自身は冪等元で
生成されるイデアルにより切り出されないものが存在する．
\end{lemma}

\begin{proof}
上を参照せよ．
\end{proof}
% END EXPANDED MODULE ch110_44.tex
% BEGIN EXPANDED MODULE ch110_45.tex
\section{局所環を同一視するが ind-エタールでない環準同型}
\label{examples-section-not-ind-etale}
% P12-ERR-0352: frozen ambiguous double-relative English title; see ERRATA_CANDIDATES.jsonl.

\noindent
Section \ref{examples-section-ideal-locally-idempotents} で構成した環準同型 $R \to B(R)$ は，
$R$ のコピーの有限積の余極限であることに注意する．したがって $R \to B(R)$ は
ind-Zariski である．Pro-\'etale Cohomology, Definition
\ref{proetale-definition-ind-zariski} を参照せよ．次に，Section
\ref{examples-section-ideal-locally-idempotents} で構成した環準同型 $A \to B^{twist}(A)$ を考える．
この環準同型は $\Spec(A)$ 上 Zariski 局所的に ind-Zariski 環準同型
$R \to B(R)$ と同型なので，局所環を同一視すると結論できる
（Pro-\'etale Cohomology, Lemma \ref{proetale-lemma-ind-zariski-implies} を参照）．
Section \ref{examples-section-ideal-locally-idempotents} の議論により，冪等元で生成されない核をもつ切断
$B^{twist}(A) \to A$ が存在する．ここでもし $A \to B^{twist}(A)$ が ind-エタール，すなわち
$A \to A_i$ がエタールであるように $B^{twist}(A) = \colim A_i$ と書けたなら，
$A_i \to A$ の核は冪等元で生成されるはずである
（Algebra, Lemmas \ref{algebra-lemma-map-between-etale} および
\ref{algebra-lemma-surjective-flat-finitely-presented}）．これは上で述べた結果に矛盾する．

\begin{lemma}
\label{examples-lemma-not-ind-etale}
局所環を同一視するが ind-エタールでない環準同型 $A \to B$ が存在する．したがって，
なおさら ind-Zariski ではない．
\end{lemma}

\begin{proof}
上の議論を参照せよ．
\end{proof}
% END EXPANDED MODULE ch110_45.tex
% BEGIN EXPANDED MODULE ch110_46.tex
\section{入射加群に付随する flasque でない準連接層}
\label{examples-section-nonflasque}

\noindent
この型のさらに多くの例については，Illusie が Verdier によるいくつかの例を説明している
\cite[Expos\'e II, Appendix I]{SGA6} を参照せよ．

\medskip\noindent
アフィンスキーム $X = \Spec(A)$ を考える．ここで
$$
A = k[x, y, z_1, z_2, \ldots]/(x^nz_n)
$$
は Properties, Example \ref{properties-example-does-not-work-in-general} の環である．
$I = (x) \subset A$ とおき，$X$ の準コンパクト開集合 $U = D(x)$ を考える．loc.\ cit.\ で見たように，
どの $n \geq 0$ に対しても $A$-加群準同型 $I^n \to A$ から来ない切断
$s \in \mathcal{O}_X(U)$ が存在する．

\medskip\noindent
$\alpha : A \to J$ を，$A$ から入射 $A$-加群への埋め込みとする．
$Q = J/\alpha(A)$ とおき，商写像を $\beta : J \to Q$ と書く．写像
$$
\Gamma(X, \widetilde{J})
\longrightarrow
\Gamma(U, \widetilde{J})
$$
は全射でないと主張する．実際，$\alpha(s)$ は像に属さないと主張する．これを見るため，
背理法で論じる．$U$ 上で $\alpha(s)$ に制限される元 $x \in J$ が存在すると仮定する．
すると $\beta(x) \in Q$ は $U$ 上で $0$ に制限される元である．したがって，ある $n$ に対して
$I^n\beta(x) = 0$ である．Properties, Lemma
\ref{properties-lemma-sections-over-quasi-compact-open-in-affine} を参照せよ．これは射
$\varphi : I^n \to A$，$h \mapsto \alpha^{-1}(hx)$ を得ることを意味する．この射
$\varphi$ が Properties, Lemma
\ref{properties-lemma-sections-over-quasi-compact-open-in-affine} の写像を通じて切断 $s$ を
与えることは容易に分かり，矛盾である．

\begin{lemma}
\label{examples-lemma-nonflasque}
アフィンスキーム $X = \Spec(A)$ と入射 $A$-加群 $J$ であって，$\widetilde{J}$ が $X$ 上
flasque 層でないものが存在する．$U$ が標準開であっても，制限写像
$\Gamma(X, \widetilde{J}) \to \Gamma(U, \widetilde{J})$ は全射とは限らない．
\end{lemma}

\begin{proof}
上を参照せよ．
\end{proof}

\noindent
実際，同様の構成により，付随する準連接層が準コンパクト開集合上で非零のコホモロジーを
もつ入射加群の例を得られる．具体的には，体 $k$ 上の環
$$
A = k[x, y, w_1, u_1, w_2, u_2, \ldots]/(x^nw_n, y^nu_n, u_n^2, w_n^2)
$$
から始める．$I$ を入射 $A$-加群として入射準同型 $A \to I$ を選ぶ．
$A_{xy} \subset I_{xy}$ の元 $1/xy$ は $I_x \oplus I_y \to I_{xy}$ の像に属さないと主張する．
背理法により，ある $n \geq 1$ と $i, j \in I$ に対して
$$
\frac{1}{xy} = \frac{i}{x^n} + \frac{j}{y^n}
$$
であると仮定する．分母を払うと，ある $m \geq 0$ に対して
$$
(xy)^{n + m - 1} = x^my^{n + m}i + x^{n + m}y^mj
$$
を得る．$u_{n + m}w_{n + m}$ を掛けると，$A$ において
$u_{n + m}w_{n + m}(xy)^{n + m - 1} = 0$ となり，所望の矛盾を得る．
% P12-ERR-0353: frozen proof omits the required nonvanishing assertion; see ERRATA_CANDIDATES.jsonl.
$U = D(x) \cup D(y) \subset X = \Spec(A)$ とする．任意の $A$-加群 $M$ に対して，
Mayer--Vietoris により完全系列
$$
0 \to H^0(U, \widetilde{M}) \to M_x \oplus M_y \to M_{xy} \to
H^1(U, \widetilde{M}) \to 0
$$
がある．したがって $H^1(U, \widetilde{I})$ は非零である．

\begin{lemma}
\label{examples-lemma-nonvanishing}
台となる位相空間が Noether であるアフィンスキーム $X = \Spec(A)$ と入射 $A$-加群 $I$ で，
$\widetilde{I}$ が $X$ のある準コンパクト開集合 $U$ 上で非消滅な $H^1$ をもつものが存在する．
\end{lemma}

\begin{proof}
上を参照せよ．位相空間として $\Spec(A) = \Spec(k[x, y])$ であることに注意する．
\end{proof}
% END EXPANDED MODULE ch110_46.tex
% BEGIN EXPANDED MODULE ch110_47.tex
\section{分離的でない平坦群スキーム}
\label{examples-section-non-separated-group-scheme}

\noindent
体上の任意の群スキームは分離的である．Groupoids, Lemma
\ref{groupoids-lemma-group-scheme-over-field-separated} を参照せよ．基底上の群スキームについては，
これは正しくない．

\medskip\noindent
$k$ を体とし，$S = \Spec(k[x]) = \mathbf{A}^1_k$ とする．$G$ を，$0$ を二重化したアフィン直線
（Schemes, Example \ref{schemes-example-affine-space-zero-doubled} を参照）とし，$S$ 上のスキームと
みなす．したがって $G \to S$ の各ファイバーは，一点集合または二元集合
（一方は $U$，他方は $V$ に属する）である．ゆえに，$U$ の元を群構造の零元とすることで，
これらのファイバーに群の構造を入れられる．より正確には，$G$ は $S$ に同型に写る二つの
開集合 $U, V$ をもち，$U \cap V$ は $S \setminus \{0\}$ に同型に写る．このとき
$$
G \times_S G = U \times_S U \cup V \times_S U \cup
U \times_S V \cup V \times_S V
$$
であり，各部分は $S$ と同型である．したがって，第一および最後の部分を $U$ に，中央の
二つの部分を $V$ に写す一意な $S$-射として乗法 $m : G \times_S G \to G$ を定義できる．
これは上の各点ごとの記述と一致する．これが群スキーム構造を定めることの確認は省略する．

\begin{lemma}
\label{examples-lemma-non-separated-group-scheme}
アフィン直線上有限型の平坦群スキームで，分離的でないものが存在する．
\end{lemma}

\begin{proof}
上の議論を参照せよ．
\end{proof}

\begin{lemma}
\label{examples-lemma-non-quasi-separated-group-scheme}
無限次元アフィン空間上有限型の平坦群スキームで，準分離的でないものが存在する．
\end{lemma}

\begin{proof}
上と同じ構成を，無限次元アフィン空間
$S = \mathbf{A}^\infty_k = \Spec k[x_1, x_2, \ldots]$ を基底とし，極大イデアル
$(x_1, x_2, \ldots)$ に対応する原点 $0 \in S$ を $G$ で二重化される閉点として行える．
得られる群スキーム $G \rightarrow S$ は，Schemes, Example
\ref{schemes-example-not-quasi-separated} で説明したように準分離的でない．
\end{proof}
% END EXPANDED MODULE ch110_47.tex
% BEGIN EXPANDED MODULE ch110_48.tex
\section{単位成分は平坦だが平坦でない群スキーム}
\label{examples-section-non-flat-group-scheme}

\noindent
$X \to S$ をスキームの単射とし，$G = S \amalg X$ とする．$m : G \times_S G \to G$ を
$S$-射
$$
G \times_S G = X \times_S X \amalg X \amalg X \amalg S
\longrightarrow G = X \amalg S
$$
であって，直和因子 $X \times_S X$ と $S$ を $S$ に写し，二つの直和因子 $X$ を恒等射で
$X$ に写すものとする．これは群法則を定める．これを見るには，写像
$G \times_S G \times_S G \to G$ として
$m \circ (m \times \text{id}_G) = m \circ (\text{id}_G \times m)$ を示さなければならない．
$G \times_S G \times_S G$ を上のように成分に分解すると，各 $8$ 個の部分への制限について
これを確認すればよい．各部分は $S$，$X$，$X \times_S X$，または
$X \times_S X \times_S X$ のいずれかに同型である．さらに，二つの写像はいずれも，
これらの部分をそれぞれ $S$，$X$，$S$，$X$ に写す．ここで $X \to S$ が単射であるため，
$X \times_S X \cong X$ および $X \times_S X \times_S X \cong X$ であり，いずれの場合にも
$S$-射 $S \to S$ または $X \to X$ はちょうど一つである．したがって
$m \circ (m \times \text{id}_G) = m \circ (\text{id}_G \times m)$ である．ゆえに，
$X \to S$ としてスキームの任意の平坦でない単射（例えば閉埋め込み）を取れば，
基底 $S$ 上の群スキームで，その単位成分は $S$（したがって平坦）だが，全体は平坦でない
例を得る．

\begin{lemma}
\label{examples-lemma-non-flat-group-scheme}
基底 $S$ 上の群スキーム $G$ で，その単位成分は $S$ 上平坦だが，その全体は $S$ 上
平坦でないものが存在する．
\end{lemma}

\begin{proof}
上の議論を参照せよ．
\end{proof}
% END EXPANDED MODULE ch110_48.tex
% BEGIN EXPANDED MODULE ch110_49.tex
\section{体上の分離的でない群代数空間}
\label{examples-section-non-separated-group-space}

\noindent
体上の任意の群スキームは分離的である．Groupoids, Lemma
\ref{groupoids-lemma-group-scheme-over-field-separated} を参照せよ．体上の群代数空間については
これは正しくない（肯定的な結果についてはこの節の末尾を参照せよ）．

\medskip\noindent
$k$ を標数零の体とする．Spaces, Example
\ref{spaces-example-affine-line-translation} の代数空間
$G = \mathbf{A}^1_k/\mathbf{Z}$ を考える．構成により，$G$ は $k$ 上のスキームの圏における
前層
$$
T \longmapsto \Gamma(T, \mathcal{O}_T) / \mathbf{Z}
$$
に付随する fppf 層である．前層上の明らかな加法則は加法 $m : G \times G \to G$ を誘導し，
$G$ を $\Spec(k)$ 上の群代数空間にする．$G$ は分離的でないことに注意する
（実際，準分離的でも局所分離的でもない）．一方，$G \to \Spec(k)$ は有限型である．

\begin{lemma}
\label{examples-lemma-non-separated-group-space}
体上有限型の群代数空間で，分離的でないものが存在する
（実際，準分離的でも局所分離的でもないものが存在する）．
\end{lemma}

\begin{proof}
上の議論を参照せよ．
\end{proof}

\noindent
肯定的な結果を述べる．群代数空間 $G$ が準分離的，局所分離的，またはより一般に decent な
代数空間ならば，$G$ は実際に分離的である．More on Groupoids in Spaces, Lemma
\ref{spaces-more-groupoids-lemma-group-scheme-over-field-separated} を参照せよ．さらに，
体上有限型で分離的な群代数空間は，More on Groupoids in Spaces, Lemma
\ref{spaces-more-groupoids-lemma-group-space-scheme-locally-finite-type-over-k} により実際に
スキームである．証明の考え方は，schematic locus が開かつ稠密であるということである．
Properties of Spaces, Proposition
\ref{spaces-properties-proposition-locally-quasi-separated-open-dense-scheme} または
Decent Spaces, Theorem \ref{decent-spaces-theorem-decent-open-dense-scheme} を参照せよ．
この開集合を平行移動することで，$G$ の各点がスキームである開近傍をもつと分かる．
% END EXPANDED MODULE ch110_49.tex
% BEGIN EXPANDED MODULE ch110_50.tex
\section{エタール射のファイバー内の点の間の特殊化}
\label{examples-section-specializations-fibre-etale}
% P12-ERR-0354: frozen English title omits connecting words; see ERRATA_CANDIDATES.jsonl.

\noindent
$f : X \to Y$ がスキームのエタール射，またはより一般に局所準有限射ならば，ファイバーの
点の間に特殊化はない．Morphisms, Lemma
\ref{morphisms-lemma-locally-quasi-finite-fibres} を参照せよ．しかし，代数空間の射については，
一般にこれは成り立たない．

\medskip\noindent
例を与えるため，$k$ を体とし，
$$
P = k[u, u^{-1}, y, \{x_n\}_{n \in \mathbf{Z}}].
$$
とおく．自己同型 $\tau$ が $\tau(u) = u$，$\tau(y) = uy$，および
$\tau(x_n) = x_{n + 1}$ という規則で与えられるとき，この自己同型が生成する $k$-代数準同型による
$\mathbf{Z}$ の $P$ への作用を考える．$d \geq 1$ に対して
$I_d = ((1 - u^d)y, x_n - x_{n + d}, n \in \mathbf{Z})$ とおく．すると
$V(I_d) \subset \Spec(P)$ は $\tau^d$ の不動点軌跡である．$S \subset P$ を，$y$ とすべての
% P12-ERR-0355: frozen English "fix point locus" wording; see ERRATA_CANDIDATES.jsonl.
$1 - u^d$，$d \in \mathbf{N}$ で生成される乗法的部分集合とする．すると
$\mathbf{Z}$ は $U = \Spec(S^{-1}P)$ に自由に作用する．$X = U/\mathbf{Z}$ を商代数空間とする．
Spaces, Definition \ref{spaces-definition-quotient} を参照せよ．

\medskip\noindent
$S^{-1}P$ の素イデアル $\mathfrak p_n = (x_n, x_{n + 1}, \ldots)$ を考える．
$\tau(\mathfrak p_n) = \mathfrak p_{n + 1}$ であることに注意する．したがって，これらはそれぞれ
これらはそれぞれ点 $\xi_n \in U$ を定め，$X$ における像はいずれも $X$ の同じ点 $x$ である．さらに特殊化
% P12-ERR-0356: frozen English agreement/article defect; see ERRATA_CANDIDATES.jsonl.
$$
\ldots \leadsto \xi_n \leadsto \xi_{n - 1} \leadsto \ldots
$$
がある．したがって $U \to X$ は約束した型の例である．

\begin{lemma}
\label{examples-lemma-specializations-fibre-etale}
代数空間のエタール射 $f : X \to Y$ と，$|X|$ における点の非自明な特殊化
$x \leadsto x'$ であって，$|Y|$ において $f(x) = f(x')$ となるものが存在する．
\end{lemma}

\begin{proof}
上の議論を参照せよ．
\end{proof}
% END EXPANDED MODULE ch110_50.tex
% BEGIN EXPANDED MODULE ch110_51.tex
\section{fppf トーサーでないトーサー}
\label{examples-section-torsor-not-fppf}

\noindent
Groupoids, Remark \ref{groupoids-remark-fun-with-torsors} では，任意の $G$-トーサーが fppf 位相に
関する $G$-トーサーであるかという問いを提起した．この節では，必ずしもそうではないことを示す．

\medskip\noindent
$k$ を体とする．以下，すべてのスキームとスタックは $k$ 上のものとする．群スキーム
$G \to \Spec(k)$ を
$$
G = (\mu_{2, k})^\infty =
\mu_{2, k} \times_k \mu_{2, k} \times_k \mu_{2, k} \times_k \ldots =
\lim_n (\mu_{2, k})^n
$$
とする．ここで $\mu_{2, k}$ は $\Spec(k)$ 上の二乗根の群スキームである．Groupoids, Example
% P12-ERR-0358: frozen nonstandard "second roots of unity" wording; see ERRATA_CANDIDATES.jsonl.
\ref{groupoids-example-roots-of-unity} を参照せよ．アフィンスキームの逆極限なので，$G$ はアフィン
群スキームである．実際，これは環 $k[t_1, t_2, t_3, \ldots]/(t_i^2 - 1)$ のスペクトルである．
乗法写像 $m : G \times_k G \to G$ は代数の側では $t_i \mapsto t_i \otimes t_i$ により与えられる．

\medskip\noindent
$k$ 上の任意の $G$-トーサーは，ある $a_i \in k^*$ に対する
$$
P = \Spec(k[x_1, x_2, x_3, \ldots]/(x_i^2 - a_i))
$$
という形であり，その $G$-作用 $G \times_k P \to P$ は，代数の側で
$x_i \to t_i \otimes x_i$ により与えられると主張する．証明は省略する．実際，例のためには
$P$ が $G$-トーサーであることだけが必要であり，これは
$k' = k(\sqrt{a_1}, \sqrt{a_2}, \ldots)$ 上でスキーム $P$ が $G$-同変に $G$ と同型になることから
明らかである．$P$ が自明であるための必要十分条件は $k' = k$ である．実際，$P$ が
$k$-有理点をもてば，すべての $a_i$ が平方だからである．

\medskip\noindent
$P$ が fppf トーサーであるための必要十分条件は，体拡大
$k' = k(\sqrt{a_1}, \sqrt{a_2}, \ldots)/k$ が有限であることだと主張する．$k'$ が $k$ 上有限なら，
$\{\Spec(k') \to \Spec(k)\}$ は $P$ を自明化する fppf 被覆であり，したがって $P$ は確かに fppf
トーサーである．逆に，$P$ が fppf $G$-トーサーであると仮定する．これは，各 $P_{S_i}$ が
自明となる fppf 被覆 $\{S_i \to \Spec(k)\}$ が存在することを意味する．$S_i$ が空でないような
$i$ を選び，閉点 $s \in S_i$ を取る．Varieties, Lemma
\ref{varieties-lemma-locally-finite-type-Jacobson} により，体拡大 $\kappa(s)/k$ は有限であり，
構成により $P_{\kappa(s)}$ は $\kappa(s)$-有理点をもつ．したがって
$k \subset k' \subset \kappa(s)$ であり，$k'$ は $k$ 上有限である．

\medskip\noindent
明示的な例を得るには，例えば $k = \mathbf{Q}$ および $a_i = i$ と取ればよい
（望むなら，$a_i$ を $i$ 番目の素数としてもよい）．

\begin{lemma}
\label{examples-lemma-torsors-principal-spaces-not-equal}
$S$ をスキーム，$G$ を $S$ 上の群スキームとする．主等質 $G$-空間を分類するスタック
$G\textit{-Principal}$（Examples of Stacks, Subsection
\ref{examples-stacks-subsection-principal-homogeneous-spaces} を参照）と，fppf $G$-トーサーを分類する
スタック $G\textit{-Torsors}$（Examples of Stacks, Subsection
\ref{examples-stacks-subsection-fppf-torsors} を参照）は，一般には同値でない．
\end{lemma}

\begin{proof}
上の議論は，関手 $G\textit{-Torsors} \to G\textit{-Principal}$ が一般には本質的全射でないことを
示している．
\end{proof}
% END EXPANDED MODULE ch110_51.tex
% BEGIN EXPANDED MODULE ch110_52.tex
\section{準コンパクト平坦被覆をもつが代数的でないスタック}
\label{examples-section-not-algebraic-stack}

\noindent
この節では Brian Conrad による例を簡潔に記述する．この例はオンラインの
\href{https://mathoverflow.net/questions/15082/fpqc-covers-of-stacks/15269#15269}{この場所}
で見られる．ここで扱う例はわずかに異なる．

\medskip\noindent
$k$ を代数閉体とする．以下，すべてのスキームとスタックは $k$ 上のものとする．
$G \to \Spec(k)$ をアフィン群スキームとする．Examples of Stacks, Lemma
\ref{examples-stacks-lemma-classifying-stacks} では，$\mathcal{X} = [\Spec(k)/G]$ を
$(\Sch/\Spec(k))_{fppf}$ 上の亜群のスタックとみなす互いに同値ないくつかの方法を与えた．
特に $\mathcal{X}$ は fppf $G$-トーサーを分類する．より正確には，$1$-射
$T \to \mathcal{X}$ は $T$ 上の fppf $G_T$-トーサー $P$ に対応し，$2$-矢はトーサーの同型に
対応する．したがって対角 $1$-射
$$
\Delta :
\mathcal{X}
\longrightarrow
\mathcal{X} \times_{\Spec(k)} \mathcal{X}
$$
は表現可能かつアフィンである．実際，スキーム $T/k$ 上の任意の fppf $G_T$-トーサーの対
$P_1, P_2$ に対し，スキーム $\mathit{Isom}(P_1, P_2)$ は $T$ 上アフィンである．
$\Spec(k)$ 上の自明な $G$-トーサーは $1$-射
$$
f : \Spec(k) \longrightarrow \mathcal{X}.
$$
を定める．これは全射な $1$-射であると主張する．理由は単純で，定義により任意の $1$-射
$T \to \mathcal{X}$ に対し，$P_{T_i}$ が自明な $G_{T_i}$-トーサーと同型になるような fppf 被覆
$\{T_i \to T\}$ が存在する．したがって合成 $T_i \to T \to \mathcal{X}$ は $f$ を経由する．
ゆえに射影 $T \times_\mathcal{X} \Spec(k) \to T$ は明らかに全射である
（これが $f$ の全射性の定義である．Algebraic Stacks, Definition
\ref{algebraic-definition-relative-representable-property} を参照）．同様にして，$f$ は
準コンパクトかつ平坦であることを示せる（詳細は省略する）．さらに，$k$ 上のスキームとして
$$
\Spec(k) \times_\mathcal{X} \Spec(k) \cong G
$$
であることも記録しておく．

\medskip\noindent
$U$ がスキームであるような全射滑らかな射 $p : U \to \mathcal{X}$ が存在すると仮定する．
ファイバー積
$$
\xymatrix{
W \ar[d] \ar[r] & U \ar[d] \\
\Spec(k) \ar[r] & \mathcal{X}
}
$$
を考える．すると $W$ は $k$ 上の空でない滑らかなスキームであり，したがって $k$-点をもつ．
これは $f$ が $U$ を経由することを意味する．ゆえに
$$
G \cong
\Spec(k) \times_\mathcal{X} \Spec(k) \cong
(\Spec(k) \times_k \Spec(k))
\times_{(U \times_k U)}
(U \times_\mathcal{X} U)
$$
を得る．射影 $U \times_\mathcal{X} U \to U$ は滑らかと仮定したので，Morphisms, Lemma
\ref{morphisms-lemma-permanence-finite-type} により，
$U \times_\mathcal{X} U \to U \times_k U$ は局所有限型であると結論する．したがって，この場合
$G$ は $k$ 上局所有限型である．以上により次の補題を証明した（これは大幅に一般化できる）．

\begin{lemma}
\label{examples-lemma-BG-algebraic}
$k$ を体，$G$ を $k$ 上のアフィン群スキームとする．スタック $[\Spec(k)/G]$ がスキームによる
滑らかな被覆をもつならば，$G$ は $k$ 上有限型である．
\end{lemma}

\begin{proof}
上の議論を参照せよ．
\end{proof}

\noindent
この節の表題どおりの明示例を得るには，例えば Section
\ref{examples-section-torsor-not-fppf} の群スキーム $G = (\mu_{2, k})^\infty$ を取ればよい．これは
$k$ 上局所有限型でない．上の議論により，$\mathcal{X} = [\Spec(k)/G]$ は
Algebraic Stacks, Definition \ref{algebraic-definition-algebraic-stack} の性質 (1) と (2) を
満たすが，性質 (3) を満たさない．したがって $\mathcal{X}$ は代数スタックでない．一方，
全射，平坦，準コンパクトな射 $U \to \mathcal{X}$ をもつスキーム $U$ は存在する．すなわち，
上で調べた射 $f : \Spec(k) \to \mathcal{X}$ である．
% END EXPANDED MODULE ch110_52.tex
% BEGIN EXPANDED MODULE ch110_53.tex
\section{対象について極限保存だが，極限保存でない例}
\label{examples-section-limit-preserving}

\noindent
$S$ を空でないスキーム，$\mathcal{G}$ を $(\Sch/S)_{fppf}$ 上の入射 Abel 層とする．
Examples of Stacks, Lemma
\ref{examples-stacks-lemma-torsors-sheaf-stack-in-groupoids} により，$S$ 上の亜群のスタック
$$
\mathcal{G}\textit{-Torsors} \longrightarrow (\Sch/S)_{fppf}
$$
を得る．すべての $\mathcal{G}$-トーサーは自明なので，このスタックは
$(\Sch/S)_{fppf}$ 上で対象について極限保存である
（Criteria for Representability, Section \ref{criteria-section-limit-preserving} を参照）．一方，
$\mathcal{G}$ は層として極限保存とは限らないため，$\mathcal{G}\textit{-Torsors}$ は一般には
極限保存でない（Artin's Axioms, Definition \ref{artin-definition-limit-preserving} を参照）．
例えば，任意の非零入射層 $\mathcal{I}$ を取り，
$\mathcal{G} = \prod_{n \in \mathbf{Z}} \mathcal{I}$ とおけば例を得る．

\begin{lemma}
\label{examples-lemma-limit-preserving-on-objects-not-limit-preserving}
$S$ を空でないスキームとする．亜群のスタック $p : \mathcal{X} \to (\Sch/S)_{fppf}$ であって，
$p$ は対象について極限保存だが，$\mathcal{X}$ は極限保存でないものが存在する．
\end{lemma}

\begin{proof}
上の議論を参照せよ．
\end{proof}
% END EXPANDED MODULE ch110_53.tex
% BEGIN EXPANDED MODULE ch110_54.tex
\section{代数的でない分類スタック}
\label{examples-section-non-algebraic}
% P12-ERR-0362: frozen duplicate label shared with the later sheaf section; see ERRATA_CANDIDATES.jsonl.

\noindent
$S = \Spec(\mathbf{F}_p)$ とし，$\mu_p$ を $S$ 上の $p$ 乗根の群スキームとする．
Groupoids in Spaces, Section \ref{spaces-groupoids-section-stacks} では商スタック $[S/\mu_p]$ を導入し，
Examples of Stacks, Section \ref{examples-stacks-section-group-quotient-stacks} では
$[S/\mu_p]$ が fppf $\mu_p$-トーサーの分類スタックであることを示した．すなわち，
$S$ 上のスキーム $T$ が与えられると，圏 $\Mor_S(T, [S/\mu_p])$ は $T$ 上の fppf
$\mu_p$-トーサーの圏と標準的に同値である．最後に Criteria for Representability, Theorem
\ref{criteria-theorem-flat-groupoid-gives-algebraic-stack} で，$[S/\mu_p]$ が代数スタックであることを見た．

\medskip\noindent
そこで「エタール $\mu_p$-トーサーを分類する亜群にファイバー化された圏 $\mathcal{S}$ はどうか」
と問える（言い換えると，$\mathcal{S}$ は $\Sch/S$ 上の圏で，スキーム $T$ 上のファイバー圏は
$T$ 上のエタール $\mu_p$-トーサーの圏である）．

\medskip\noindent
第一の異議は，これが fppf 位相に対するスタックでないことである．対象の降下が成り立たないからである．
例えば $T = \Spec(\mathbf{F}_p(T))$ 上の $\mu_p$-トーサー
$\Spec(\mathbf{F}_p(t)[x]/(x^p - t))$ は fppf 局所的には自明だが，エタール局所的には自明でない．
% P12-ERR-0359: frozen t/T inconsistency in the base field; see ERRATA_CANDIDATES.jsonl.

\medskip\noindent
この第一の問題を直すにはエタール位相で作業すればよく，この場合には対象の降下が成り立つ．
実際，$\mathcal{S}$ は $(\Sch/S)_\etale$ 上の亜群のスタックである．さらに対角
$\Delta : \mathcal{S} \to \mathcal{S} \times \mathcal{S}$ は（スキームにより）表現可能である．
これは，二つの $\mu_p$-トーサー（エタール局所的に自明であるか否かを問わず）が与えられたとき，
それらの間の同型の層がスキームで表現可能だからである．

\medskip\noindent
そこで最後に，スキーム $U$ と滑らかかつ全射な $1$-射 $U \to \mathcal{S}$ が存在するか問える．
これが不可能であることを二通りに示す．直接の議論（読者には飛ばすことを勧める）と，一般結果を
用いる議論である．

\medskip\noindent
直接の議論（概略）：$1$-射 $\mathcal{S} \to \Spec(\mathbf{F}_p)$ は，形式的滑らかさの無限小
持ち上げ判定法を満たすことに注意する．実際，スキームの一階無限小肥厚 $T \to T'$ が与えられると，
$\mu_p(T') \to \mu_p(T)$ の核は $T'$ における $T$ のイデアル層の切断と同型であり，したがって
$H^1_\etale(T, \mu_p) = H^1_\etale(T', \mu_p)$ である．さらに，$\mathcal{S}$ は極限保存スタックである．
ゆえに $U \to \mathcal{S}$ が滑らかならば，$U \to \Spec(\mathbf{F}_p)$ は極限保存であり，
形式的滑らかさの無限小持ち上げ判定法を満たす．
% P12-ERR-0360: frozen duplicated English interword space; see ERRATA_CANDIDATES.jsonl.
したがって $U$ は $\mathbf{F}_p$ 上滑らかである．特に $U$ は被約なので
$H^1_\etale(U, \mu_p) = 0$ である．ゆえに $U \to \mathcal{S}$ は
$U \to \Spec(\mathbf{F}_p) \to \mathcal{S}$ と分解し，第一の矢は滑らかである．滑らかさの降下により，
$U \to \mathcal{S}$ が滑らかなら $\Spec(\mathbf{F}_p) \to \mathcal{S}$ も滑らかでなければならない．
しかしこれは成り立たない．なぜなら
$\Spec(\mathbf{F}_p) \times_\mathcal{S} \Spec(\mathbf{F}_p)$ は $\mu_p$ であり，
$\Spec(\mathbf{F}_p)$ 上滑らかでないからである．

\medskip\noindent
構造的な議論：Criteria for Representability, Section
\ref{criteria-section-stacks-etale} では，代数スタックを，エタール位相に関する亜群のスタックで，
その対角が代数空間により表現可能で滑らかな被覆をもつものと考えられることを見た．したがって，
滑らかで全射な $U \to \mathcal{S}$ が存在すれば，$\mathcal{S}$ は代数スタックであり，特に fppf
位相における降下を満たす．しかし上で見たように，$\mathcal{S}$ は fppf 位相における降下を満たさない．
% P12-ERR-0361: frozen English does-not-satisfies agreement defect; see ERRATA_CANDIDATES.jsonl.

\medskip\noindent
大まかに言えば，上の議論は，基底 $B$ 上の群スキーム $G$ に対し，エタール局所的に自明なトーサーを
分類するエタール位相の分類スタックが代数的であるための必要十分条件が，$G$ が $B$ 上滑らかである
ことだと示している．fppf 位相で作業する利点の一つは，$G \to B$ が平坦かつ局所有限表示であると
仮定すれば十分なことである．実際，fppf 位相に関する商スタック $[B/G]$ が代数的であるための
必要十分条件は $G \to B$ が平坦かつ局所有限表示であることである．Criteria for Representability,
Lemma \ref{criteria-lemma-BG-algebraic} を参照せよ．
% END EXPANDED MODULE ch110_54.tex
% BEGIN EXPANDED MODULE ch110_55.tex
\section{準コンパクト平坦被覆をもつが代数的でない層}
\label{examples-section-not-algebraic}
% P12-ERR-0362: frozen duplicate label shared with the preceding classifying-stack section; see ERRATA_CANDIDATES.jsonl.

\noindent
関手 $F = (\mathbf{P}^1)^\infty$ を考える．すなわち，スキーム $T$ に対して $F(T)$ は，各
$f_i : T \to \mathbf{P}^1$ がスキームの射であるような $f = (f_1, f_2, f_3, \ldots)$ 全体の集合である．
$\mathbf{P}^1$ は fpqc 被覆に関する層条件を満たすことに注意する．Descent, Lemma
\ref{descent-lemma-fpqc-universal-effective-epimorphisms} を参照せよ．層の積は層なので，$F$ も fpqc
位相に関する層条件を満たす．$F$ の対角は表現可能である．実際，射 $f : T \to F$ と
$g : S \to F$ が与えられると，$T \times_F S$ はスキーム $T \times S$ 内の閉部分スキーム
$T \times_{f_i, \mathbf{P}^1, g_i} S$ のスキーム論的共通部分である．全射滑らかなアフィン射
$\text{SL}_2 \to \mathbf{P}^1$ を備えた群スキーム $\text{SL}_2$ を考える．次に，標準的な積射
$U \to F$ をもつ $U = (\text{SL}_2)^\infty$ を考える．$U$ はアフィンスキームであることに注意する．
射 $U \to F$ は平坦，全射，かつ普遍開であると主張する．実際，射 $f : T \to F$ が与えられたとする．
このとき $Z = T \times_F U$ は，$T$ 上のスキーム
$Z_i = T \times_{f_i, \mathbf{P}^1} \text{SL}_2$ の無限ファイバー積である．各射
$Z_i \to T$ は全射，滑らか，かつアフィンであるため，
$$
Z = Z_1 \times_T Z_2 \times_T Z_3 \times_T \ldots
$$
は $Z$ 上平坦かつアフィンなスキームである．単純な極限の議論により $Z \to T$ も開である．
% P12-ERR-0363: frozen tautological "flat and affine over Z" base; see ERRATA_CANDIDATES.jsonl.

\medskip\noindent
一方，$F$ は代数空間でないと主張する．実際，$F$ が代数空間ならば，これは準コンパクトで分離的な
（$F$ 上のファイバー積の記述による）代数空間となる．したがって準連接層のコホモロジーは，ある
次数より上で消滅するはずである（Cohomology of Spaces, Proposition
\ref{spaces-cohomology-proposition-vanishing} およびその前の注意を参照）．しかし，切断をもつ射影
$$
(\mathbf{P}^1)^\infty \longrightarrow (\mathbf{P}^1)^n
$$
のもとで $\mathcal{O}(-2, -2, \ldots, -2)$ を引き戻すことにより，次数 $n$ でコホモロジーが
非零となる準連接層を明らかに得られる．以上により，MathOverflow で Anton Geraschenko が提起した
問いへの答えを得る．

\begin{lemma}
\label{examples-lemma-not-algebraic}
fpqc 位相に関する層条件を満たし，表現可能な対角 $\Delta : F \to F \times F$ をもち，かつ
$U$ がスキームであるような全射，平坦，普遍開，準コンパクトな射 $U \to F$ が存在するが，
$F$ は代数空間ではないような関手 $F : \Sch^{opp} \to \textit{Sets}$ が存在する．
\end{lemma}

\begin{proof}
上の議論を参照せよ．
\end{proof}
% END EXPANDED MODULE ch110_55.tex
% BEGIN EXPANDED MODULE ch110_56.tex
\section{エタール位相と Zariski 被覆・有限エタール被覆の比較}
\label{examples-section-etale-zariski-finite-etale}

\noindent
この節では，C.~Deninger の問いに答えて B.~Bhatt が見いだした例を与える．これは
$\mathbf{C}$ 上の代数多様体のエタール位相が，開埋め込みと有限エタール射で生成される位相より
真に細かいことを示す．このような例は専門家には確実に知られており，おそらく Artin と
Grothendieck にまで遡る\footnote{伝聞によれば，Grothendieck は当初，Zariski 開被覆と有限エタール
被覆を用いてエタール位相を定義しようとしたが，Artin がそうでないことを彼に納得させた
（おそらく，ここに記録するような例に基づいていた）．}．公刊された例を知っているなら，
\href{mailto:stacks.project@gmail.com}{stacks.project@gmail.com} まで電子メールを送ってほしい．

\medskip\noindent
$X$ を $\mathbf{C}$ 上の滑らかで連結なアフィン曲線とする．滑らかな連結曲線 $Y$ からの
次数 $3$ の有限射 $g: Y \to X$ で，非分岐点 $u$（次数 1）と分岐点 $y$（次数 2）からなる
ファイバー $g^{-1}(x) = \{u, y\}$ をもつ点 $x \in X$ が存在するものを選ぶ．$X$ を $x$ の
Zariski 開近傍で置き換えることで，$g$ は $y$ 以外のすべての点で不分岐と仮定してよい．
$U = Y \setminus \{y\}$ とする．すると $f=g|_U: U \to X$ はエタール全射であり，
$f^{-1}(x) = \{u\}$ である．

\begin{lemma}
\label{examples-lemma-etale-not-dominated-zariski-finite-etale}
エタール被覆 $U \to X$ は，開埋め込みと有限エタール射の有限個の合成によって細分できない．
\end{lemma}

\begin{proof}
背理法のため，射の塔
$$
W_m \to W_{m-1} \to \dots \to W_1 \to W_0 = X
$$
が存在し，各写像 $W_i \to W_{i - 1}$ は開埋め込みまたは有限エタール射であり，さらに
$f : U \to X$ に沿った $W_m \to X$ の持ち上げ $W_m \to U$ と，$u \in U$（したがって
$x \in X$）に写る点 $w \in W_m$ が存在すると仮定する．各 $W_k$ を，$w$ の像
$w_k \in W_k$ を含む連結成分で置き換えることにより，各 $W_k$ は滑らかで連結と仮定してよい．

\medskip\noindent
各 $0 \le i \le m$ に対し，ファイバー積 $Z_i = Y \times_X W_i$ を考える．
$W_i \to X$ はエタールで $Y$ は滑らかなので，曲線 $Z_i$ は滑らかな連結曲線の非交和であり，
基底変換により写像 $Z_i \to W_i$ は次数 $3$ の有限射である．$w_i$ 上の
$Z_i \to W_i$ のファイバーは，$Y$ の $y, u$ に写る二点 $z_i, v_i$ からなる．
射 $Z_i \to W_i$ は $v_i$ で不分岐であり，$z_i$ では次数 $2$ で分岐する．
$n_i = \#\pi_0(Z_i)$ とおく．$Z_i \to W_i$ のファイバーが $2$ 点をもつので
$n_i \in \{1, 2\}$ である．さらに，遷移写像 $Z_i \to Z_{i - 1}$ は支配的なので，$n_i$ は
$i$ とともに増加することしかできない．明らかに $n_0 = 1$ である．さらに $n_m \geq 2$ で
なければならない．実際，$X$-射 $W_m \to U \subset Y$ があるので，基底変換
$Z_m \to W_m$ は切断をもつ．したがって，$n_{i - 1} = 1$ かつ $n_i = 2$ となる最小の添字
$i \ge 1$ が存在する．$n_j$ は $Z_i$ が既約かどうかによって定まるので，写像
$W_i \to W_{i - 1}$ は開埋め込みではあり得ず，有限エタールでなければならない．
% P12-ERR-0364: frozen Z_i/Z_j subscript inconsistency; see ERRATA_CANDIDATES.jsonl.
$n_i = 2$ なので，滑らかな曲線 $Z_i$ は二つの連結成分をもつ．$W_i$ への写像は有限なので，
これらの成分の一つは $W_i$ に同型に写らなければならない．したがって
$Z_i = W_i \sqcup Z'_i$ と書け，$Z'_i \to W_i$ の次数は $2$ である．誘導写像
$W_i \subset Z_i \to Z_{i - 1}$ は，いずれも $W_{i - 1}$ 上有限である連結曲線の間の
支配的写像である．このような写像は全射でなければならない．ゆえに $z_{i - 1}$ に写る点
$t \in W_i$ を選べる．離散付値環の拡大
$$
\mathcal{O}_{W_{i - 1}, w_{i - 1}} \to
\mathcal{O}_{Z_i, z_i} \to
\mathcal{O}_{W_i, z}
$$
を得る．第一の拡大は分岐する一方，$W_i \to W_{i - 1}$ はエタールなので合成は不分岐である．
これは不可能であり，矛盾を得る．
% P12-ERR-0365: frozen local-ring points conflict with the chosen t; see ERRATA_CANDIDATES.jsonl.
\end{proof}
% END EXPANDED MODULE ch110_56.tex
% BEGIN EXPANDED MODULE ch110_57.tex
\section{層と特殊化}
\label{examples-section-sheaves}

\noindent
以下では，Topologies, Definition \ref{topologies-definition-big-etale-site} で構成した大エタールサイト
$\Sch_\etale$ を固定する．さらに，スキームはこのサイトの対象であるものとする．スキーム $X$ の点
$x, x'$ に対し，$x \in \overline{\{x'\}}$ ならば $x$ は $x'$ の {\it 特殊化} であるといい，
$x' \leadsto x$ と書くことを思い出そう．特に $x = x'$ のときこれは成り立つ．

\medskip\noindent
次の規則で定義される関手 $F : \Sch_\etale \to \textit{Ab}$ を考える：
$$
F(X) = \prod\nolimits_{x \in X}
\prod\nolimits_{x' \in X, x' \leadsto x, x' \not = x}
\mathbf{Z}/2\mathbf{Z}
$$
スキーム $X$ に対して，その台集合を $|X|$ と書く．元 $a \in F(X)$ は集合の写像
$|X| \times |X| \to \mathbf{Z}/2\mathbf{Z}$，$(x, x') \mapsto a(x, x')$ とみなす．これは
$x = x'$ の場合，または $x$ が $x'$ の特殊化でない場合には零である．スキームの射
$f : X \to Y$ が与えられたとき，
$$
F(f) : F(Y) \longrightarrow F(X)
$$
を，$b \in F(Y)$ に対して
$$
F(f)(b)(x, x') =
\left\{
\begin{matrix}
0 & \text{if }x\text{ is not a specialization of }x' \\
b(f(x), f(x')) & \text{else.}
\end{matrix}
\right.
$$
とする規則で定める．これは実際に $F(X)$ の元を定めることに注意する．合成可能な射
$f : X \to Y$ と $g : Y \to Z$ に対して $F(f) \circ F(g) = F(g \circ f)$ であると主張する．
実際，$c \in F(Z)$ とし，$X$ の点の特殊化 $x' \leadsto x$ を取ると，
$$
F(g \circ f)(x, x') = c(g(f(x)), g(f(x'))) = F(g)(F(f)(c))(x, x')
$$
である．なぜなら $f(x') \leadsto f(x)$ だからである（これは $f(x) = f(x')$ の場合にも成り立つ）．
% P12-ERR-0366: frozen ill-typed functoriality display; see ERRATA_CANDIDATES.jsonl.

\medskip\noindent
$G$ をエタール位相における $F$ の層化とする．

\medskip\noindent
$X$ がスキームで $x' \leadsto x$ が特殊化かつ $x' \not = x$ ならば，$G(X) \not = 0$ であると
主張する．実際，$a$ を，関数 $|X| \times |X| \to \mathbf{Z}/2\mathbf{Z}$ とみなしたとき
$(x, x')$ で非零となる元 $a \in F(X)$ とする．$\{f_i : U_i \to X\}$ を $X$ のエタール被覆とする．
$f_i(u_i) = x$ となる $i$ と点 $u_i \in U_i$ を選べる．一般化は平坦射に沿って持ち上がるので
（Morphisms, Lemma \ref{morphisms-lemma-generalizations-lift-flat} を参照），
$f_i(u'_i) = x'$ となる特殊化 $u'_i \leadsto u_i$ を取れる．上の構成により
$F(f_i)(a) \not = 0$ と分かる．したがって $a$ は $G(X)$ の非零元を定める．

\medskip\noindent
$X = \Spec(k)$ で $k$ が体なら（より一般には，すべての素イデアルが極大である環なら），
$F(X) = 0$ であり，任意のエタール射 $U \to X$ に対しても $F(U) = 0$ であることに注意する．
エタール射のファイバー内の相異なる点の間に特殊化がないからである．したがって $G(X) = 0$ である．

\medskip\noindent
$X \subset X'$ が肥厚であると仮定する．More on Morphisms, Definition
\ref{more-morphisms-definition-thickening} を参照せよ．基底変換関手
$U' \mapsto U = U' \times_{X'} X$ は，$X'$ 上エタールなスキームの圏から $X$ 上エタールな
スキームの圏への同値である．\'Etale Cohomology, Theorem
\ref{etale-cohomology-theorem-topological-invariance} を参照せよ．この状況では常に
$F(U) = F(U')$ なので，$G(X) = G(X')$ でもある．

\medskip\noindent
変種として，スキーム $X$ に写像 $a : |X|^{n + 1} \to \mathbf{Z}/2\mathbf{Z}$ 全体を対応させる
前層 $F_n$ を考えられる．ただし $a(x_0, \ldots, x_n)$ は，
$x_n \leadsto \ldots \leadsto x_0$ が特殊化の列であり，かつ
$x_n \not = x_{n - 1} \not = \ldots \not = x_0$ である場合にのみ非零とする．
$G_n$ を $F_n$ に付随する層とする．上とまったく同じ方法で，$\dim(X) \geq n$ なら $G_n$ は
非零であり，$\dim(X) < n$ なら零であることを示せる．

\begin{lemma}
\label{examples-lemma-sheaf-zero-on-low-dimension}
$\Sch_\etale$ 上の Abel 群の層 $G$ で，次の性質をもつものが存在する：
\begin{enumerate}
\item $\dim(X) < n$ ならば常に $G(X) = 0$ である，
\item $\dim(X) \geq n$ ならば $G(X)$ は零でない，かつ
\item $X \subset X'$ が肥厚ならば $G(X) = G(X')$ である．
\end{enumerate}
\end{lemma}

\begin{proof}
上の議論を参照せよ．
\end{proof}

\begin{remark}
\label{examples-remark-specialization}
いくつか注意を述べる：
\begin{enumerate}
\item 前層 $F$ と $F_n$ は分離前層である．
\item $F$, $F_n$ は層ではないことが分かる．
\item $G$, $G_n$ は実際には fppf 位相に関する層であることを示せる．
% P12-ERR-0367: frozen plural-subject agreement defect; see ERRATA_CANDIDATES.jsonl.
\end{enumerate}
必要になれば，これらの結果を証明する．
\end{remark}
% END EXPANDED MODULE ch110_57.tex
% BEGIN EXPANDED MODULE ch110_58.tex
\section{層と構成可能関数}
\label{examples-section-constructible-functions}

\noindent
以下では，Topologies, Definition \ref{topologies-definition-big-etale-site} で構成した大エタールサイト
$\Sch_\etale$ を固定する．さらに，スキームはこのサイトの対象であるものとする．この節では，
スキーム $X$ の {\it 構成可能分割} とは，局所有限な非交和分解
$X = \coprod_{i \in I} X_i$ で，各 $X_i \subset X$ が $X$ の局所構成可能部分集合であるものをいう．
局所有限とは，任意の準コンパクト開集合 $U \subset X$ に対し，$X_i \cap U$ が空でないような
$i \in I$ が有限個しかないことを意味する．$f : X \to Y$ がスキームの射で
$Y = \coprod Y_j$ が構成可能分割ならば，$X = \coprod f^{-1}(Y_j)$ は $X$ の構成可能分割である．
集合 $S$ とスキーム $X$ が与えられたとき，{\it 構成可能関数} $f : |X| \to S$ とは，
$X = \coprod_{s \in S} f^{-1}(s)$ が $X$ の構成可能分割となる写像をいう．$G$ が（抽象）群で
$a, b : |X| \to G$ が構成可能関数ならば，$ab : |X| \to G, x \mapsto a(x)b(x)$ も構成可能関数である．
その理由は，任意の二つの構成可能分割に対し，両方を細分する第三の分割が存在することにある．

\medskip\noindent
$A$ を任意の Abel 群とする．任意のスキーム $X$ に対して
$$
F(X) =
\frac{\{a : |X| \to A \mid a \text{ is a constructible function}\}}{\{\text{locally constant functions }|X| \to A\}}
$$
と定義する．$F(X)$ の元 $a$ は，局所定数関数を加える違いを除いて定まる関数と単に考える．
スキームの射 $f : X \to Y$ と元 $b \in F(Y)$ が与えられたとき，$F(f)(b) = b \circ f$ と定める．
こうして $F$ は $\Sch_\etale$ 上の前層となる．

\medskip\noindent
$\{f_i : U_i \to X\}$ が fppf 被覆で，$a \in F(X)$ が各添字に対して
$F(f_i)(a) = 0$ を $F(U_i)$ において満たすならば，$a \circ f_i$ は各 $i$ に対して局所定数関数である．
これはさらに，射 $f_i$ が開なので $a$ が局所定数関数であることを意味する．したがって
$F(X)$ において $a = 0$ である．ゆえに $F$ は分離前層である
（fppf 位相において，したがってなおさらエタール位相において）．

\medskip\noindent
$G$ をエタール位相における $F$ の層化とする．$F$ は分離的であり，例えば $X$ が離散付値環の
スペクトルならば $F(X) \not = 0$ なので，$G$ は零でない．

\medskip\noindent
$X = \Spec(k)$ とし，$k$ を体とする．このとき $X$ の任意のエタール被覆は，$k'/k$ が体の有限分離
拡大であるような被覆 $\{\Spec(k') \to \Spec(k)\}$ により支配される．
$F(\Spec(k')) = 0$ なので，$G(X) = 0$ と分かる．

\medskip\noindent
$X \subset X'$ が肥厚であると仮定する．More on Morphisms, Definition
\ref{more-morphisms-definition-thickening} を参照せよ．基底変換関手
$U' \mapsto U = U' \times_{X'} X$ は，$X'$ 上エタールなスキームの圏から $X$ 上エタールな
スキームの圏への同値である．\'Etale Cohomology, Theorem
\ref{etale-cohomology-theorem-topological-invariance} を参照せよ．この状況では
$F(U) = F(U')$ なので，$G(X) = G(X')$ でもある．

\medskip\noindent
層 $G$ は極限保存である．Limits of Spaces, Definition
\ref{spaces-limits-definition-locally-finite-presentation} を参照せよ．実際，環 $R$ が環の有向余極限
$R = \colim_i R_i$ として書かれているとする．$X = \Spec(R)$ および $X_i = \Spec(R_i)$ とおけば，
$X = \lim_i X_i$ である．このとき $G(X) = \colim_i G(X_i)$ である．これを証明するには，まず
$\Spec(R)$ の構成可能分割が，ある $\Spec(R_i)$ の構成可能分割から来ることを示す．
% P12-ERR-0368: frozen singular-article/plural-noun defect; see ERRATA_CANDIDATES.jsonl.
これにより $F$ に対する結果を得る．層化に対する結果を得るには，任意のエタール環準同型
$R \to R'$ が，ある $i$ に対するエタール環準同型 $R_i \to R_i'$ から来ることを用いる．詳細は省略する．

\begin{lemma}
\label{examples-lemma-weird-sheaf}
$\Sch_\etale$ 上の Abel 群の層 $G$ で，次の性質をもつものが存在する：
\begin{enumerate}
\item $k$ が体ならば常に $G(\Spec(k)) = 0$ である，
\item $G$ は極限保存である，
\item $X \subset X'$ が肥厚ならば $G(X) = G(X')$ である，かつ
\item $G$ は零でない．
\end{enumerate}
\end{lemma}

\begin{proof}
上の議論を参照せよ．
\end{proof}
% END EXPANDED MODULE ch110_58.tex
% BEGIN EXPANDED MODULE ch110_59.tex
\section{lisse-エタールサイトは関手的でない}
\label{examples-section-lisse-etale-not-functorial}

\noindent
$X$ の {\it lisse-エタール} サイト $X_{lisse,\etale}$ は，$X$ 上滑らかなスキームの圏に
（通常の）エタール被覆を入れたものである．Cohomology of Stacks, Section
\ref{stacks-cohomology-section-lisse-etale} を参照せよ．$f : X  \to Y$ をスキームの射とする．
連続関手
$$
u : Y_{lisse,\etale} \longrightarrow X_{lisse,\etale},\quad
V/Y \longmapsto V \times_Y X
$$
がある．したがって随伴関手の対
$$
u^s :
\Sh(X_{lisse,\etale})
\longrightarrow
\Sh(Y_{lisse,\etale}),
\quad
u_s :
\Sh(Y_{lisse,\etale})
\longrightarrow
\Sh(X_{lisse,\etale}),
$$
を得る．Sites, Section \ref{sites-section-continuous-functors} を参照せよ．一般に $u$ はサイトの射を
定め{\bf ない}と主張する．Sites, Definition \ref{sites-definition-morphism-sites} を参照せよ．
言い換えると，$u_s$ は一般には左完全でないと主張する．表現可能前層は lisse-エタールサイト上の
層であることに注意する．したがって Sites, Lemma
\ref{sites-lemma-pullback-representable-sheaf} により $u_sh_V = h_{V \times_Y X}$ である．いま，
$Y$ 上滑らかなスキーム $V_1, V_2$ の二つの射
$$
\xymatrix{
V_1 \ar[rd] \ar@<1ex>[rr]^a \ar@<-1ex>[rr]_b & & V_2 \ar[ld] \\
& Y
}
$$
を考える．もし $u_s$ が左完全ならば，
$$
u_s \text{Equalizer}(h_a, h_b : h_{V_1} \to h_{V_2})
=
\text{Equalizer}(h_{a \times 1}, h_{b \times 1} :
h_{V_1 \times_Y X} \to h_{V_2 \times_Y X})
$$
となるはずである．射 $a, b : V_1 \to V_2$ として，$Y$ 上滑らかなスキームから
$(a = b) \subset V_1$ への射が存在しない，すなわち左辺が空層となるが，$X$ への基底変換後には
等化子が空でなく $X$ 上滑らかとなるものを取る．簡単な例は，
$X = \Spec(\mathbf{F}_p)$，$Y = \Spec(\mathbf{Z})$，および
$V_1 = V_2 = \mathbf{A}^1_\mathbf{Z}$ とし，射を $a(x) = x$，$b(x) = x + p$ とすることである．
$a$ と $b$ の等化子は $(p)$ 上の $\mathbf{A}^1_\mathbf{Z}$ のファイバーであることに注意する．

\begin{lemma}
\label{examples-lemma-lisse-etale-not-functorial}
lisse-エタールサイトは，スキームの射についてさえ関手的でない．
\end{lemma}

\begin{proof}
上の議論を参照せよ．
\end{proof}
% END EXPANDED MODULE ch110_59.tex
% BEGIN EXPANDED MODULE ch110_60.tex
\section{Noether スキームの圏上の層}
\label{examples-section-sheaves-locally-Noetherian}

\noindent
$S$ を局所 Noether スキームとする．Artin's Axioms, Section
\ref{artin-section-noetherian} と同様に，$S$ 上局所 Noether なスキームの fppf サイトから
$S$ の大 fppf サイトへの包含関手
$$
u : (\textit{Noetherian}/S)_{fppf} \longrightarrow (\Sch/S)_{fppf}
$$
を考える．参照した節で説明したように，この関手は連続である．したがって随伴関手の対
$$
u^s :
\Sh((\Sch/S)_{fppf})
\longrightarrow
\Sh((\textit{Noetherian}/S)_{fppf})
$$
および
$$
u_s :
\Sh((\textit{Noetherian}/S)_{fppf})
\longrightarrow
\Sh((\Sch/S)_{fppf})
$$
を得る．Sites, Section \ref{sites-section-continuous-functors} を参照せよ．しかし一般に $u$ は
サイトの射を定めないと主張する．Sites, Definition \ref{sites-definition-morphism-sites} を参照せよ．
言い換えると，関手 $u_s$ は一般には左完全でないと主張する．

\medskip\noindent
$p$ を素数とし，$S = \Spec(\mathbf{F}_p)$ とおく．$(\textit{Noetherian}/S)_{fppf}$ 上の層の
単射
$$
a : \mathcal{F} \longrightarrow \mathcal{G}
$$
を次のように定義する．$U$ を $S$ 上局所 Noether なスキームとして，
$$
\mathcal{G}(U) = \Gamma(U, \mathcal{O}_U)^* =
\Mor_S(U, \mathbf{G}_{m, S})
$$
と定め，
$$
\mathcal{F}(U) = \{f \in \mathcal{G}(U) \mid \text{fppf locally }f
\text{ has arbitrary }p\text{-power roots}\}
$$
とする．Noether $\mathbf{F}_p$-代数 $A$ の冪零根基 $I \subset A$ は冪零であり，$A$ における
$1$ の $p$-冪根は $1 + a$，$a \in I$ という形の元である．したがって $1$ の非自明な
$p$-冪根のうち，任意の $p$-冪根をもつものはない．
% P12-ERR-0369: frozen spurious English "of"; see ERRATA_CANDIDATES.jsonl.
% P12-ERR-0370: frozen double-negative hyphenation; see ERRATA_CANDIDATES.jsonl.
ゆえに，任意の局所 Noether スキーム $U$ に対して $\mathcal{F}(U)$ は $p$-捩れのない Abel 群で
あると結論する．詳細の一部は省略する．したがって
$p : \mathcal{F} \to \mathcal{F}$ は $(\textit{Noetherian}/S)_{fppf}$ 上の Abel 層の単射である．

\medskip\noindent
矛盾を得るため，$u_s$ が完全であると仮定する．
% P12-ERR-0371: frozen proof assumes exactness although the claim is failure of left exactness; see ERRATA_CANDIDATES.jsonl.
すると $p : u_s\mathcal{F} \to u_s\mathcal{F}$ も単射であり，$S$ 上の任意の $V$ に対して
$(u_s\mathcal{F})(V)$ は $p$-捩れのない Abel 群である．表現可能前層は fppf サイト上の層なので，
Sites, Lemma \ref{sites-lemma-pullback-representable-sheaf} により，$u_s\mathcal{G}$ は
$\mathbf{G}_{m, S}$ で表現される．$u_s\mathcal{F} \to u_s\mathcal{G}$ が単射であることから，
$S$ 上の各スキーム $V$ に対し，$p$-捩れのない部分群
$$
(u_s\mathcal{F})(V) \subset \Gamma(V, \mathcal{O}_V)^*
$$
で，次の性質をもつものを得る．$U$ が局所 Noether であるような $S$ 上のスキームの任意の射
$V \to U$ に対し，部分群
$$
\mathcal{F}(U) \subset \Gamma(U, \mathcal{O}_U)^*
$$
は制限写像 $\Gamma(U, \mathcal{O}_U)^* \to \Gamma(V, \mathcal{O}_V)^*$ によって部分群
$(u_s\mathcal{F})(V)$ の中に写る．

\medskip\noindent
実際の矛盾は次のように得られる．
$k = \bigcup_{n \geq 0} \mathbf{F}_p(t^{1/{p^n}})$ とし，
$$
B = k \otimes_{\mathbf{F}_p(t)} k
$$
および $V = \Spec(B)$ とおく．二つの余射 $k \to B$ に対応する二つの射影
$V \to \Spec(k)$ があり，$\Spec(k)$ は Noether なので，部分群
$$
(u_s\mathcal{F})(V) \subset B^*
$$
は $k^* \otimes 1$ と $1 \otimes k^*$ を含むと結論する．しかし
$$
(t^{1/p} \otimes 1) \cdot (1 \otimes t^{-1/p}) =
t^{1/p} \otimes t^{-1/p}
$$
は $B$ の非自明な $p$-捩れ単元なので，これは矛盾である．

\begin{lemma}
\label{examples-lemma-not-a-morphism-of-sites-noetherian-to-all}
$S = \Spec(\mathbf{F}_p)$ のとき，包含関手
$(\textit{Noetherian}/S)_{fppf} \to (\Sch/S)_{fppf}$ はサイトの射を定めない．
\end{lemma}

\begin{proof}
上の議論を参照せよ．
\end{proof}
% END EXPANDED MODULE ch110_60.tex
% BEGIN EXPANDED MODULE ch110_61.tex
\section{準連接加群の導来順像}
\label{examples-section-derived-push-quasi-coherent}

\noindent
$k$ を標数 $p > 0$ の体とし，$S = \Spec(k[x])$ とする．$G = \mathbf{Z}/p\mathbf{Z}$ を抽象群，
または $S$ 上の定数群スキームとみなす．$G$ が $S$ に自明に作用するときの代数スタック
$\mathcal{X} = [S/G]$ を考える．Examples of Stacks, Remark
\ref{examples-stacks-remark-X-mod-G-group} および Criteria for Representability, Lemma
\ref{criteria-lemma-BG-algebraic} を参照せよ．構造射
$$
f : \mathcal{X} \longrightarrow S
$$
を考える．この射は準コンパクトかつ準分離的である．したがって関手
$$
Rf_{\QCoh, *} :
D^{+}_\QCoh(\mathcal{O}_\mathcal{X})
\longrightarrow
D^{+}_\QCoh(\mathcal{O}_S),
$$
を得る．Derived Categories of Stacks, Proposition
\ref{stacks-perfect-proposition-derived-direct-image-quasi-coherent} を参照せよ．
$Rf_{\QCoh, *}\mathcal{O}_\mathcal{X}$ を計算しよう．$D_\QCoh(\mathcal{O}_S)$ は
$k[x]$-加群の導来圏と同値なので（Derived Categories of Schemes, Lemma
\ref{perfect-lemma-affine-compare-bounded} を参照），これは
$R\Gamma(\mathcal{X}, \mathcal{O}_\mathcal{X})$ の計算と同値である．そのために被覆
$S \to \mathcal{X}$ とスペクトル系列
$$
H^q(S \times_\mathcal{X} \ldots \times_\mathcal{X} S, O)
\Rightarrow H^{p + q}(\mathcal{X}, \mathcal{O}_\mathcal{X})
$$
を使える．
% P12-ERR-0372: frozen plain O coefficient in the spectral sequence; see ERRATA_CANDIDATES.jsonl.
Cohomology of Stacks, Proposition
\ref{stacks-cohomology-proposition-smooth-covering-compute-cohomology} を参照せよ．
$$
S \times_\mathcal{X} \ldots \times_\mathcal{X} S = S \times G^p
$$
はアフィンであることに注意する．したがって複体
$$
k[x] \to \text{Map}(G, k[x]) \to \text{Map}(G^2, k[x]) \to \ldots
$$
は $R\Gamma(\mathcal{X}, \mathcal{O}_\mathcal{X})$ を計算する．ここで
$\varphi \in \text{Map}(G^{p - 1}, k[x])$ に対し，その微分は
$(g_1, \ldots, g_p)$ を
$$
\varphi(g_2, \ldots, g_p) +
\sum\nolimits_{i = 1}^{p - 1}
(-1)^i\varphi(g_1, \ldots, g_i + g_{i + 1}, \ldots, g_p)
+ (-1)^p\varphi(g_1, \ldots, g_{p - 1}).
$$
に写す写像である．これは $G$ が $k[x]$ に自明に作用するときの群コホモロジーを計算する複体に
ほかならない（将来の参照をここに挿入）．巡回群 $G$ の $k[x]$ におけるコホモロジーは，
各コホモロジー次数 $\geq 0$ にちょうど一つの $k[x]$ のコピーである
（将来の参照をここに挿入）．
% P12-ERR-0373: frozen pair of unresolved editorial placeholders; see ERRATA_CANDIDATES.jsonl.
したがって
$$
Rf_*\mathcal{O}_\mathcal{X} = \bigoplus\nolimits_{n \geq 0} \mathcal{O}_S[-n]
$$
と結論する．次に複体
$$
E = \bigoplus\nolimits_{m \geq 0} \mathcal{O}_\mathcal{X}[m]
$$
を考える．これは $D_\QCoh(\mathcal{O}_\mathcal{X})$ の対象である．一般結果のために議論を中断する．

\begin{lemma}
\label{examples-lemma-is-limit}
$\mathcal{X}$ を代数スタックとする．$K$ を $D(\mathcal{O}_\mathcal{X})$ の対象で，そのコホモロジー層が
局所準連接であり（Sheaves on Stacks, Definition
\ref{stacks-sheaves-definition-locally-quasi-coherent}），かつ平坦基底変換性
（Cohomology of Stacks, Definition \ref{stacks-cohomology-definition-flat-base-change}）を満たすものとする．
このとき $D(\mathcal{O}_\mathcal{X})$ に識別三角形
$$
K \to
\prod\nolimits_{n \geq 0} \tau_{\geq -n} K \to
\prod\nolimits_{n \geq 0} \tau_{\geq -n} K \to K[1]
$$
が存在する．言い換えると，$K$ はその標準切断の導来逆極限である．
\end{lemma}

\begin{proof}
$\mathcal{X}$ の「大 fppf サイト」$\mathcal{X}_{fppf}$ 上で作業していることを思い出そう
（Sheaves on Stacks と Cohomology on Stacks の章における $\mathcal{O}_\mathcal{X}$-加群の層に関する
規約による）．$\mathcal{B}$ を，アフィンスキーム $U$ の上にある $\mathcal{X}_{fppf}$ の対象 $x$ 全体の
集合とする．Sheaves on Stacks, Lemmas
\ref{stacks-sheaves-lemma-compare-fppf-etale}，
\ref{stacks-sheaves-lemma-cohomology-restriction}，Descent, Lemma
\ref{descent-lemma-quasi-coherent-and-flat-base-change}，および Cohomology of Schemes, Lemma
\ref{coherent-lemma-quasi-coherent-affine-cohomology-zero} を合わせると，$\mathcal{F}$ が局所準連接で
$x \in \mathcal{B}$ ならば $H^p(x, \mathcal{F}) = 0$ であると分かる．したがって主張は，
$d = 0$ とした Cohomology on Sites, Lemma
\ref{sites-cohomology-lemma-is-limit-dimension} から従う．
\end{proof}

\begin{lemma}
\label{examples-lemma-sum-is-product}
$\mathcal{X}$ を代数スタックとする．$\mathcal{F}_n$ が $\mathcal{X}$ 上の，平坦基底変換性をもつ
局所準連接層の族ならば，$D(\mathcal{O}_\mathcal{X})$ において
$\oplus_n \mathcal{F}_n[n] \to \prod_n \mathcal{F}_n[n]$ は同型である．
\end{lemma}

\begin{proof}
Lemma \ref{examples-lemma-is-limit} により，直和が直積と同型であると分かるので正しい．
\end{proof}

\noindent
議論を続ける．準連接加群は局所準連接であり，平坦基底変換性を満たす
（Sheaves on Stacks, Lemma \ref{stacks-sheaves-lemma-quasi-coherent}）ので，
$$
E = \prod\nolimits_{m \geq 0} \mathcal{O}_\mathcal{X}[m]
$$
を得る．コホモロジーは逆極限と可換するので，
$$
Rf_*E = \prod\nolimits_{m \geq 0}
\left(\bigoplus\nolimits_{n \geq 0} \mathcal{O}_S[m - n]\right)
$$
と分かる．この複体は $D_\QCoh(\mathcal{O}_S)$ の対象でないことに注意する．実際，次数 $0$ の
コホモロジー層は $\mathcal{O}_S$ のコピーの無限直積であり，これは局所準連接
$\mathcal{O}_S$-加群ですらない．

\begin{lemma}
\label{examples-lemma-push-not-OK}
代数スタックの準コンパクトかつ準分離的な射 $f : \mathcal{X} \to \mathcal{Y}$ は，関手
$Rf_* : D_\QCoh(\mathcal{O}_\mathcal{X}) \to
D_\QCoh(\mathcal{O}_\mathcal{Y})$ を誘導するとは限らない．
\end{lemma}

\begin{proof}
上の議論を参照せよ．
\end{proof}
% END EXPANDED MODULE ch110_61.tex
% BEGIN EXPANDED MODULE ch110_62.tex
\section{大きな Abel 圏}
\label{examples-section-big}

\noindent
この節の目的は，「大きな」Abel 圏 $\mathcal{A}$ と対象 $M, N$ で，拡大の同型類の集まり
$\Ext_\mathcal{A}(M, N)$ が集合でないものの例を与えることである．
% P12-ERR-0374: frozen missing cohomological degree on Ext; see ERRATA_CANDIDATES.jsonl.
この例は Freyd による．\cite[page 131, Exercise A]{Freyd} を参照せよ．

\medskip\noindent
$\mathcal{A}$ を次のように定義する．$\mathcal{A}$ の対象は三つ組 $(M, \alpha, f)$ からなり，
$M$ は Abel 群，$\alpha$ は順序数，$f : \alpha \to \text{End}(M)$ は写像である．射
$(M, \alpha, f) \to (M', \alpha', f')$ は Abel 群の準同型 $\varphi : M \to M'$ であって，
{\it 任意の}順序数 $\beta$ に対して
$$
\varphi \circ f(\beta) = f'(\beta) \circ \varphi
$$
を満たすものとして与えられる．ここで，$\beta$ が $\alpha$ に属さない場合には $f(\beta) = 0$ とし，
同様に $\beta$ が $\alpha'$ の元でない場合には $f'(\beta)$ を零とする．こうして定義された圏が
Abel 圏であることの確認は省略する．

\medskip\noindent
対象 $Z = (\mathbf{Z}, \emptyset, f)$，すなわちすべての作用素が零である対象を考える．観察すべき
ことは，$\mathcal{A}$ で計算した群 $\Ext^1_\mathcal{A}(Z, Z)$ が真の類であり，集合ではないことである．
実際，各順序数 $\alpha$ に対し，台となる群が $M = \mathbf{Z} \oplus \mathbf{Z}$ で，$f$ の値が
$$
f(\alpha) =
\left(
\begin{matrix}
0 & 1 \\
0 & 0
\end{matrix}
\right).
$$
を除いて常に零となる，$Z$ による $Z$ の拡大 $(M, \alpha + 1, f)$ を見つけられる．これは明らかに
拡大の同型類の真の類を生じさせる．特に $\mathcal{A}$ の導来圏では射の集まりが真の類になる．
Derived Categories, Lemma \ref{derived-lemma-ext-1} を参照せよ．これは，三角圏の Verdier 商を
定義する際に注意を要することを意味する．

\begin{lemma}
\label{examples-lemma-big-abelian-category}
$\Ext$-群が真の類となる「大きな」Abel 圏 $\mathcal{A}$ が存在する．
\end{lemma}

\begin{proof}
上の議論を参照せよ．
\end{proof}
% END EXPANDED MODULE ch110_62.tex
% BEGIN EXPANDED MODULE ch110_63.tex
\section{弱随伴点とスキーム論的稠密性}
\label{examples-section-weak-ass-dense}

\noindent
$k$ を体とする．$R = k[z, x_i, y_i]/(z^2, zx_iy_i)$ とし，$i$ は $\mathbf{N}$ の元を走るものとする．
$R = R_0 \oplus M_0$ であることに注意する．ここで $R_0 = k[x_i, y_i]$ は部分環で，$M_0$ は
平方零イデアルであり，$R_0$-加群として $M_0 \cong R_0/(x_iy_i)$ である．素イデアル
$\mathfrak p = (z, x_i)$ は，$R$ を $R$-加群とみたときの弱随伴素イデアルである
（Algebra, Definition \ref{algebra-definition-weakly-associated}）．実際，$R_\mathfrak p$ の元 $z$ は
非零だが $\mathfrak pR_\mathfrak p$ によって零化される．一方，開部分スキーム
$$
U = \bigcup D(x_i) \subset \Spec(R) = S
$$
を考える．$U \subset S$ はスキーム論的に稠密であると主張する
（Morphisms, Definition \ref{morphisms-definition-scheme-theoretically-dense}）．これを証明するには，
$j : U \to S$ を包含射として，$\mathcal{O}_S \to j_*\mathcal{O}_U$ が単射であることを示せば十分である．
Morphisms, Lemma \ref{morphisms-lemma-characterize-scheme-theoretically-dense} を参照せよ．代数に戻すと，
すべての $g \in R$ に対して写像
$$
R_g \longrightarrow \prod R_{x_ig}
$$
が単射であることを示さなければならない．$g_0 \in R_0$，$m_0 \in M_0$ として
$g = g_0 + m_0$ と書く．このとき $R_g = R_{g_0}$ である（詳細は省略する）．したがって
$g \in R_0$ と仮定してよい．また $g$ は非零と仮定してよい．いま
$R_g = (R_0)_g \oplus (M_0)_g$ である．$R_0$ は整域なので，写像
$(R_0)_g \to \prod (R_0)_{x_ig}$ は単射である．$g \in (x_iy_i)$ ならば
$(M_0)_g = 0$ であり，示すことはない．$g \not \in (x_iy_i)$ ならば，$(x_iy_i)$ は $R_0$ の
根基イデアルなので，$M_0 \to \prod (M_0)_{x_ig}$ が単射であることを示せばよい．
$R_0 \to M_0 \to (M_0)_{x_n}$ の核は $(x_iy_i, y_n)$ である．$(x_iy_i, y_n)$ は根基イデアル
なので，$g \not \in (x_iy_i, y_n)$ ならば $R_0 \to M_0 \to (M_0)_{x_ng}$ の核は
$(x_iy_i, y_n)$ である．すべての $n \gg 0$ に対して $g \not \in (x_iy_i, y_n)$ なので，
核は所望どおり $\bigcap_{n \gg 0} (x_iy_i, y_n) = (x_iy_i)$ に含まれると結論する．

\medskip\noindent
Ofer Gabber による第二の例を述べる．$k$ を体とし，$R$ を関数 $\mathbf{N} \to k$ 全体の環，
$R'$ を最終的に定数となる関数 $\mathbf{N} \to k$ 全体の環とする．すると
$\Spec(R)$ は Stone--{\v C}ech コンパクト化\footnote{各元 $f \in R$ は，$u$ が単元，$e$ が冪等元で
あるように $ue$ の形である．すると Algebra, Lemma
\ref{algebra-lemma-ring-with-only-minimal-primes} により $\Spec(R)$ は Hausdorff である．一方，離散位相を
入れた $\mathbf{N}$ は稠密開部分集合とみなせる．Hausdorff 準コンパクト位相空間 $X$ への集合写像
$\mathbf{N} \to X$ が与えられると，環準同型 $\mathcal{C}^0(X; k) \to R$ を得る．ここで
$\mathcal{C}^0(X; k)$ は局所定数写像 $X \to k$ の $k$-代数である．これは
$\Spec(R) \to \Spec(\mathcal{C}^0(X; k)) = X$ を与え，普遍性を証明する．} $\beta\mathbf{N}$ であり，
$\Spec(R')$ は一点コンパクト化\footnote{ここでは，$R'$ には余分な極大イデアルが実際に一つしかない
ことを示す．} $\mathbf{N}^* = \mathbf{N} \cup \{\infty\}$ である．環 $R$ と $R'$ ではすべての
素イデアルが極小なので，すべての点が弱随伴点である．

\begin{lemma}
\label{examples-lemma-example-schematically-dense-missing-weakly-associated-point}
被約スキーム $X$ とスキーム論的に稠密な開集合 $U \subset X$ で，ある弱随伴点
$x \in X$ が $U$ に属さないものが存在する．
\end{lemma}

\begin{proof}
第一の例では，構成により $\mathfrak p \not \in U$ である．Gabber の例では，スキーム
$\Spec(R)$ または $\Spec(R')$ は被約である．
\end{proof}
% END EXPANDED MODULE ch110_63.tex
% BEGIN EXPANDED MODULE ch110_64.tex
\section{トレースの非加法性の例}
\label{examples-section-non-additive}

\noindent
$k$ を体とし，$R = k[\epsilon]$ を $k$ 上の双対数の環とする．言い換えると
$R = k[x]/(x^2)$ であり，$\epsilon$ は $R$ における $x$ の合同類である．複体の短完全列
$$
\xymatrix{
0 \ar[d] \ar[r] &
R \ar[d]^\epsilon \ar[r]_1 &
R \ar[d] \\
R \ar[r]^1 & R \ar[r] & 0
}
$$
を考える．ここで縦列が複体であり，第一行を次数 $0$，第二行を次数 $1$ に置く．第一の複体
（すなわち左の縦列）を $A^\bullet$，第二を $B^\bullet$，第三を $C^\bullet$ と書く．図式
\begin{equation}
\label{examples-equation-commutes-up-to-homotopy}
\vcenter{
\xymatrix{
A^\bullet \ar[d]_{1 + \epsilon} \ar[r] &
B^\bullet \ar[r] \ar[d]_1 &
C^\bullet \ar[d]_1 \\
A^\bullet \ar[r] & B^\bullet \ar[r] & C^\bullet
}
}
\end{equation}
は $K(R)$ で可換，すなわちホモトピーを除いて可換する複体の図式であると主張する．実際，右の
正方形は可換であり，左の正方形のずれはホモトピー $1 : A^1 \to B^0$ で与えられる．一方，
$$
\text{Tr}_{A^\bullet}(1 + \epsilon) + \text{Tr}_{C^\bullet}(1)
\not = \text{Tr}_{B^\bullet}(1).
$$

\begin{lemma}
\label{examples-lemma-nonadditivity-of-trace}
環 $R$，ホモトピー圏 $K(R)$ の識別三角形 $(K, L, M, \alpha, \beta, \gamma)$，およびこの識別三角形の
自己準同型 $(a, b, c)$ であって，$K$, $L$, $M$ は完全複体であり，
$\text{Tr}_K(a) + \text{Tr}_M(c) \not = \text{Tr}_L(b)$ となるものが存在する．
\end{lemma}

\begin{proof}
上の例を考える．写像 $\gamma : C^\bullet \to A^\bullet[1]$ は次数 $0$ における $\epsilon$ による
乗法で与えられる．Derived Categories, Definition
\ref{derived-definition-distinguished-triangle} を参照せよ．したがって
$$
\xymatrix{
C^\bullet \ar[d] \ar[r]_\gamma & A^\bullet[1] \ar[d] \\
C^\bullet \ar[r]^\gamma & A^\bullet[1]
}
$$
も $\epsilon(1 + \epsilon) = \epsilon$ により $K(R)$ で可換である．ゆえに，確かに識別三角形の射を得る．
\end{proof}
% END EXPANDED MODULE ch110_64.tex
% BEGIN EXPANDED MODULE ch110_65.tex
\section{射影スキームの構造層の有限でないコホモロジー}
\label{examples-section-nonfinite-h0}

\noindent
$R$ を任意の非 Noether 環，$I$ を有限生成でない $R$ のイデアルとする．
$S = R[x, y]/xyI$ を，$x$ と $y$ の次数が $1$ である次数付き環とみなす．すると
$X = \text{Proj}(S)$ は $R$ 上の射影スキームである．$H^0(X, \mathcal{O}_X)$ は有限生成
$R$-加群でないと主張する．

\medskip\noindent
環 $S_x$，$S_y$，$S_{xy}$ を二重次数で分解すると，
$$
S_x = \bigoplus_{\substack{(a, b) \in \mathbf{Z}^2 \\ b \geq 0}} R'_{a, b},
\quad\quad
S_y = \bigoplus_{\substack{(a, b) \in \mathbf{Z}^2 \\ a \geq 0}} R''_{a, b},
\quad\quad
S_{xy} = \bigoplus_{\substack{(a, b) \in \mathbf{Z}^2}} R/I,
$$
と書ける．ここで $(a, b)$-直和因子は次数 $a + b$ に置かれ，
$$
R'_{a, b} =
\left\{
\begin{matrix}
R & b = 0, \\
R/I & b \ne 0,
\end{matrix}
\right.
\quad\quad
R''_{a, b} =
\left\{
\begin{matrix}
R & a = 0, \\
R/I & a \ne 0.
\end{matrix}
\right.
$$
である．$R$-加群 $H^0(X, \mathcal{O}_X)$ は，$(F, G) \mapsto F - G$ で定義される写像
$$
(S_x)_0 \times (S_y)_0 \to (S_{xy})_0
$$
の核として計算できる．ここで記号 $A_0$ は，$\mathbf{Z}$-次数付き環 $A$ の次数 $0$ 成分を表す．
上の $S_x, S_y, S_{xy}$ の明示的記述により，この核は，$F - G \in I$ を満たす対 $(F, G)$ 全体から
なる部分加群 $M \subset R^2$ と同型である．$(F, G) \mapsto F - G$ で与えられる $R$-加群の写像
$M \to I$ は全射なので，$M$ は有限生成 $R$-加群でない．

\medskip\noindent
全射 $R[x, y] \to S$ は閉埋め込み $i : X \to \mathbf{P}^1_R$ に対応する．すると
$\mathcal{F} = i_*\mathcal{O}_X$ は有限型準連接 $\mathcal{O}_{\mathbf{P}^1_R}$-加群であり，その
$H^0$ は上で計算した有限生成でない加群 $M$ に等しい．

\begin{lemma}
\label{examples-lemma-nonfinite-h0}
有限でない $H^0$．
\begin{enumerate}
\item 環 $R$ と $R$ 上の射影スキーム $X$ で，$H^0(X, \mathcal{O}_X)$ が有限 $R$-加群でないものが存在する．
\item 環 $R$ と $\mathbf{P}^1_R$ 上の有限型準連接加群 $\mathcal{F}$ で，
$H^0(\mathbf{P}^1_R, \mathcal{F})$ が有限 $R$-加群でないものが存在する．
\end{enumerate}
\end{lemma}

\begin{proof}
上の議論を参照せよ．
\end{proof}
% END EXPANDED MODULE ch110_65.tex
% BEGIN EXPANDED MODULE ch110_66.tex
\section{射影的であることは基底上局所的でない}
\label{examples-section-non-descending-property-projective}

\noindent
降下の章では，射の多くの性質が，fpqc 位相においてさえ基底上局所的であることを見た．
Descent, Sections \ref{descent-section-descending-properties-morphisms}，
\ref{descent-section-descending-properties-morphisms-fpqc}，および
\ref{descent-section-descending-properties-morphisms-fppf} を参照せよ．射影性についてはこれは正しくない．

\begin{lemma}
\label{examples-lemma-non-descending-property-projective}
性質
\begin{enumerate}
\item[] $\mathcal{P}(f) =$「$f$ は射影的である」，および
\item[] $\mathcal{P}(f) =$「$f$ は準射影的である」
\end{enumerate}
は基底上 Zariski 局所的でない．したがって，なおさら基底上 fpqc 局所的でない．
\end{lemma}

\begin{proof}
Hironaka \cite[Example B.3.4.1]{H} に従い，滑らかな複素 3 次元多様体の固有射
$f:V_Y\to Y$ で，Zariski 局所的には射影的だが射影的でないものを定義する．$f$ は固有かつ
射影的でないので，準射影的でもない．

\medskip\noindent
$Y$ を複素数体 ${\mathbf C}$ 上の射影 3 空間とする．$C$ と $D$ を $Y$ の滑らかな二次曲線で，
閉部分スキーム $C\cap D$ が被約で二つの複素点 $P$ と $Q$ からなるものとする．例えば
$C=\{ [x,y,z,w]: xy=z^2, w=0\}$，$D=\{ [x,y,z,w]:
xy=w^2, z=0\}$，$P=[1,0,0,0]$，$Q=[0,1,0,0]$ とする．$Y-Q$ 上では，まず曲線 $C$ を
吹き上げ，次に曲線 $D$ の狭義変換を吹き上げる
（Divisors, Definition \ref{divisors-definition-strict-transform}）．$Y-P$ 上では，まず曲線 $D$ を
吹き上げ，次に曲線 $C$ の狭義変換を吹き上げる．$Y-P-Q$ 上では，構成した二つの多様体は標準的に
同型なので，$Y-P-Q$ 上で貼り合わせられる．得られる $V_Y$ は ${\mathbf C}$ 上の滑らかな固有
3 次元多様体である．射 $f:V_Y\to Y$ は固有であり，$Y-P$ と $Y-Q$ 上で吹き上げなので，
Divisors, Lemma \ref{divisors-lemma-blowing-up-projective} により Zariski 局所的に射影的である．
$V_Y$ が ${\mathbf C}$ 上射影的でないことを示す．それにより $f$ が射影的でないことが従う．

\medskip\noindent
そのため，$L$ を $C-P-Q$ の複素点の $V_Y$ における逆像，$M$ を $D-P-Q$ の複素点の逆像とする．
すると $L$ と $M$ は射影直線 ${\mathbf P}^1_{{\mathbf C}}$ と同型である．次に $E$ を
$C\cup D\subset Y$ の $V_Y$ における逆像，すなわち $V_Y$ 内のその逆像とする．したがって $E\rightarrow C\cup D$ は固有射で，
$(C\cup D)-\{P,Q\}$ 上のファイバーは ${\mathbf P}^1$ と同型である．$E$ における $P$ の逆像は
二直線 $L_0$ と $M_0$ の合併であり，$E$ 上のサイクルの有理同値
$L\sim L_0+M_0$ および $M\sim M_0$ がある（$C$ と $D$ が ${\mathbf P}^1$ と同型であることを用いる）．
二曲線を吹き上げた順序から生じる非対称性に注意する．$Q$ の近くでは逆のことが起こる．したがって
$Q$ の逆像は二直線 $L_0'$ と $M_0'$ の合併であり，$E$ 上に有理同値
$L\sim L_0'$ および $M\sim L_0'+M_0'$ がある．これらを合わせると，$E$ 上，したがって
$V_Y$ 上で $L_0+M_0'\sim 0$ を得る．もし $V_Y$ が ${\mathbf C}$ 上射影的なら，豊富な直線束
$H$ をもち，これは $V_Y$ のすべての曲線上で次数 $> 0$ をもつ．特に $H$ は
$L_0+M_0'$ 上で正の次数をもつはずである．これは，体上の固有スキーム上で直線束の次数が有理同値を
法とする 1-サイクル上でよく定義されることに矛盾する
（Chow Homology, Lemma \ref{chow-lemma-proper-pushforward-rational-equivalence} および
Lemma \ref{chow-lemma-factors}）．したがって $V_Y$ は ${\mathbf C}$ 上射影的でない．
\end{proof}

\noindent
別の用語では，Hironaka の 3 次元多様体 $V_Y$ は，被約部分スキーム $C\cup D$ に沿う $Y$ の
吹き上げ $Y'$ の小解消である．ここで $Y'$ は二つの通常二重点をもつ．$Y$ を $C$ に沿って，次に
$D$ の狭義変換に沿って吹き上げて $Z$ を定義すると，$Z$ は滑らかな射影 3 次元多様体であり，
射影的でない 3 次元多様体 $V_Y$ は $Y-P$ 上の「フロップ」によって $Z$ と異なる．
% END EXPANDED MODULE ch110_66.tex
% BEGIN EXPANDED MODULE ch110_67.tex
\section{射影スキームの非有効降下データ}
\label{examples-section-non-effective-descent-projective}

\noindent
降下の章では，fpqc 射に関するスキームの降下データが，いくつかの射の類について有効であることを見た．
特に，アフィン射，より一般に準アフィン射は fpqc 被覆に関する降下を満たす
（Descent, Lemma \ref{descent-lemma-quasi-affine}）．射影射についてはこれは正しくない．

\begin{lemma}
\label{examples-lemma-non-effective-descent-projective}
スキームのエタール被覆 $X\to S$ と，$X\to S$ に関する降下データ $(V/X,\varphi)$ で，
$V\to X$ は射影的だが，その降下データはスキームの圏で有効でないものが存在する．
\end{lemma}

\begin{proof}
滑らかで分離的な複素 3 次元代数空間でスキームでないものに関する Hironaka の例
\cite[Example B.3.4.2]{H} を模倣する．

\medskip\noindent
群 $G = \mathbf{Z}/2 = \{1, g\}$ の複素数体上の射影 3 空間 $\mathbf{P}^3$ への作用
$$
g[x,y,z,w] = [y,x,w,z].
$$
を考える．${\mathbf P}^3$ 内の互いに交わらない二直線の一つは
$L_1=\{ [x,x,z,z]\}$ であり，
\par\noindent
もう一つは $L_2=\{ [x,-x,z,-z]\}$ である．この作用はこれら二直線の外では自由である．
$Y={\mathbf P}^3-(L_1\cup L_2)$ とする．${\mathbf C}$ 上の滑らかな準射影スキーム $S=Y/G$ で，
$Y\to S$ が $G$-トーサーとなるものが存在する
（Groupoids, Definition \ref{groupoids-definition-principal-homogeneous-space}）．明示的には，
射
\begin{align*}
{\mathbf P}^3 & \to \text{Proj } {\mathbf C}[x,y,z,w]^G\\
   & = \text{Proj } {\mathbf C}[u_0,u_1,v_0,v_1,v_2]/(v_0v_1=v_2^2),
\end{align*}
による ${\mathbf P}^3$ 内の開部分集合 $Y$ の像として $S$ を定義できる．ここで
$u_0=x+y$，$u_1=z+w$，$v_0=(x-y)^2$，$v_1=(z-w)^2$，$v_2=(x-y)(z-w)$ であり，
環は $u_0,u_1$ の次数を 1，$v_0,v_1,v_2$ の次数を 2 として次数付けられる．

\medskip\noindent
$C=\{ [x,y,z,w]: xy=z^2, w=0\}$ および
$D=\{ [x,y,z,w]: xy=w^2, z=0\}$ とする．これらは ${\mathbf P}^3$ の滑らかな二次曲線で，
$G$-不変開部分集合 $Y$ に含まれ，$g(C)=D$ である．また $C\cap D$ は二点
$P:=[1,0,0,0]$ と $Q:=[0,1,0,0]$ からなり，この二点は $G$ の作用によって交換される．

\medskip\noindent
$V_Y\to Y$ を，$Y-P$ 上ではまず $D$ を吹き上げ，次に $C$ の狭義変換を吹き上げることで，
$Y-Q$ 上ではまず $C$ を吹き上げ，次に $D$ の狭義変換を吹き上げることで定義されるスキームとする
（これは Lemma \ref{examples-lemma-non-descending-property-projective} の証明と同じ構成だが，ここでの $Y$ は
${\mathbf P}^3$ の開部分集合を表し，${\mathbf P}^3$ 全体ではない）．すると $Y$ 上の $G$ の作用は $V_Y$ 上の
$G$ の作用へ持ち上がり，$Y-P$ と $Y-Q$ の逆像を交換する．$V_Y$ 上のこの $G$-作用は，
$G$-トーサー $Y\to S$ に関する $V_Y$ 上の降下データ $(V_Y/Y,\varphi_Y)$ を与える．射
$V_Y\to Y$ は，Lemma \ref{examples-lemma-non-descending-property-projective} の証明で示したように固有だが
射影的でない．

\medskip\noindent
$X$ を開部分集合 $Y-P$ と $Y-Q$ の非交和とする．すると全射エタール射 $X\to Y\to S$ がある．
$V$ を $V_Y\to Y$ の $X$ への引き戻しとする．$V_Y\to Y$ は各開部分集合 $Y-P$ と $Y-Q$ 上で
吹き上げなので，射 $V\to X$ は射影的である．さらに，降下データ $(V_Y/Y,\varphi_Y)$ は，
エタール被覆 $X\to S$ に関する降下データ $(V/X,\varphi)$ に引き戻される．

\medskip\noindent
この降下データがスキームの圏で有効であると仮定する．すなわち，その降下データとともに射
$V\to X$ に引き戻されるスキーム $U\to S$ が存在するとする．このとき $U$ は $V_Y$ の
$G$-作用による商となる．
$$
\xymatrix{
V \ar[r]\ar[d]&  X\ar[d] \\
V_Y \ar[r]\ar[d]& Y\ar[d] \\
U \ar[r]& S
}
$$

\medskip\noindent
$E$ を $C\cup D\subset Y$ の $V_Y$ における逆像とする．したがって $E\rightarrow C\cup D$ は固有射で，
$(C\cup D)-\{P,Q\}$ 上のファイバーは ${\mathbf P}^1$ と同型である．$E$ における $P$ の逆像は
二直線 $L_0$ と $M_0$ の合併である．したがって $E$ における $Q=g(P)$ の逆像は二直線
$L_0'=g(M_0)$ と $M_0'=g(L_0)$ の合併である．Lemma
\ref{examples-lemma-non-descending-property-projective} の証明で示したように，$E$ 上に有理同値
$L_0+M_0'=L_0+g(L_0)\sim 0$ がある．

\medskip\noindent
閉部分スキームの降下により，曲線 $L_1\subset U$（${\mathbf P}^1$ と同型）が存在し，
その $V_Y$ における逆像は $L_0\cup g(L_0)$ である
（閉埋め込みがアフィン射であることに注意し，Descent, Lemma
\ref{descent-lemma-affine} を用いる）．$R$ を $L_1$ の複素点とする．$U$ はスキームであると仮定した
ので，局所環 $O_{U,R}$ の関数 $f$ で，$R$ では消えるが曲線 $L_1$ 全体では消えないものを選べる．
% P12-ERR-0375: frozen plain O local-ring notation; see ERRATA_CANDIDATES.jsonl.
$D_{\text{loc}}$ を $\Spec O_{U,R}$ の閉部分集合 $\{f = 0\}$ の既約成分とする．すると
$D_{\text{loc}}$ の余次元は 1 である．$U$ における $D_{\text{loc}}$ の閉包は，点 $R$ を含むが
曲線 $L_1$ 全体を含まない $U$ の既約因子 $D_U$ である．$V_Y$ における $D_U$ の逆像は有効因子
$D$ であり，$L_0\cup g(L_0)$ と交わるが，曲線 $L_0$ と $g(L_0)$ のいずれも含まない．

\medskip\noindent
複素 3 次元多様体 $V_Y$ は滑らかなので，$O(D)$ は $V_Y$ 上の直線束である．ここでは正則局所環が
階乗的，言い換えると UFD であることを用いる．More on Algebra, Lemma
\ref{more-algebra-lemma-regular-local-UFD} を参照せよ．$O(D)$ の固有曲面 $E\subset V_Y$ への制限は，
$D$ に関する情報により 1-サイクル $L_0+g(L_0)$ 上で正の次数をもつ直線束である．
% P12-ERR-0376: frozen plain O(D) line-bundle notation; see ERRATA_CANDIDATES.jsonl.
$E$ 上で $L_0+g(L_0)\sim 0$ なので，これは，体上の固有スキーム上で直線束の次数が有理同値を法とする
1-サイクル上でよく定義されることに矛盾する
（Chow Homology, Lemma \ref{chow-lemma-proper-pushforward-rational-equivalence} および
Lemma \ref{chow-lemma-factors}）．したがって降下データ $(V/X,\varphi)$ は実際に有効でない．すなわち，
$U$ はスキームとして存在しない．
\end{proof}

\noindent
この例では，降下データは代数空間の圏では{\it 有効である}．より正確には，例えば Algebraic Spaces,
Lemma \ref{spaces-lemma-quotient} により，$U$ は ${\mathbf C}$ 上 3 次元の滑らかで分離的な代数空間として
存在する．Hironaka の 3 次元多様体 $U$ は，滑らかな準射影 3 次元多様体 $S$ を既約な節点曲線
$(C\cup D)/G$ に沿って吹き上げた $S'$ の小解消である．3 次元多様体 $S'$ は節点特異点をもつ．
$S'$ のもう一つの小解消（「フロップ」によって $U$ と異なる）も，やはりスキームでない代数空間である．
% END EXPANDED MODULE ch110_67.tex
% BEGIN EXPANDED MODULE ch110_68.tex
\section{全空間がスキームでない曲線族}
\label{examples-section-family-of-curves}

\noindent
Quot, Section \ref{quot-section-curves} では，スキーム $S$ 上の曲線族を，代数空間 $X$ から $S$ への
相対次元 $\leq 1$ の固有，平坦，有限表示な射と定義する．$S$ が完備 Noether 局所環のスペクトルなら，
$X$ はスキームである．More on Morphisms of Spaces, Lemma
\ref{spaces-more-morphisms-lemma-projective-over-complete} を参照せよ．この節では，一般にはこれは
正しくないことを示す．

\medskip\noindent
$k$ を体とする．まず，固有平坦射
$$
Y \longrightarrow \mathbf{A}^1_k
$$
と $0 \in \mathbf{A}^1_k(k)$ の上にある点 $y \in Y(k)$ で，次の性質をもつものから始める：
\begin{enumerate}
\item ファイバー $Y_0$ は $k$ 上の滑らかな幾何学的既約曲線である，
\item $\mathbf{A}^1_k$ を支配する任意の真の閉部分スキーム $T \subset Y$ に対し，交わり
$T \cap Y_0$ は $y$ と異なる点を少なくとも一つ含む．
\end{enumerate}
このような曲面が与えられたとき，次のように例を構成する．
$$
\xymatrix{
Y \ar[rd] & Z \ar[d] \ar[l] \ar[r] & X \ar[ld] \\
& \mathbf{A}^1_k
}
$$
ここで $Z \to Y$ は $y$ における $Y$ の吹き上げである．$E \subset Z$ を例外因子，$C \subset Z$ を
$Y_0$ の狭義変換とする．スキーム論的に $Z_0 = E \cup C$ である
（これを見るには，$Y$ が $y$ で滑らかで，さらに $Y \to \mathbf{A}^1_k$ が $y$ で滑らかなことを用いる）．
Artin の結果により（\cite{ArtinII}；$C$ の法束が負であることを見るには Semistable Reduction, Lemma
\ref{models-lemma-properties-form} を用いる），$Z$ 内の曲線 $C$ を吹き下ろして，図式のような代数空間
$X$ を得られる．$x \in X(k)$ を $C$ の像とする．

\medskip\noindent
$X$ はスキームでないと主張する．実際，もしスキームならば，$x$ のアフィン開近傍 $U \subset X$ が
存在する．$T = X \setminus U$ とおく．すると $T$ は $\mathbf{A}^1_k$ を支配する
（$X \to \mathbf{A}^1_k$ のファイバーは固有で次元 $1$ だが，$U \to \mathbf{A}^1_k$ のファイバーは
アフィンであり，したがって異なる）．$T' \subset Z$ を $T$ に同型に写る閉部分スキームとする
（$x \not \in T$ だからである）．すると $X$ における $T'$ の像は上の条件 (2) に矛盾する
（$T' \cap Z_0$ は吹き上げ $Z \to Y$ の例外因子 $E$ に含まれるからである）．
% P12-ERR-0377: frozen contradiction refers to image in X although condition (2) is on Y; see ERRATA_CANDIDATES.jsonl.

\medskip\noindent
議論を終えるには $Y$ を構成する必要がある．$k$ の標数は $3$ でないと仮定する．
$\mathbf{A}^1_k = \Spec(k[t])$ と書き，$\mathbf{P}^2_{k[t]}$ 内で
$$
Y \quad : \quad T_0^3 + T_1^3 + T_2^3 - tT_0T_1T_2 = 0
$$
とする．$t = 0$ におけるこのファイバーは，滑らかな射影種数 $1$ 曲線である．アフィン部分
$V_+(T_0)$ 上ではアフィン方程式
% P12-ERR-0378: frozen V_+(T_0) instead of the standard affine chart D_+(T_0); see ERRATA_CANDIDATES.jsonl.
$$
1 + x^3 + y^3 - txy = 0
$$
を得て，これは $k$ 上の滑らかな曲面を定める．対称性により他のアフィン部分でも同じことが
成り立つので，$Y$ は滑らかな曲面である．最後に，このアフィン方程式から分数体が $k(x, y)$ である
ことも分かり，したがって $Y$ は有理曲面である．有理曲面の Picard 群は有限生成である
（将来の参照をここに挿入）．
% P12-ERR-0379: frozen unresolved editorial placeholder; see ERRATA_CANDIDATES.jsonl.
したがって，性質 (2) をもつ $y \in Y_0(k)$ を選ぶには，すべての正整数に対して
\begin{equation}
\label{examples-equation-three}
\mathcal{O}_{Y_0}(ny) \not \in \Im(\Pic(Y) \to \Pic(Y_0))
\text{ for all }n > 0
\end{equation}
を満たす $y$ を選べば十分である．実際，(2) に反する $T$ の $1$ 次元既約成分の和は，ある
$n \geq 1$ に対し，因子 $ny$ における $Y_0$ の有効 Cartier 因子交叉を与え，
$\mathcal{O}_{Y_0}(ny)$ が制限写像の像に属すると結論することになる．
% P12-ERR-0380: frozen malformed Cartier-intersection sentence; see ERRATA_CANDIDATES.jsonl.
$Y_0$ の種数は $\geq 1$ なので，写像
$$
Y_0(k) \to \Pic(Y_0),\quad y \mapsto \mathcal{O}_{Y_0}(y)
$$
は単射であることに注意する．いま $k$ が非可算代数閉体ならば，$\Pic(Y)$ の可算性と今の注意を用いて，
(\ref{examples-equation-three})，したがって (2) を満たす $y \in Y_0(k)$ を見つけられる．

\begin{lemma}
\label{examples-lemma-family-of-curves-not-scheme}
体 $k$ と曲線族 $X \to \mathbf{A}^1_k$ で，$X$ がスキームでないものが存在する．
\end{lemma}

\begin{proof}
上の議論を参照せよ．
\end{proof}
% END EXPANDED MODULE ch110_68.tex
% BEGIN EXPANDED MODULE ch110_69.tex
\section{導来基底変換}
\label{examples-section-derived-base-change}

\noindent
$R \to R'$ を環準同型とする．More on Algebra, Section
\ref{more-algebra-section-derived-base-change} において，導来基底変換関手
$- \otimes_R^\mathbf{L} R' : D(R) \to D(R')$ を構成する．
次に，$R \to A$ をもう一つの環準同型とする．図式
$$
\xymatrix{
	A \ar[r] & A \otimes_R R' \ar@{=}[r] & A' \\
R \ar[u] \ar[r] & R' \ar[u] \ar[ur]
}
$$
を考える．$A$-加群 $M$ が与えられると，テンソル積 $M \otimes_R R'$ は
$A \otimes_R R'$-加群，すなわち $A'$-加群である．環準同型
$A \to A'$ に対して導来関手
$$
- \otimes_A^\mathbf{L} A' : D(A) \longrightarrow D(A')
$$
が存在するが，一般にはこの関手は $- \otimes_R^\mathbf{L} R'$ と一致しない．
より正確には，$K \in D(A)$ に対し，More on Algebra, Equation
(\ref{more-algebra-equation-comparison-map}) で構成される $D(R')$ における標準写像
$$
K \otimes_R^{\mathbf{L}} R' \longrightarrow
K \otimes_A^{\mathbf{L}} A'
$$
は一般には同型でない．そこで「導来基底変換関手」$T : D(A) \to D(A')$，すなわち
$T(K)$ が $D(R')$ において $K \otimes_R^\mathbf{L} R'$ に写るような関手が存在するか，
という疑問が生じる．この節では，一般にはそのような関手が存在しないことを示す．

\medskip\noindent
$k$ を体とする．$R = k[x, y]$，$R' = R/(xy)$，$A = R/(x^2)$ とおく．
$D(R')$ の対象 $A \otimes_R^\mathbf{L} R'$ は
$$
x^2 : R' \longrightarrow R'
$$
によって表され，$H^0(A \otimes_R^\mathbf{L} R') = A \otimes_R R'$ である．
$D(A \otimes_R R')$ の対象 $E$ で，$D(R')$ において
$A \otimes_R^\mathbf{L} R'$ に写るものは存在しないと主張する．実際，そのような $E$ に対しては
加群 $H^0(E)$ が自由となり，したがって $E$ は
$H^0(E)[0] \oplus H^{-1}(E)[1]$ と分解するはずである．しかし
$A \otimes_R^\mathbf{L} R'$ が $D(R')$ においてそのコホモロジー群の直和と同型でないことは
容易に分かる．

\begin{lemma}
\label{examples-lemma-no-derived-base-change}
$R \to R'$ および $R \to A$ を環準同型とする．一般には，三角圏の関手
$T : D(A) \to D(A \otimes_R R')$ であって，$A$-加群 $M$ に対し，
$D(A \otimes_R R') \to D(R')$ の下で $M \otimes_R^\mathbf{L} R'$ に写る
$D(A \otimes_R R')$ の対象 $T(M)$ を与えるものは存在しない．
\end{lemma}

\begin{proof}
上の議論を参照せよ．
\end{proof}
% END EXPANDED MODULE ch110_69.tex
% BEGIN EXPANDED MODULE ch110_70.tex
\section{興味深いコンパクト対象}
\label{examples-section-interesting-compact}


\noindent
$R$ を環とし，$(A, \text{d})$ を微分付き次数付き $R$-代数とする．
$A = R$ のとき，$D(A, \text{d}) = D(R)$ の任意のコンパクト対象は，有限生成射影加群からなる
有限複体によって表されることが知られている．言い換えると，コンパクト対象は perfect である．
More on Algebra, Proposition \ref{more-algebra-proposition-perfect-is-compact} を参照せよ．
微分付き次数付き加群の言葉における類似の問いは，次のようになる：
「$D(A, \text{d})$ の任意のコンパクト対象は，有限かつ次数付き射影的な微分付き次数付き
$A$-加群 $P$ によって表されるか？」

\medskip\noindent
一般の微分付き次数付き代数に対して，これは正しくない．実際，$k$ を標数 $2$ の体とする
（これにより符号を気にせずに済む）．
$A = k[x, y]/(y^2)$ とし，次の次数と微分を入れる：
\begin{enumerate}
\item $x$ の次数は $0$，
\item $y$ の次数は $-1$，
\item $\text{d}(x) = 0$，かつ
\item $\text{d}(y) = x^2 + x$．
\end{enumerate}
このとき $x : A \to A$ は $K(A, \text{d})$ における射影子である．したがって
$$
A = \Ker(x) \oplus \Im(1 - x)
$$
が $K(A, \text{d})$ において成り立つ．Differential Graded Algebra, Lemma
\ref{dga-lemma-homotopy-direct-sums} および Derived Categories, Lemma
\ref{derived-lemma-projectors-have-images-triangulated} を参照せよ．
$A$ が $D(A, \text{d})$ のコンパクト対象であることは明らかである．すると
Derived Categories, Lemma \ref{derived-lemma-compact-objects-subcategory} により，
$\Ker(x)$ は $D(A, \text{d})$ のコンパクト対象である．

\medskip\noindent
次に，$\Ker(x)$ を表す微分付き次数付き（右）$A$-加群を $M$ とし，$M$ は次数付き $A$-加群として
有限かつ射影的であると仮定する．$k[x, y]/(y^2)$ 上の任意の有限次数付き射影加群は次数付き自由なので，
$M$ は次数付き $k[x, y]/(y^2)$-加群として有限自由である（すなわち微分を忘れれば有限自由である）．
$N = M/M(x^2 + x)$ とおく．完全列
$$
0 \to M \xrightarrow{x^2 + x} M \to N \to 0
$$
を考える．$x^2 + x$ は次数 $0$ で $A$ の中心に属し，さらに
$\text{d}(x^2 + x) = 0$ なので，これは微分付き次数付き $A$-加群の短完全列である．
さらに $\text{d}(y) = x^2 + x$ だから，$N$ 上の微分は線形である．写像
$$
H^{-1}(N) \to H^0(M)
\quad\text{and}\quad
H^0(M) \to H^0(N)
$$
は同型である．実際，$D(A, \text{d})$ において $M \cong \Ker(x)$ なので
$H^*(M) = H^0(M) = k$ である．境界写像を計算すると，次数付き加群として
$H^*(N) = k[x, y]/(x, y^2)$ となるが，詳細は省略する．
$N$ は自由 $k[x, y]/(y^2, x^2 + x)$-加群なので，微分と両立する分解
$$
\ldots \to N[2] \xrightarrow{y} N[1] \xrightarrow{y} N \to N/Ny \to 0
$$
をもつ．$N$ は有界であり，$H^0(N) = k[x,y]/(x, y^2)$ なので，Homology, Lemma
\ref{homology-lemma-first-quadrant-ss} から $H^0(N/Ny) = k[x]/(x)$ が従う．
しかし $N/Ny$ は自由 $k[x]/(x^2 + x) = k \times k$-加群からなる有限複体なので，
そのコホモロジーの次元は偶数でなければならず，矛盾する．

\begin{lemma}
\label{examples-lemma-no-good-representatif-compact-object}
微分付き次数付き代数 $(A, \text{d})$ と $D(A, \text{d})$ のコンパクト対象 $E$ で，
$E$ を有限かつ次数付き射影的な微分付き次数付き $A$-加群によって表すことのできないものが存在する．
\end{lemma}

\begin{proof}
上の議論を参照せよ．
\end{proof}
% END EXPANDED MODULE ch110_70.tex
% BEGIN EXPANDED MODULE ch110_71.tex
\section{二つの微分付き次数付き圏}
\label{examples-section-nongraded-differential-graded}

\noindent
この節では，Differential Graded Algebra, Situation \ref{dga-situation-ABC} の公理
(A)，(B)，(C) を満たすが，その対象には $\mathbf{Z}$-次数が付随しない二つの微分付き次数付き圏を構成する．

\medskip\noindent
{\bf 例 I．} $X$ を位相空間とする．
% P12-ERR-0381: frozen missing preposition in the English notation sentence; see ERRATA_CANDIDATES.jsonl.
$\underline{\mathbf{Z}}$ で値 $\mathbf{Z}$ の定数層を表す．$A$ を
$\underline{\mathbf{Z}}$-トーサーとする．この状況で，アーベル群の層 $\mathcal{F}$ が
{\it $A$-次数付き}であるとは，局所切断 $a \in A(U)$ が与えられるごとに射影子
$p_a : \mathcal{F}|_U \to \mathcal{F}|_U$ が存在し，局所同型
$\underline{\mathbf{Z}}|_U \to A|_U$ があるときには必ず
$\mathcal{F}|_U = \bigoplus_{n \in \mathbf{Z}} p_n(\mathcal{F})$ となることをいう．
別の言い方をすれば，$X$ 上局所的にはアーベル層 $\mathcal{F}$ は
$\mathbf{Z}$-次数をもつが，重なり上では，異なる次数付けの選択はトーサー $A$ の遷移関数が与える
次数のシフトだけ異なる．組 $(\mathcal{F}, \text{d})$ が
{\it アーベル層の $A$-次数付き複体}であるとは，$\mathcal{F}$ が $A$-次数付きアーベル層であり，
$\text{d} : \mathcal{F} \to \mathcal{F}$ が微分，すなわち $\text{d}^2 = 0$ を満たし，かつ
$A$ の任意の局所切断 $a$ に対して
$p_{a + 1} \circ \text{d} = \text{d} \circ p_a$ となることをいう．言い換えると，
$\text{d}(p_a(\mathcal{F}))$ は $p_{a + 1}(\mathcal{F})$ に含まれる．

\medskip\noindent
次に，圏 $\mathcal{A}$ を次のように定める．
\begin{enumerate}
\item 対象はアーベル層の $A$-次数付き複体である．
\item 対象 $(\mathcal{F}, \text{d})$，$(\mathcal{G}, \text{d})$ に対して
$$
\Hom_\mathcal{A}((\mathcal{F}, \text{d}), (\mathcal{G}, \text{d})) =
\bigoplus \Hom^n(\mathcal{F}, \mathcal{G})
$$
とおく．ここで $\Hom^n(\mathcal{F}, \mathcal{G})$ は，$A$ のすべての局所切断 $a$ に対して
$f(p_a(\mathcal{F})) \subset p_{a + n}(\mathcal{G})$ を満たすアーベル層の写像 $f$ 全体の群である．
微分として $\text{d}(f) = \text{d} \circ f - (-1)^n f \circ \text{d}$ をとる．
Differential Graded Algebra, Example \ref{dga-example-category-complexes} を参照せよ．
\end{enumerate}
これが実際に (A)，(B)，(C) を満たす微分付き次数付き圏であることの確認は省略する．
すべての性質は $X$ 上局所的に確認でき，そこではアーベル層の複体からなる微分付き次数付き圏が
そのまま得られる．したがって三角圏 $K(\mathcal{A})$ を得る．

\medskip\noindent
$X$ の捻れ導来圏．$\mathcal{A}$ の対象 $(\mathcal{F}, \text{d})$ が与えられると，
よく定義された $A$-次数付きコホモロジー層 $H(\mathcal{F}, \text{d})$ があることに注意する．
したがって $K(\mathcal{A})$ における準同型が何を意味するかは明らかである．準同型を逆にすることにより，
{\it $A$-次数付き層の複体の導来圏 $D(\mathcal{A})$} を得る．$A$ が自明なトーサーなら
$D(\mathcal{A})$ は $D(X)$ に等しいが，
% P12-ERR-0382: frozen additive ``nonzero torsor'' wording; see ERRATA_CANDIDATES.jsonl.
非零トーサーに対しては $X$ の一種の{\it 捻れ}導来圏が得られる．

\medskip\noindent
{\bf 例 II．} $C$ を標数 $2$ の完全体 $k$ 上の滑らかな曲線とする．このとき
$\Omega_{C/k}$ には標準的な平方根が備わる．すなわち，
$\Omega_{C/k} = \mathcal{L}^{\otimes 2}$ と書くことができ，$C$ 上の任意の局所関数 $f$ に対し，
ある $\mathcal{L}$ の局所切断 $s$ が存在して $\text{d}(f)$ は $s^{\otimes 2}$ に等しい．
その「理由」は
$$
\text{d}(a_0 + a_1t + \ldots +a_dt^d) =
(\sum\nolimits_{i\text{ は奇数}} a_i^{1/2} t^{(i - 1)/2})^2\text{d}t
$$
である
% P12-ERR-0383: frozen unresolved editorial placeholder; see ERRATA_CANDIDATES.jsonl.
（ここに将来の参照を挿入）．とりわけ，これにより標準接続
$$
\nabla_{can} :
\Omega_{C/k}
\longrightarrow
\Omega_{C/k} \otimes_{\mathcal{O}_C} \Omega_{C/k}
$$
で，$2$-曲率が零であるもの（すなわち，平方の微分が零となる一意な接続）が定まる．
接続をもつベクトル束の圏はテンソル圏なので，すべての $n \in \mathbf{Z}$ に対して可逆層
$\Omega_{C/k}^{\otimes n}$ 上の標準接続 $\nabla_{can}$ も得られる．

\medskip\noindent
$\mathcal{A}$ を次の圏とする．
\begin{enumerate}
\item 対象は，有限局所自由層 $\mathcal{F}$ と，$2$-曲率が零である接続
$$
\nabla :
\mathcal{F}
\longrightarrow
\mathcal{F} \otimes_{\mathcal{O}_C} \Omega_{C/k}
$$
からなる組 $(\mathcal{F}, \nabla)$ である．
\item $(\mathcal{F}, \nabla_\mathcal{F})$ と $(\mathcal{G}, \nabla_\mathcal{G})$ の間の射を
$$
\Hom_\mathcal{A}((\mathcal{F}, \nabla_\mathcal{F}),
(\mathcal{G}, \nabla_\mathcal{G})) =
\bigoplus \Hom_{\mathcal{O}_C}(\mathcal{F},
\mathcal{G} \otimes_{\mathcal{O}_C} \Omega_{C/k}^{\otimes n})
$$
で与える．次数 $n$ の元
$f : \mathcal{F} \to \mathcal{G} \otimes \Omega_{C/k}^{\otimes n}$ に対し，適切な同一視の下で
$$
\text{d}(f) =
\nabla_{\mathcal{G} \otimes \Omega_{C/k}^{\otimes n}} \circ f +
f \circ \nabla_\mathcal{F}
$$
とおく．
\end{enumerate}
これが性質 (A)，(B)，(C) をもつ微分付き次数付き圏をなすことの確認は省略する．
したがって三角ホモトピー圏 $K(\mathcal{A})$ を得る．

\medskip\noindent
$C = \mathbf{P}^1_k$ ならば $K(\mathcal{A})$ は零圏である．一方，$C$ が種数 $> 1$ の
滑らかな固有曲線ならば $K(\mathcal{A})$ は零でない．実際，$\mathcal{N}$ を次数
$0 \leq d < g - 1$ の可逆層で，零でない切断 $\sigma$ をもつものとする．このとき
$(\mathcal{F}, \nabla_\mathcal{F}) = (\mathcal{O}_C, \text{d})$ および
$(\mathcal{G}, \nabla_\mathcal{G}) = (\mathcal{N}^{\otimes 2}, \nabla_{can})$ とおく．すると
% P12-ERR-0384: frozen three-case display omits n > 1; see ERRATA_CANDIDATES.jsonl.
$$
\Hom_\mathcal{A}^n((\mathcal{F}, \nabla_\mathcal{F}),
(\mathcal{G}, \nabla_\mathcal{G})) =
\left\{
\begin{matrix}
0 & \text{if} & n < 0 \\
\Gamma(C, \mathcal{N}^{\otimes 2}) & \text{if} & n = 0 \\
\Gamma(C, \mathcal{N}^{\otimes 2} \otimes \Omega_{C/k}) & \text{if} & n = 1
\end{matrix}
\right.
$$
% P12-ERR-0385: frozen missing verb in the English explanation; see ERRATA_CANDIDATES.jsonl.
最初の $0$ は，条件 $d < g - 1$ により
$\mathcal{N}^{\otimes 2} \otimes \Omega_{C/k}^{\otimes -1}$ の次数が負だからである．
さて，切断 $\sigma^{\otimes 2}$ の微分は零なので，準同型群
$$
\Hom_{K(\mathcal{A})}((\mathcal{F}, \nabla_\mathcal{F}),
(\mathcal{G}, \nabla_\mathcal{G}))
$$
は零でない．
% END EXPANDED MODULE ch110_71.tex
% BEGIN EXPANDED MODULE ch110_72.tex
\section{固有代数空間のスタックは代数的でない}
\label{examples-section-proper-spaces-not-algebraic}

\noindent
Quot, Section \ref{quot-section-stack-of-spaces} において，群亜群のスタック
$$
p'_{fp, flat, proper} :
\Spacesstack'_{fp, flat, proper}
\longrightarrow
\Sch_{fppf}
$$
を導入して調べた．これは，スキーム $S$ 上の切断の圏が，$S$ 上平坦，固有かつ有限表示な
代数空間の圏であるスタックである．このスタックが Artin の公理の多くを満たすことを証明した．
% P12-ERR-0386: frozen missing verb in the English introduction; see ERRATA_CANDIDATES.jsonl.
この節では，形式的有効性が一般には成り立たないことを示すことにより，このスタックが代数的でないことを示す．

\medskip\noindent
標準的な例では，次の事実を用いる：次元 $g$ のアーベル多様体の普遍変形空間は $g^2$ 個の形式的
パラメータをもつのに対し，任意の有効な形式的変形は次元 $\leq g(g + 1)/2$ の完備局所環上で定義できる．
ここでは，K3 曲面の適当な非有効変形を書き下すことで例を構成する．議論の概略だけを述べ，詳細はすべては与えない．

\medskip\noindent
$k = \mathbf{C}$ を複素数体とする．$X \subset \mathbf{P}^3_k$ を $k$ 上の滑らかな次数
$4$ の曲面とする．このとき $\omega_X \cong \Omega^2_{X/k} \cong \mathcal{O}_X$ である．さらに
$\dim_k H^0(X, T_{X/k}) = 0$，
$\dim_k H^1(X, T_{X/k}) = 20$，および
$\dim_k H^2(X, T_{X/k}) = 0$ である．$X$ は $k$ 上滑らかなので
$L_{X/k} = \Omega_{X/k}$ であり，
$\Ext^i_{\mathcal{O}_X}(\Omega_{X/k}, \mathcal{O}_X) =
H^i(X, T_{X/k})$ であり，さらに Cotangent, Lemma
\ref{cotangent-lemma-find-obstruction-ringed-topoi} があるので，$X$ の普遍変形が
$$
k[[x_1, \ldots, x_{20}]]
$$
上に存在することが分かる．この普遍変形が有効であると仮定する
（Artin's Axioms, Section \ref{artin-section-formal-objects} の意味で）．すると，普遍変形を回収する
代数空間 $Y$ と平坦固有射
$$
f : Y \longrightarrow \Spec(k[[x_1, \ldots, x_{20}]])
$$
が得られるはずである．これは次の理由で不可能である．$Y$ は分離的なので，アフィン開部分スキーム
$V \subset Y$ をとれる．$Y$ の特殊ファイバー $X$ は滑らかなので，$f$ は滑らかである．したがって
$Y$ は正則なもの上滑らかであることから正則であり，$Y$ における $V$ の補集合 $D$ は有効 Cartier 因子である．
このとき $\mathcal{O}_Y(D)$ は $\Pic(Y)$ の非自明な元である
% P12-ERR-0387: frozen incomplete English Picard-class phrase; see ERRATA_CANDIDATES.jsonl.
（これを証明するには，体上の固有滑らかな代数空間の空でないアフィン開集合の補集合が Picard 群の非自明な類を
常に定めることを示し，それを $f$ の一般ファイバーに適用する）．最後に矛盾を得るため，
$\Pic(Y) = 0$ を示す．実際，写像 $\Pic(Y) \to \Pic(X)$ は単射である．これは
$H^1(X, \mathcal{O}_X) = 0$ であること（したがって
% P12-ERR-0388: frozen unbased product in the thickening expression; see ERRATA_CANDIDATES.jsonl.
$\mathcal{O}_X$ の $Y \times \Spec(k[[x_i]]/\mathfrak m^n)$ へのすべての変形は自明である）と，
Grothendieck の存在定理（厚化上で同じ層を生じる連接加群は同型であるという定理）による．
$X$ が十分一般ならば，$\Pic(X) = \mathbf{Z}$ は $\mathcal{O}_X(1)$ によって生成される．
したがって，$\mathcal{O}_X(n)$，$n > 0$ が一次近傍へ変形しないことを示せば十分である
\footnote{この議論は，写像 $c_1 : \Pic(X) \to H^1(X, \Omega_{X/k})$ が単射である限り成り立つ．
標数零の任意の体 $k$ と $\mathbf{P}^3_k$ 内の任意の滑らかな次数 $4$ の超曲面 $X$ に対して，これは真である．}．
カップ積
$$
H^1(X, \Omega_{X/k}) \times H^1(X, T_{X/k})
\longrightarrow H^2(X, \mathcal{O}_X)
$$
を考える．連接双対性により，これは非退化なペアリングである．計算により，Hodge コホモロジーにおける
Chern 類 $c_1(\mathcal{O}_X(n)) \in H^1(X, \Omega_{X/k})$ は零でない．したがって，
$c_1(\mathcal{O}_X(n))$ とのカップ積が零でない一次変形が存在する．最後に，このカップ積が持ち上げの
障害類であることを示す．

\begin{lemma}
\label{examples-lemma-proper-spaces-not-algebraic}
スキーム $S$ 上の切断の圏が $S$ 上平坦，固有かつ有限表示な代数空間の圏である群亜群のスタック
$$
p'_{fp, flat, proper} :
\Spacesstack'_{fp, flat, proper}
\longrightarrow
\Sch_{fppf}
$$
（Quot, Section \ref{quot-section-stack-of-spaces} を参照）は代数スタックでない．
\end{lemma}

\begin{proof}
% P12-ERR-0389: frozen English conditional and agreement errors; see ERRATA_CANDIDATES.jsonl.
これが代数スタックならば，すべての形式的対象は有効となる．Artin's Axioms, Lemma
\ref{artin-lemma-effective} を参照せよ．上の議論により，$\Spec(\mathbf{C})$ への基底変換後には
そうならないことが分かる．したがって結論を得る．
\end{proof}
% END EXPANDED MODULE ch110_72.tex
% BEGIN EXPANDED MODULE ch110_73.tex
\section{代数的でない Hom スタックの例}
\label{examples-section-non-algebraic-hom-stack}

\noindent
$\mathcal{Y}, \mathcal{Z}$ をスキーム $S$ 上の代数スタックとする．
{\it Hom スタック} $\underline{\Mor}_S(\mathcal{Y}, \mathcal{Z})$ とは，スキーム $T$ 上の
切断の圏が
$$
\Mor_T(\mathcal{Y} \times_S T, \mathcal{Z} \times_S T)
$$
で与えられる $S$ 上の群亜群のスタックである．ここで対象は $1$-射，射は $2$-射である．
これが実際に $(\Sch/S)_{fppf}$ 上の群亜群のスタックであることの証明は省略する
% P12-ERR-0390: frozen unresolved editorial placeholder; see ERRATA_CANDIDATES.jsonl.
（ここに将来の参照を挿入）．もちろん，一般には Hom スタックは代数的でない．この節では，
実際に代数的でなく，しかも $\mathcal{Y}$ は $S$ 上固有平坦なスキームによって表現され，
$\mathcal{Z}$ は $S$ 上滑らかかつ固有である例を与える．

\medskip\noindent
$k$ を有限体の代数閉包でない代数閉体とする．$S = \Spec(k[[t]])$ および
% P12-ERR-0391: frozen t^n indexing makes S_0 the zero-ring spectrum; see ERRATA_CANDIDATES.jsonl.
$S_n = \Spec(k[t]/(t^n)) \subset S$ とする．写像 $f : X \to S$ が次を満たすとする．
\begin{enumerate}
\item $f$ は射影的かつ平坦で，そのファイバーは幾何学的に連結な曲線である．
\item ファイバー $X_0 = X \times_S S_0$ は，滑らかな既約成分をもつ節点曲線であり，その双対グラフは
有理曲線からなるループをもつ．
\item $X$ は正則スキームである．
\end{enumerate}
このような曲面 $X$ の例として，$\mathbf{P}^2_{k[[t]]}$ 内の
$$
X\quad :\quad T_0T_1T_2 - t(T_0^3 + T_1^3 + T_2^3) = 0
$$
をとれる．$A_0$ を $k$ 上の非零アーベル多様体，例えば楕円曲線とする．
$A = A_0 \times_{\Spec(k)} S$ を $A_0$ に付随する $S$ 上の定数アーベルスキーム
\footnote{基底 $S$ 上のアーベルスキームとは，$S$ 上平坦，固有かつ有限表示な群スキームで，
すべてのファイバーがアーベル多様体であるものをいう．}とする．
% P12-ERR-0392: frozen defining formula has one unmatched closing parenthesis; see ERRATA_CANDIDATES.jsonl.
スタック $\mathcal{X} = \underline{\Mor}_S(X, [S/A]))$ が代数的でないことを示す．

\medskip\noindent
$[S/A]$ は，一方では $A$ が $S$ に自明に作用する商スタックであり，他方では fppf
$A$-トーサーを分類するスタックに等しい．Examples of Stacks, Proposition
\ref{examples-stacks-proposition-equal-quotient-stacks} を参照せよ．
$[S/A] = [\Spec(k)/A_0] \times_{\Spec(k)} S$ であることに注意する．これにより，スキーム $T$ 上の
ファイバー圏は次のように記述できる：
\begin{align*}
\mathcal{X}_T
& =
\underline{\Mor}_S(X, [S/A])_T \\
& =
\Mor_T(X \times_S T, [S/A] \times_S T) \\
& =
\Mor_S(X \times_S T, [S/A]) \\
& =
\Mor_{\Spec(k)}(X \times_S T, [\Spec(k)/A_0]).
\end{align*}
これは任意の $S$-スキーム $T$ に対して成り立つ．言い換えると，群亜群 $\mathcal{X}_T$ は
$X \times_S T$ 上の fppf $A_0$-トーサーの群亜群である．$\mathcal{X}$ が代数スタックでない理由を
論じる前に，いくつかの補題が必要である．

\begin{lemma}
\label{examples-lemma-torsors-over-two-dimensional-regular}
$W$ を関数体 $K$ をもつ 2 次元正則整 Noether スキームとし，$G \to W$ をアーベルスキームとする．
このとき写像 $H^1_{fppf}(W, G) \to H^1_{fppf}(\Spec(K), G)$ は単射である．
\end{lemma}

\begin{proof}[証明の概略]
$P \to W$ を，一般点で自明となる fppf $G$-トーサーとする．このとき $W$ 上の射
$\Spec(K) \to P$ があり，そのスキーム論的像 $Z \subset P$ をとれる．$P \to W$ は固有なので
（$W$ 上固有な群代数空間のトーサーだからである），$Z \to W$ は固有双有理射である．
Spaces over Fields, Lemma \ref{spaces-over-fields-lemma-finite-in-codim-1} により，射
$Z \to W$ は $W$ の有限個の閉点を除いて有限である．
% P12-ERR-0393: frozen unresolved surface-resolution placeholder; see ERRATA_CANDIDATES.jsonl.
（曲面に対し二次変換を吹き上げることで射の不定性を解消することについて，ここに将来の参照を挿入）により，
$Z \to W$ の幾何学的ファイバーの既約成分は有理曲線である．More on Groupoids in Spaces, Lemma
\ref{spaces-more-groupoids-lemma-no-nonconstant-morphism-from-P1-to-group} により，有理曲線から群スキーム，
またはそのトーサーへの非定数射は存在しない．したがって $Z \to W$ は有限であり，ゆえに $Z$ はスキームで，
Morphisms, Lemma \ref{morphisms-lemma-finite-birational-over-normal} により $Z \to W$ は同型である．
言い換えると，トーサー $P$ は自明である．
\end{proof}

\begin{remark}
\label{examples-remark-torsors-over-regular}
Lemma \ref{examples-lemma-torsors-over-two-dimensional-regular} における $W$ の $2$ 次元という条件は除ける．
実際，$W$ が Noether，正則かつ整ならば，(1) 任意の $G$-トーサー $P$ に対して
$P(W)\to P(K)$ は全単射であり，(2) $H^1(W,G)\to H^1(K,G)$ は単射である．
証明の概略は次のとおりである．まず (2) は明らかに (1) から従う．(1) について，単射性は自明であり，
全射性は（現在の証明と同様に）有理切断のスキーム論的閉包が切断であることを示す問題である．
この問題は $W$ 上 fppf 局所的なので，\'etale 基底変換の後，$P=G$ と仮定してよい
（\'etale 被覆は正則性を保つことに注意する）．このとき
$G=\underline{\mathrm{Pic}}^0_{G'/W}$ であり，$G'$ は双対アーベルスキームである．$G'$ は正則スキームなので，
一般的に定義された可逆層が延長することから結論が従う．実際，「$W$ は正則」を
「滑らかな $W$-スキームは局所階乗的である」まで弱められる．Laurent Moret-Bailly は Raynaud から
Picard の技巧を学んだ；Raynaud の仕事のどこかに参照があるはずである．
\end{remark}

\begin{lemma}
\label{examples-lemma-torsors-over-field-torsion}
$G$ を体 $K$ 上の滑らかな可換群代数空間とする．このとき
$H^1_{fppf}(\Spec(K), G)$ は捩れ群である．
\end{lemma}

\begin{proof}
$\Spec(K)$ 上の任意の $G$-トーサー $P$ は，$G$ の一つの形式として $K$ 上滑らかである．
したがって $P$ は有限分離拡大 $L/K$ 上に点をもつ．$[L : K] = n$ とする．
$[n](P)$ で，その類が $H^1_{fppf}(\Spec(K), G)$ における $P$ の類の $n$ 倍である
$G$-トーサーを表す．$K$ 上の代数空間の標準射
$$
P \times_{\Spec(K)} \ldots \times_{\Spec(K)} P \to [n](P)
$$
がある．$G$ はアーベルなので，この射は対称である．したがって商
$$
(P \times_{\Spec(K)} \ldots \times_{\Spec(K)} P)/S_n
$$
を経由する．一方，射 $\Spec(L) \to P$ は射
$$
(\Spec(L) \times_{\Spec(K)} \ldots \times_{\Spec(K)} \Spec(L))/S_n
\longrightarrow (P \times_{\Spec(K)} \ldots \times_{\Spec(K)} P)/S_n
$$
を定め，左辺のスキームが $K$-有理点をもつことは読者が確認できる．したがって
$[n](P)$ は自明なトーサーである．
\end{proof}

\noindent
$\mathcal{X} = \underline{\Mor}_S(X, [S/A])$ が代数スタックでないことを証明するには，
Artin's Axioms, Lemma \ref{artin-lemma-effective} により，次を示せば十分である．

\begin{lemma}
\label{examples-lemma-not-essentially-surjective}
標準写像 $\mathcal{X}(S) \to \lim \mathcal{X}(S_n)$ は本質的全射でない．
\end{lemma}

\begin{proof}[証明の概略]
定義を展開すると，$H^1(X, A_0) \to \lim H^1(X_n, A_0)$ が全射でないことを確かめれば十分である．
$X$ は正則かつ射影的なので，Lemmas \ref{examples-lemma-torsors-over-field-torsion} および
\ref{examples-lemma-torsors-over-two-dimensional-regular} により，$X$ 上の各 $A_0$-トーサーは捩れである．
特に，群 $H^1(X, A_0)$ は捩れ群である．したがって，次を示せば十分である：
(a) 群 $H^1(X_0, A_0)$ は非捩れである；
(b) すべての $n$ に対して写像 $H^1(X_{n + 1}, A_0) \to H^1(X_n, A_0)$ は全射である．

\medskip\noindent
(a) について．節点曲線 $X_0$ の各成分上の自明な $A_0$-トーサーを，節点上での同型として
$A_0$ の非捩れ点を用いて貼り合わせることにより，$X_0$ 上の非捩れ $A_0$-トーサー $P_0$ を構成する．
より正確には，$x \in X_0$ を，有理曲線からなるループに現れる節点とする．
$X'_0 \to X_0$ を $X_0 \setminus \{x\}$ における $X_0$ の正規化とする．
% P12-ERR-0394: frozen x_0 target though the node is named x; see ERRATA_CANDIDATES.jsonl.
$x', x'' \in X'_0$ を $x_0$ に写る二点とする．次に $A_0 \times_{\Spec(k)} X'_0$ をとり，
非捩れ点 $a_0 \in A_0(k)$ による平行移動 $A_0 \to A_0$ を用いて
$A_0 \times {x'}$ と $A_0 \times \{x''\}$ を同一視する
（$k$ は代数閉体で有限体の代数閉包でないので，このような非捩れ点が存在する；これは実際，自明な事実ではない）．
貼り合わせは代数空間（実際にはスキーム）であり，$X_0$ 上の非捩れ $A_0$-トーサーであることを示せる．
% P12-ERR-0395: frozen wrong article in the English torsor phrase; see ERRATA_CANDIDATES.jsonl.
これが非捩れである理由は，$[n](P_0)$ が切断をもてば，その切断は射 $s : X'_0 \to A_0$ で，
$A_0(k)$ の群法において $[n](a_0) = s(x') - s(x'')$ を満たすものを生じるからである．しかし，
% P12-ERR-0396: frozen malformed English subject in the constancy sentence; see ERRATA_CANDIDATES.jsonl.
ループの既約成分は有理的なので，切断 $s$ はそれらの上で定数である
（More on Groupoids in Spaces, Lemma
\ref{spaces-more-groupoids-lemma-no-nonconstant-morphism-from-P1-to-group}）．
したがって $s(x') = s(x'')$ となり，矛盾を得る．

\medskip\noindent
(b) について．変形理論によれば，$A_0$-トーサー $P_n \to X_n$ を $A_0$-トーサー
$P_{n + 1} \to X_{n + 1}$ へ変形する障害は，$X_0$ 上の適当なベクトル束 $\omega$ に対する
$H^2(X_0, \omega)$ に属する．後者は $X_0$ が曲線なので消え，主張が従う．
\end{proof}

\begin{proposition}
\label{examples-proposition-nonalghomstack}
スタック $\mathcal{X} = \underline{\Mor}_S(X, [S/A])$ は代数的でない．
\end{proposition}

\begin{proof}
上の議論を参照せよ．
\end{proof}

\begin{remark}
\label{examples-remark-contradict-aoki}
Proposition \ref{examples-proposition-nonalghomstack} は \cite[Theorem 1.1]{AokiHomStacks} と矛盾する．
問題は $\underline{\Mor}_S(X, [S/A])$ の形式的対象の非有効性である．同じ問題は Erratum
\cite{AokiHomStacksErr} でも言及されており，これは \cite{AokiHomStacks} に対するものである．
残念ながら，その Erratum は
$\mathcal{Z}$ が分離的なら $\underline{\Mor}_S(\mathcal{Y}, \mathcal{Z})$ は代数的であるとも主張しているが，
$[S/A]$ は分離的なので，これも Proposition \ref{examples-proposition-nonalghomstack} と矛盾する．
\end{remark}
% END EXPANDED MODULE ch110_73.tex
% BEGIN EXPANDED MODULE ch110_74.tex
\section{強い形式的有効性を満たさない代数スタック}
\label{examples-section-non-formal-effectiveness}

\noindent
これは \cite[Example 4.12]{Bhatt-Algebraize} である．$k$ を代数閉体とし，$A$ を $k$ 上の
アーベル多様体とする．$A(k)$ が捩れ群でないと仮定する（$k$ が有限体の代数閉包でなければ常に成り立つ）．
$\mathcal{X} = [\Spec(k)/A]$ とおく．あるイデアル $I \subset k[x, y]$ が存在して，
$$
\mathcal{X}_{\Spec(k[x, y]^\wedge)}
\longrightarrow
\lim \mathcal{X}_{\Spec(k[x, y]/I^n)}
$$
が本質的全射でないと主張する．実際，$I$ を $xy(x + y - 1)$ で生成されるイデアルとする．
このとき $X_0 = V(I)$ は，三点で三角形状に貼り合わされた三つの $\mathbf{A}^1_k$ からなる．
したがって，$X_0$ の既約成分上では自明なトーサーをとり，非捩れ点による平行移動で貼り合わせることにより，
$X_0$ 上の $A$ に対する無限位数のトーサー $P_0$ を作れる．Lemma
\ref{examples-lemma-not-essentially-surjective} の証明とまったく同様に，$P_0$ を
$X_n = \Spec(k[x, y]/I^n)$ 上のトーサー $P_n$ へ持ち上げられる．
$k[x, y]^\wedge$ は 2 次元正則整域なので，$\Spec(k[x, y]^\wedge)$ 上の $A$ に対する任意のトーサー
$P$ は捩れである（Lemmas \ref{examples-lemma-torsors-over-two-dimensional-regular} および
\ref{examples-lemma-torsors-over-field-torsion}）．したがって，このトーサーの系は表示された関手の像に属さない．

\begin{lemma}
\label{examples-lemma-non-formal-effectiveness}
% P12-ERR-0397: frozen ``closure'' rather than ``algebraic closure''; see ERRATA_CANDIDATES.jsonl.
$k$ を有限体の閉包でない代数閉体とし，$A$ を $k$ 上のアーベル多様体とする．
$\mathcal{X} = [\Spec(k)/A]$ とおく．局所冪零な核をもつ全射遷移写像を備えた
$k$-代数の逆系 $R_n$ と，
% P12-ERR-0398: frozen sentence omits ``objects of'' the stack; see ERRATA_CANDIDATES.jsonl.
系 $(\Spec(R_n))$ 上にある $\mathcal{X}$ の系 $(\xi_n)$ で，Artin's Axioms, Remark
\ref{artin-remark-strong-effectiveness} の意味で有効でないものが存在する．
\end{lemma}

\begin{proof}
上の議論を参照せよ．
\end{proof}
% END EXPANDED MODULE ch110_74.tex
% BEGIN EXPANDED MODULE ch110_75.tex
\section{Grothendieck の存在定理に対する反例}
\label{examples-section-Grothendieck-existence}

\noindent
$k$ を体，$A = k[[t]]$ とする．$X$ を $U = \Spec(A[x])$ と $V = \Spec(A[y])$ を，同一視
$$
U \setminus \{0_U\} \longrightarrow V \setminus \{0_V\}
$$
によって貼り合わせたものとする．この同一視は $x$ を $y$ に写す．ここで
% P12-ERR-0399: frozen capital O_V token; see ERRATA_CANDIDATES.jsonl.
$0_U \in U$ と $O_V \in V$ はそれぞれ極大イデアル $(x, t)$ と $(y, t)$ に対応する点である．
$A_n = A/(t^n)$ および $X_n = X \times_{\Spec(A)} \Spec(A_n)$ とおく．
$\mathcal{F}_n$ を，$A_n[x]$-加群 $A_n[x]/(x) \cong A_n$ と $A_n[y]$-加群 $0$ を明らかな方法で
貼り合わせることに対応する $X_n$ 上の連接層とする．
% P12-ERR-0400: frozen bare verb in the English ideal-sheaf definition; see ERRATA_CANDIDATES.jsonl.
$\mathcal{I} \subset \mathcal{O}_X$ を $t$ で生成されるイデアル層とする．すると
$(\mathcal{F}_n)$ は，Cohomology of Schemes, Section
\ref{coherent-section-existence-proper-support} で定義される圏
$\textit{Coh}_{\text{support proper over } A}(X, \mathcal{I})$ の対象である．一方，この対象は
Cohomology of Schemes, Equation
(\ref{coherent-equation-completion-functor-proper-over-A}) の関手の像に属さない．
% P12-ERR-0401: frozen ``if it where'' conditional; see ERRATA_CANDIDATES.jsonl.
実際，像に属するとすれば，有限 $A[x]$-加群 $M$，有限 $A[y]$-加群 $N$，および同型
$M[1/t] \cong N[1/t]$ で，すべての $n$ に対して
$M/t^nM \cong A_n[x]/(x)$ かつ $N/t^nN = 0$ を満たすものが存在するはずである．
これが不可能であることは容易に分かる．

\begin{lemma}
\label{examples-lemma-counter-Grothendieck-existence}
連接層の代数化に対する反例として，次が成り立つ．
\begin{enumerate}
\item Cohomology of Schemes, Theorem \ref{coherent-theorem-grothendieck-existence} の
Grothendieck 存在定理は，$X \to \Spec(A)$ が分離的であるという仮定を除けば偽である．
\item Quot, Theorems \ref{quot-theorem-coherent-algebraic-general} および
\ref{quot-theorem-coherent-algebraic} の連接層のスタック $\Cohstack_{X/B}$ は，
$X \to S$ が分離的であるという仮定を除けば一般には代数的でない．
\item Quot, Proposition \ref{quot-proposition-quot} の関手
$\Quotfunctor_{\mathcal{F}/X/B}$ は，$X \to B$ が分離的であるという仮定を除けば
一般には代数空間でない．
\end{enumerate}
\end{lemma}

\begin{proof}
(1) は上で見た．これは $\textit{Coh}_{X/A}$ が Artin's Axioms, Section
\ref{artin-section-axioms} の公理 [4] を満たさないことを示す．したがって Artin's Axioms, Lemma
\ref{artin-lemma-effective} により，これは代数スタックではありえない．このようにして (2) が分かる．
(3) を見るには，すべての $n$ に対して両立する全射
$\mathcal{O}_{X_n} \to \mathcal{F}_n$ があることに注意する．したがって
$\Quotfunctor_{\mathcal{O}_X/X/A}$ は公理 [4] を満たさず，前と同様に (3) が分かる．
\end{proof}
% END EXPANDED MODULE ch110_75.tex
% BEGIN EXPANDED MODULE ch110_76.tex
\section{アフィン形式代数空間}
\label{examples-section-affine-formal-algebraic-space}

\noindent
$K$ を体とし，$(V_i)_{i \in I}$ を，遷移写像が全射で $\lim V_i = 0$ を満たす，$K$ 上の
非零ベクトル空間の有向逆系とする．Section \ref{examples-section-zero-limit} を参照せよ．
$V_i$ を平方零イデアルとして $R_i = K \oplus V_i$ を $K$-代数とみなす．すると $R_i$ は，
遷移写像が全射で核が冪零，かつ $\lim R_i = K$ を満たす $K$-代数の逆系である．
アフィン形式代数空間 $X = \colim \Spec(R_i)$ は，McQuillan でないアフィン形式代数空間の例である．

\begin{lemma}
\label{examples-lemma-affine-not-mcquillan}
McQuillan でないアフィン形式代数空間が存在する．
\end{lemma}

\begin{proof}
上の議論を参照せよ．
\end{proof}

\noindent
Section \ref{examples-section-zero-limit} のような完全列の系
$0 \to W_i \to V_i \to K \to 0$ をとる．
$A_i = K[V_i]/(ww'; w, w' \in W_i)$ とおく．すると，冪零な核をもつ両立する全射の系
$A_i \to K[t]$ があり，遷移写像 $A_i \to A_j$ も全射で冪零な核をもつ．
$V_i$ は $s \in S_i$ を基底とする自由 $K$-加群であることを思い出そう．すると $K$ の標数が零なら，
$A_i$ の次数 $d$ 部分は，それぞれ $t^d$ に写る $s^d$，$s \in S_i$ を基底とする自由 $K$-加群である．
したがって $A_i$ の次数 $d$ 部分の逆系はベクトル空間 $V_i$ の逆系と同型である．
$\lim V_i = 0$ なので，少なくとも $K$ の標数が零なら $\lim A_i = K$ と結論できる．
これは，「正則関数」が点を分離しないアフィン形式代数空間の例を与える．

\begin{lemma}
\label{examples-lemma-affine-formal-functions-do-not-separate-points}
正則関数が次の意味で点を分離しないアフィン形式代数空間 $X$ が存在する：
Formal Spaces, Definition \ref{formal-spaces-definition-affine-formal-algebraic-space} のように
$X = \colim X_\lambda$ と書けば，$\lim \Gamma(X_\lambda, \mathcal{O}_{X_\lambda})$ は体であるが，
$X_{red}$ は無限個の点をもつ．
\end{lemma}

\begin{proof}
上の議論を参照せよ．
\end{proof}

\noindent
$K$，$I$，$(V_i)$ を上のとおりとする．次の系を考える：
$$
\Phi = (\Lambda, J_i \subset \Lambda, (M_i) \to (V_i))
$$
ここで $\Lambda$ は増大写像をもつ $K$-代数，$J_i \subset \Lambda$ は各 $i \in I$ に対する
平方零イデアル，$(M_i) \to (V_i)$ は $K$-ベクトル空間の逆系の写像で，次を満たすものとする：
各 $i$ に対して $M_i \to V_i$ は全射である；$M_i$ は $\Lambda$-加群構造をもつ；遷移写像
$M_i \to M_j$，$i > j$ は $\Lambda$-線形である；$i > j$ に対して
$J_j M_i \subset \Ker(M_i \to M_j)$ である．主張：すべての $i > j$ に対して
$M_j = M_i/J_j M_i$ となる上のような系が存在する．

\medskip\noindent
この主張が正しければ，アフィン形式代数空間の表現可能な射
$$
\colim_{i \in I} \Spec(\Lambda/J_i \oplus M_i)
\longrightarrow
\text{Spf}(\lim \Lambda/J_i)
$$
で，始域は McQuillan でないが終域は McQuillan であるものが得られる．ここで
$\Lambda/J_i \oplus M_i$ は，$M_i$ を平方零イデアルとする通常の $\Lambda/J_i$-代数構造をもつ．
表現可能性は，$i > j$ に対する条件 $M_i/J_jM_i = M_j$ にちょうど対応する．
% P12-ERR-0402: frozen ungrammatical nonsurjectivity predicate; see ERRATA_CANDIDATES.jsonl.
射の始域が McQuillan でないのは，射影 $\lim_{i \in I} M_i \to M_i$ が全射でないからである．
これは，写像 $\lim V_i \to V_i$ が全射でなく，かつ全射 $M_i \to V_i$ があるためである．
いくつかの詳細は省略する．

\medskip\noindent
主張の証明．まず，少なくとも一つの系，すなわち
$$
\Phi_0 = (K, J_i = (0), (V_i) \xrightarrow{\text{id}} (V_i))
$$
が存在することに注意する．系 $\Phi$ が与えられたとき，系の射
$\Phi \to \Phi'$（系の射は明らかな仕方で定義する）で，
$\Ker(M_i/J_j M_i \to M_j)$ が $M'_i/J'_j M'_i$ の中で零に写るものが存在することを示す．
これができれば，$n \geq 1$ に対して帰納的に $\Phi_n = (\Phi_{n - 1})'$ とおき，
$\Phi = \colim \Phi_n$ をとる通常の技巧により，望ましい性質をもつ系が得られる．詳細は省略する．

\medskip\noindent
$\Phi$ からの $\Phi'$ の構成．$i > j$ かつ $\xi \in \Ker(M_i \to M_j)$ である三つ組
$u = (i, j, \xi)$ 全体の集合を $U$ とする．$s, t : U \to I$ を写像
$s(i, j, \xi) = i$ および $t(i, j, \xi) = j$ とする．そして $\xi_u \in M_{s(u)}$ を $u$ の
第三成分とする．
$$
\Lambda' = \Lambda[x_u; u \in U]/(x_u x_{u'}; u, u' \in U)
$$
をとり，増大写像 $\Lambda' \to K$ は $\Lambda$ の増大写像と $x_u$ を零に送ることによって定める．
% P12-ERR-0403: frozen malformed generator subscript; see ERRATA_CANDIDATES.jsonl.
$J'_k = J_k \Lambda' + (x_{u,\ t(u) \geq k})$ とおく．さらに
$$
M'_i = M_i \oplus \bigoplus\nolimits_{s(u) \geq i} K\epsilon_{i, u}
$$
とおく．$i > j$ に対する遷移写像 $M'_i \to M'_j$ として，所与の写像
$M_i \to M_j$ を用い，$\epsilon_{i, u}$ を $\epsilon_{j, u}$ に送る．写像
$M'_i \to V_i$ は所与の写像 $M_i \to V_i$ を誘導し，$\epsilon_{i, u}$ を零に送るものとする．
最後に，$\Lambda'$ を次のように $M'_i$ に作用させる：$\lambda \in \Lambda$ に対し，$M_i$ 上では
$\Lambda$-加群構造によって作用させ，$\epsilon_{i, u}$ 上では増大写像 $\Lambda \to K$ を通じて作用させる．
すべての $i$ に対し，元 $x_u$ は $M_i$ 上 $0$ として作用する．最後に
$$
x_u \epsilon_{i, u} = \text{image of }\xi_u\text{ in }M_i
$$
と定め，他のすべての積 $x_{u'} \epsilon_{i, u}$ を零とおく．表示式は
$s(u) \geq i$ かつ $\xi_u \in M_{s(u)}$ なので意味をもつ．
% P12-ERR-0404: frozen missing infinitive marker in the English checks sentence; see ERRATA_CANDIDATES.jsonl.
確認すべき主要な点は，$i > j$ に対して $J'_j M'_i \subset M'_i$ が $M'_j$ において零に写ること，および
% P12-ERR-0405: frozen unprimed J_j in the target quotient; see ERRATA_CANDIDATES.jsonl.
$\Ker(M_i \to M_j)$ が $M'_i/J_j M'_i$ において零に写ることである．後者の理由は，任意の
$\xi \in \Ker(M_i \to M_j)$ に対して
$\xi = x_{(i, j, \xi)} \epsilon_{i, (i, j, \xi)} \in J'_j M'_i$ だからである．詳細は省略する．

\begin{lemma}
\label{examples-lemma-representable-morphism-affine-formal-not-mcquillan-top}
アフィン形式代数空間の表現可能な射 $f : X \to Y$ で，$Y$ は McQuillan だが $X$ は
McQuillan でないものが存在する．
\end{lemma}

\begin{proof}
上の議論を参照せよ．
\end{proof}
% END EXPANDED MODULE ch110_76.tex
% BEGIN EXPANDED MODULE ch110_77.tex
\section{平坦写像は有限表示平坦写像の有向極限でない}
\label{examples-section-flat-not-colimit-flat-finitely-presented}

\noindent
この節の目的は，平坦かつ有限表示な環準同型のフィルター付き余極限でない平坦環準同型の例を与えることである．
\cite{gabber-nonexcellent} では，$A$ が次元 $1$，剰余標数零の非優秀局所環ならば，その完備化への
（平坦な）環準同型 $A \to A^\wedge$ は有限型平坦環準同型のフィルター付き余極限でないことが示されている．
この節の例では始域が優秀環となる．読者にも別の例を投稿することを勧める；例があれば
\href{mailto:stacks.project@gmail.com}{stacks.project@gmail.com} へ電子メールを送られたい．

\medskip\noindent
構成のため，素数 $p$ を固定し，$A = \mathbf{F}_p[x_1, \ldots, x_n]$ とする．
$A$ の絶対整閉包 $A^+$，すなわち $A^+$ はその分数体の代数閉包における $A$ の整閉包であるものを選ぶ．
\cite[\S 6.7]{HHBigCM} では $A \to A^+$ が平坦であることが示されている．

\medskip\noindent
$n \geq 3$ ならば，$A$-代数 $A^+$ は有限表示平坦 $A$-代数のフィルター付き余極限でないと主張する．

\medskip\noindent
$n = 3$ の場合の議論を概説し，より大きい $n$ への一般化は読者に委ねる．原点における $A$ の狭義
Hensel 化を $R$，その絶対整閉包を $R^+$ とすると，写像 $R \to R^+$ に対する類似の主張を示せば十分である．
$R$ は Hensel 正則局所環であり，その剰余体 $k$ は $\mathbf{F}_p$ の代数閉包であることに注意する．

\medskip\noindent
$k$ 上の ordinary アーベル曲面 $X$ と，$X$ 上の非常に豊富な直線束 $L$ を選ぶ．切断環
$\Gamma_*(X, L) = \bigoplus_n H^0(X,L^n)$ は，$L$ に関する $X$ 上のアフィン錐の座標環である．
$L$ が十分正ならば，これは正規環である．
% P12-ERR-0406: frozen missing article in the English cone-vertex phrase; see ERRATA_CANDIDATES.jsonl.
$S$ を錐の頂点における $\Gamma_*(X, L)$ の Hensel 化とする．すると $S$ は次元 $3$ の Hensel
Noether 正規整域である．有限型 $k$-代数 $\Gamma_*(X, L)$ に対する Noether 正規化を Hensel 化することで，
有限単射 $R \to S$ を得る．$R^+$ は $R$ の絶対整閉包なので，埋め込み $S \to R^+$ も固定できる．
したがって $R^+$ は $S$ の絶対整閉包でもある．$R^+$ が平坦 $R$-代数のフィルター付き余極限でないことを
示すには，次を示せば十分である：
\begin{enumerate}
\item $S \to P$ が環準同型で，$P$ が $R$ 上平坦かつ有限型ならば，環準同型 $S \to T$ で
$T$ が $R$ 上有限平坦であるものが存在する．
\item 任意の有限環準同型 $S \to T$ に対し，環 $T$ は $R$-平坦でない．
\end{enumerate}
実際，$S$ は $R$ 上有限表示なので，$R^+ = \colim_i P_i$ を有限表示平坦 $R$-代数 $P_i$ の
フィルター付き余極限として書けるならば，$i \gg 0$ に対して $S \to R^+$ は
$S \to P_i \to R^+$ と分解し，上の二主張に矛盾する．(1) は $R$ が Hensel であることと切断の議論から従う．
More on Morphisms, Lemma \ref{more-morphisms-lemma-qf-fp-flat-neighbourhood-dominates-fppf} を参照せよ．
(2) は \cite{BhattSmallCMMod} で証明されているが，読者の便宜のため議論を思い出そう．

\medskip\noindent
$U \subset \Spec(S)$ を穴あきスペクトルとする．すると自然な写像
$X \leftarrow U \subset \Spec(S)$ がある．第一の写像は同一視
$H^1(U, \mathcal{O}_U) \simeq H^1(X, \mathcal{O}_X)$ を与える．完全化の Witt ベクトルを経由し，
Artin--Schreier 列を用いることにより
\footnote{ここで，$S$ が標数 $p$ の狭義 Hensel 局所環であり，したがって
$S \to S$，$f \mapsto f^p - f$ が全射であることを用いる．また $S$ は正規整域なので
$\Gamma(U, \mathcal{O}_U) = S$ である．したがって $H^1_\etale(U, \mathbf{Z}/p)$ は，
$f \mapsto f^p - f$ が誘導する写像
$H^1(U, \mathcal{O}_U) \to H^1(U, \mathcal{O}_U)$ の核である．}，同一視
$H^1_\etale(U, \mathbf{Z}_p) \simeq H^1_\etale(X, \mathbf{Z}_p)$ を得る．特に，この群は
階数 $2$ の有限自由 $\mathbf{Z}_p$-加群である（$X$ は ordinary だからである）．矛盾を得るため，
上の (2) のような $R$-平坦な $T$ が存在すると仮定する．$V \subset \Spec(T)$ を $U$ の逆像とし，
誘導される有限全射を $f : V \to U$ と書く．$U$ は正規なので，$U_\etale$ 上のトレース写像
$f_*\mathbf{Z}_p \to \mathbf{Z}_p$ があり，引き戻し
$\mathbf{Z}_p \to f_*\mathbf{Z}_p$ との合成は $d = \deg(f)$ 倍である．コホモロジーをとり，
$H^1_\etale(U, \mathbf{Z}_p)$ が非捩れであることを用いると，
$H^1_\etale(V, \mathbf{Z}_p)$ が零でないことが分かる．$R^1\lim$ の干渉がないため
$H^1_\etale(V, \mathbf{Z}_p) \simeq \lim H^1_\etale(V, \mathbf{Z}/p^n)$ なので，群
$H^1(V_\etale,\mathbf{Z}/p)$ は零でない．$T$ は $R$-平坦なので
$\Gamma(V, \mathcal{O}_V) = T$ は狭義 Hensel であり，Artin--Schreier 列から
$H^1(V, \mathcal{O}_V) \neq 0$ が分かる．これは $H^2_\mathfrak m(T) \neq 0$ と同値である．ここで
$\mathfrak m \subset R$ は極大イデアルである．一方，$T$ は $R$-加群として有限平坦（すなわち有限自由）であり，
$H^2_\mathfrak m(R) = 0$ なので矛盾である．この矛盾が (2) を証明する．

\begin{lemma}
\label{examples-lemma-weird-flat-map}
可換環 $A$ と平坦 $A$-代数 $B$ で，有限表示平坦 $A$-代数のフィルター付き余極限として書けないものが存在する．
実際，$A$ として，有限型 $\mathbf{F}_p$-代数と，剰余体の標数が $0$ である $1$ 次元 Noether 局所環の
いずれかを選べる．
\end{lemma}

\begin{proof}
上の議論を参照せよ．
\end{proof}
% END EXPANDED MODULE ch110_77.tex
% BEGIN EXPANDED MODULE ch110_78.tex
\section{捩れ加群を法とする加群の圏}
\label{examples-section-serre-quotient-modulo-torsion-modules}

\noindent
捩れ群の圏は，すべてのアーベル群の圏の Serre 部分圏である
（Homology, Definition \ref{homology-definition-serre-subcategory}）．より一般に，任意の環 $A$ に対し，
捩れ $A$-加群の圏はすべての $A$-加群の圏の Serre 部分圏である．More on Algebra, Section
\ref{more-algebra-section-abelian-categories-modules} を参照せよ．$A$ が整域なら，商圏
（Homology, Lemma \ref{homology-lemma-serre-subcategory-is-kernel}）は分数体上のベクトル空間の圏と同値である．
これは次のより一般的な命題から従う．

\begin{proposition}
\label{examples-proposition-localization-and-serre-quotients}
$A$ を環，$S$ を $A$ の乗法的部分集合とする．$\text{Mod}_A$ で $A$-加群の圏を表し，
$\mathcal{T}$ を，任意の元が $S$ のある元によって零化される加群からなるその Serre 部分圏とする．
このとき標準同値 $\text{Mod}_A/\mathcal{T} \rightarrow \text{Mod}_{S^{-1}A}$ が存在する．
\end{proposition}

\begin{proof}
$M \mapsto M \otimes_A S^{-1}A$ で与えられる関手
$\text{Mod}_A \to \text{Mod}_{S^{-1}A}$ は完全であり（Algebra, Proposition
\ref{algebra-proposition-localization-exact}），$\mathcal{T}$ の加群を零に写す．したがって Homology, Lemma
\ref{homology-lemma-serre-subcategory-is-kernel} の普遍性により，この関手は関手
$\text{Mod}_A/\mathcal{T} \to \text{Mod}_{S^{-1}A}$ に降下する．

\medskip\noindent
逆に，$M \otimes_A S^{-1}A = 0$ を満たす任意の $A$-加群 $M$ は $\mathcal{T}$ の対象である．実際，
$M \otimes_A S^{-1}A \cong S^{-1} M$ である
（Algebra, Lemma \ref{algebra-lemma-tensor-localization}）．したがって Homology, Lemma
\ref{homology-lemma-quotient-by-kernel-exact-functor} により，関手
$\text{Mod}_A/\mathcal{T} \to \text{Mod}_{S^{-1}A}$ は忠実である．

\medskip\noindent
さらに，この埋め込みは本質的全射である：$S^{-1}A$-加群 $N$ の逆像として $N_A$，すなわち
$N$ を $A$-加群とみなしたものをとれる．実際，$x \otimes a/s$ を $(a/s) \cdot x$ に写す標準写像
$N_A \otimes_A S^{-1}A \to N$ は $S^{-1}A$-加群の同型である．
\end{proof}

\begin{proposition}
\label{examples-proposition-quotient-by-torsion-modules}
$A$ を環とし，$Q(A)$ でその全商環を表す
（Algebra, Example \ref{algebra-example-localize-at-prime} の意味で）．$\text{Mod}_A$ で $A$-加群の圏を，
$\mathcal{T}$ でその捩れ加群からなる Serre 部分圏を表し，$\text{Mod}_{Q(A)}$ で
$Q(A)$-加群の圏を表す．このとき標準同値
$\text{Mod}_A/\mathcal{T} \rightarrow \text{Mod}_{Q(A)}$ が存在する．
\end{proposition}

\begin{proof}
Proposition \ref{examples-proposition-localization-and-serre-quotients} を乗法的部分集合
$S = \{f \in A \mid f \text{ が零因子でない環 }A\}$ に適用すれば直ちに従う．
実際，加群が捩れ加群であることと，その各元が $S$ のある元によって零化されることは同値である．
\end{proof}

\begin{proposition}
\label{examples-proposition-quotient-by-finitely-generated-torsion-modules}
$A$ を Noether 整域とし，$K$ をその分数体とする．$\text{Mod}_A^{fg}$ で有限生成 $A$-加群の圏を表し，
$\mathcal{T}^{fg}$ をその有限生成捩れ加群からなる Serre 部分圏とする．このとき
$\text{Mod}_A^{fg}/\mathcal{T}^{fg}$ は有限次元 $K$-ベクトル空間の圏と標準的に同値である．
\end{proposition}

\begin{proof}
Proposition \ref{examples-proposition-quotient-by-torsion-modules} の同値を考える．
\par\noindent
これは埋め込み $\text{Mod}_A^{fg}/\mathcal{T}^{fg} \to \text{Mod}_A/\mathcal{T}$ に沿って，同値
$\text{Mod}_A^{fg}/\mathcal{T}^{fg} \to \text{Vect}_K^{fd}$ に制限される．Noether 性の仮定により
$\text{Mod}_A^{fg}$ はアーベル圏であり（More on Algebra, Section
\ref{more-algebra-section-abelian-categories-modules} を参照），また標準関手
$\text{Mod}_A^{fg}/\mathcal{T}^{fg} \to \text{Mod}_A/\mathcal{T}$ は充満である
（そうでなければ，有限生成加群の捩れ部分加群が $\mathcal{T}^{fg}$ の対象でないことがありうる）．
\end{proof}

\begin{proposition}
\label{examples-proposition-quotient-abelian-groups-by-torsion-groups}
アーベル群の圏をその捩れ群からなる Serre 部分圏で割った商は，
$\mathbf{Q}$-ベクトル空間の圏である．
\end{proposition}

\begin{proof}
主張は Proposition \ref{examples-proposition-quotient-by-torsion-modules} から直ちに従う．
\end{proof}
% END EXPANDED MODULE ch110_78.tex
% BEGIN EXPANDED MODULE ch110_79.tex
\section{異なる余極限位相}
\label{examples-section-colimit-topology}

\noindent
この例は \cite[Example 1.2, page 553]{TSH} である．
$G_n = \mathbf{Q} \times \mathbf{R}^n$，$n \geq 1$ を，加法に関する位相群で，通常の
（Euclid）位相を備えたものとみなす．$(x_0, \ldots, x_n)$ を
$(x_0, \ldots, x_n, 0)$ に写す閉埋め込み $G_n \to G_{n + 1}$ を考える．
$G = \colim G_n$ に位相
$$
U \subset G\text{ が開} \Leftrightarrow G_n \cap U\text{ が開 }\forall n
$$
を入れたものは位相群でないと主張する．

\medskip\noindent
これを見るため，集合
$$
U = \{(x_0, x_1, x_2, \ldots)\text{ で，}
|x_j| < |\cos(jx_0)| \text{ が } j > 0\}
$$
を考える．
% P12-ERR-0407: frozen overly broad integral-multiple claim; see ERRATA_CANDIDATES.jsonl.
$\pi$ は有理数でないので $jx_0$ は決して $\pi/2$ の整数倍にならないことを用いると，
$U \cap G_n$ が開であることは容易に示せる．$0 \in U$ なので，上の位相が $G$ を位相群にするならば，
$0$ の開近傍 $V \subset G$ で $V + V \subset U$ を満たすものが存在する．すると各
$j \geq 0$ に対し，$|x_j| < \epsilon_j$ ならば
$(0, \ldots, 0, x_j, 0, \ldots) \in V$ となるような $\epsilon_j > 0$ が存在する．
$V + V \subset U$ だから，$|x_0| < \epsilon_0$ および $|x_j| < \epsilon_j$ に対して
$$
(x_0, 0, \ldots, 0, x_j, 0, \ldots) \in U
$$
となるはずである．しかし，$j \epsilon_0 > \pi/2$ となるほど十分大きい $j$ をとれば，
$|\cos(jx_0)|$ が $\epsilon_j$ より小さくなるような $x_0 \in \mathbf{Q}$ を選べる．したがって
$|\cos(jx_0)| < |x_j| < \epsilon_j$ を満たす $x_j$ が存在する．この矛盾が主張を証明する．

\begin{lemma}
\label{examples-lemma-colimit-topology}
（アーベル）位相群の系 $G_1 \to G_2 \to G_3 \to \ldots$ で，位相空間の圏でとった
$\colim G_n$ と位相群の圏でとった $\colim G_n$ が異なるものが存在する．
\end{lemma}

\begin{proof}
上の議論を参照せよ．
\end{proof}
% END EXPANDED MODULE ch110_79.tex
% BEGIN EXPANDED MODULE ch110_80.tex
\section{普遍的商写像だが V 被覆でない射}
\label{examples-section-universally-submersive-not-V}

\noindent
$A$ を付値環とする．$\mathfrak p \subset A$ を，極小素イデアルでも極大イデアルでもない素イデアルとする．
（念頭に置くとよい場合は，$A$ が三つの素イデアルをもち，$\mathfrak p$ がその「中央」のものである場合である．）
アフィンスキームの射
$$
\Spec(A_\mathfrak p) \amalg \Spec(A/\mathfrak p) \longrightarrow \Spec(A)
$$
を考える．これは普遍的商写像であると主張する．これを証明するため，環準同型 $A \to B$ が与える
アフィンスキームの射 $\Spec(B) \to \Spec(A)$ をとる．このとき
$$
\Spec(B_\mathfrak p) \amalg \Spec(B/\mathfrak pB) \to \Spec(B)
$$
が商写像であることを示さなければならない．まず，これは全射である．次に，部分集合
$T \subset \Spec(B)$ で
$T_1 = \Spec(B_\mathfrak p) \cap T$ および
$T_2 = \Spec(B/\mathfrak p B) \cap T$ が閉であるものをとる．$T_1$ と $T_2$ はいずれも
$B$-代数のスペクトルなので，この $T$ はある $B$-代数のスペクトルの像である．したがって $T$ が閉であることを
示すには，$T$ が特殊化について安定であることを示せば十分である．Algebra, Lemma
\ref{algebra-lemma-image-stable-specialization-closed} を参照せよ．これを見るため，
$p \leadsto q$ を $\Spec(B)$ の点の特殊化で $p \in T$ を満たすものとする．
$A'$ を付値環とし，$\Spec(A')$ の一般点 $\eta$ が $p$ に，$\Spec(A')$ の閉点 $s$ が $q$ に写るような射
$\Spec(A') \to \Spec(B)$ をとる．Schemes, Lemma \ref{schemes-lemma-points-specialize} を参照せよ．
合成 $\gamma : \Spec(A') \to \Spec(A)$ の像は，ちょうど
$\gamma(\eta) \leadsto \xi \leadsto \gamma(s)$ を満たす点 $\xi \in \Spec(A)$ 全体である
（詳細は省略する）．$\mathfrak p \not \in \Im(\gamma)$ なら，$p, q \in \Spec(B_\mathfrak p)$ がともに
成り立つか，$p, q \in \Spec(B/\mathfrak pB)$ がともに成り立つ．この場合，$T_1$，それぞれ
$T_2$ が閉であることから，$q \in T_1$，それぞれ $q \in T_2$，したがって $q \in T$ が従う．
最後に $\mathfrak p \in \Im(\gamma)$ とし，$\mathfrak p = \gamma(r)$ と書く．このとき特殊化
$p \leadsto r$ および $r \leadsto q$ がある．この場合，
$p, r \in \Spec(B_\mathfrak p)$ および $r, q \in \Spec(B/\mathfrak pB)$ である．
% P12-ERR-0408: frozen ``fist'' typo in the English sequence; see ERRATA_CANDIDATES.jsonl.
そこでまず $r \in T_1 \subset T$，次に $r$ は $\mathfrak p$ に写るので $r \in T_2$，さらに
$q \in T_2 \subset T$ を結論し，望む結果を得る．

\medskip\noindent
一方，一元族
$$
\{\Spec(A_\mathfrak p) \amalg \Spec(A/\mathfrak p) \longrightarrow \Spec(A)\}
$$
は V 被覆でないと主張する．Topologies, Definition \ref{topologies-definition-V-covering} を参照せよ．
% P12-ERR-0409: frozen ``if it where'' conditional; see ERRATA_CANDIDATES.jsonl.
実際，これが V 被覆ならば，付値環の拡大 $A \subset B$ で，
$\Spec(B) \to \Spec(A)$ が
$\Spec(A_\mathfrak p) \amalg \Spec(A/\mathfrak p)$ を経由するものが存在する．
% P12-ERR-0410: frozen undefined A' rather than B; see ERRATA_CANDIDATES.jsonl.
これは $\Spec(A')$ が非連結であることを意味するが，不合理である．

\begin{lemma}
\label{examples-lemma-universally-submersive-not-V}
普遍的商写像であるアフィンスキームの射 $X \to Y$ で，$\{X \to Y\}$ が V 被覆でないものが存在する．
\end{lemma}

\begin{proof}
上の議論を参照せよ．
\end{proof}
% END EXPANDED MODULE ch110_80.tex
% BEGIN EXPANDED MODULE ch110_81.tex
\section{整数環のスペクトルは準コンパクトでない}
\label{examples-section-canonical}

\noindent
もちろん，この節の表題は整数環のスペクトルを位相空間として見たものを指していない．任意のスペクトルは
位相空間として準コンパクトだからである（Algebra, Lemma \ref{algebra-lemma-quasi-compact}）．そうではなく，
スキームの圏の標準位相における整数環のスペクトルと，サイトにおける準コンパクト対象の定義
（Sites, Definition \ref{sites-definition-quasi-compact}）を指している．

\medskip\noindent
素数全体の集合 $P$ 上の非単項超フィルターを $U$ とする．部分集合 $T \subset P$ に対し，
$T^c = P \setminus T$ で補集合を表す．$A \in U$ に対し，$S_A \subset \mathbf{Z}$ を
$p \in A$ によって生成される乗法的部分集合とする．
$$
\mathbf{Z}_A = S_A^{-1}\mathbf{Z}
$$
とおく．$P$ を $\Spec(\mathbf{Z})$ の閉点全体の集合とみなせば，
$\Spec(\mathbf{Z}_A) = \{(0)\} \cup A^c \subset \Spec(\mathbf{Z})$ であることに注意する．
$A, B \in U$ ならば $A \cap B \in U$ および $A \cup B \in U$ であり，
$$
\mathbf{Z}_{A \cap B} =
\mathbf{Z}_A \times_{\mathbf{Z}_{A \cup B}} \mathbf{Z}_B
$$
が成り立つ（環のファイバー積）．特に，任意の整数 $n$ と元 $A_1, \ldots, A_n \in U$ に対し，射
$$
\Spec(\mathbf{Z}_{A_1}) \amalg \ldots \amalg \Spec(\mathbf{Z}_{A_n})
\longrightarrow \Spec(\mathbf{Z})
$$
は，ある $p$ に対して $\Spec(\mathbf{Z}[1/p])$ を経由する
% P12-ERR-0411: frozen singular agreement in the English sentence; see ERRATA_CANDIDATES.jsonl.
（実際，任意の $p \in A_1 \cap \ldots \cap A_n$ が使える）．したがって，平坦射の族
$\{\Spec(\mathbf{Z}_A) \to \Spec(\mathbf{Z})\}_{A \in U}$ は共同全射であるが，その有限部分族は
どれも共同全射でない．

\medskip\noindent
$\mathbf{Z}$-加群 $M$ に対して
$$
M_A = S_A^{-1}M = M \otimes_{\mathbf{Z}} \mathbf{Z}_A
$$
とおく．主張 I：任意の $\mathbf{Z}$-加群 $M$ に対して
$$
M =
\text{Equalizer}\left(
\xymatrix{
\prod\nolimits_{A \in U} M_A \ar@<1ex>[r] \ar@<-1ex>[r] &
\prod\nolimits_{A, B \in U} M_{A \cup B}
}
\right)
$$
が成り立つ．まず $M$ が捩れ自由であると仮定する．このときすべての $A \in U$ に対して
$M_A \subset M_P$ である．したがって，示すべきことは
$$
M = \bigcap\nolimits_{A \in U} M_A\text{ inside }M_P = M \otimes \mathbf{Q}
$$
である．実際，$U$ は非単項なので，任意の素数 $p$ に対して $\{p\}^c \in U$ である．また
$M_{\{p\}^c} = M_{(p)}$ は素イデアル $(p)$ における局所化に等しい．
% P12-ERR-0412: frozen false two-prime intersection assertion; see ERRATA_CANDIDATES.jsonl.
したがって，すでに $M_{(2)} \cap M_{(3)} = M$ なので上の主張は明らかである．
次に $M$ が捩れ加群であると仮定する．このとき
$$
M = \bigoplus\nolimits_{p \in P} M[p^\infty]
$$
であり，これに対応して
$$
M_A = \bigoplus\nolimits_{p \not \in A} M[p^\infty]
$$
となる．これは $A$ に属する素数で局所化するからである．
$(x_A) \in \prod M_A$ が等化子に属するとする．
$x_p = x_{\{p\}^c} \in M[p^{\infty}]$ と書く．すると等化子の性質は
$$
x_A = (x_p)_{p \not \in A}
$$
を述べ，特に $p \not \in A$ のうち有限個を除いて $x_p$ が零であると述べる．捩れの場合の証明を
終えるには，有限個を除くすべての素数 $p$ に対して $x_p$ が零であることを示せば十分である．
そうでないとし，$\{p \in P \mid x_p \not = 0\} = T \amalg T'$ を二つの無限集合の非交和として書く．
すると $U$ は超フィルターなので，$T \not \in U$ または $T' \not \in U$ である
（実際，$T, T'$ の両方が $U$ に属すれば，$U$ は $T \cap T' = \emptyset$ を含むことになり，これは許されない）．
$T \not \in U$ とする．このとき $T = A^c$ となり，上で述べた有限性に矛盾する．
最後に $M$ を一般の加群とする．短完全列
$$
0 \to M_{tors} \to M \to M/M_{tors} \to 0
$$
と，次の大きな図式を考える：
$$
\xymatrix{
M_{tors} \ar[r] \ar[d] &
\prod\nolimits_{A \in U} M_{tors, A} \ar@<1ex>[r] \ar@<-1ex>[r] \ar[d] &
\prod\nolimits_{A, B \in U} M_{tors, A \cup B} \ar[d] \\
M \ar[r] \ar[d] &
\prod\nolimits_{A \in U} M_A \ar@<1ex>[r] \ar@<-1ex>[r] \ar[d] &
\prod\nolimits_{A, B \in U} M_{A \cup B} \ar[d] \\
M/M_{tors} \ar[r] &
\prod\nolimits_{A \in U} (M/M_{tors})_A \ar@<1ex>[r] \ar@<-1ex>[r] &
\prod\nolimits_{A, B \in U} (M/M_{tors})_{A \cup B} \\
}
$$
% P12-ERR-0413: frozen English participial sentence lacks a finite verb; see ERRATA_CANDIDATES.jsonl.
列の完全性と，捩れ加群 $M_{tors}$ および捩れ自由加群 $M/M_{tors}$ に対する結果を用いて図式追跡を
行うと，$M$ に対する主張 I が証明される．これは次の補題の現象の例を与える．

\begin{lemma}
\label{examples-lemma-non-fpqc-descent}
環 $A$ と平坦環準同型の無限族 $\{A \to A_i\}_{i \in I}$ が存在し，任意の $A$-加群 $M$ に対して
$$
M =
\text{Equalizer}\left(
\xymatrix{
\prod\nolimits_{i \in I} M \otimes_A A_i \ar@<1ex>[r] \ar@<-1ex>[r] &
\prod\nolimits_{i, j \in I} M \otimes_A A_i \otimes_A A_j
}
\right)
$$
が成り立つが，同じことが成り立つ有限部分族は存在しない．
\end{lemma}

\begin{proof}
上の議論を参照せよ．
\end{proof}

\noindent
素数全体の集合 $P$ 上の非単項超フィルター $U$ を引き続き用いる．$R$ を環とする．
$A \in U$ に対し $R_A = S_A^{-1}R = R \otimes \mathbf{Z}_A$ と書く．
主張 II：閉部分集合 $T_A \subset \Spec(R_A)$，$A \in U$ で，すべての $A, B \in U$ に対して
$$
(\Spec(R_{A \cup B}) \to \Spec(R_A))^{-1}T_A =
(\Spec(R_{A \cup B}) \to \Spec(R_B))^{-1}T_B
$$
を満たすものが与えられたとする．このとき閉部分集合 $T \subset \Spec(R)$ で，すべての
$A \in U$ に対して $T_A = (\Spec(R_A) \to \Spec(R))^{-1}(T)$ を満たすものが存在する．
$A \in U$ に対し，$I_A \subset R_A$ を $T_A$ を切り出す根基イデアルとする．貼り合わせ条件から，
すべての $A, B \in U$ に対して
$S_{A \cup B}^{-1}I_A = S_{A \cup B}^{-1}I_B$ が $R_{A \cup B}$ において成り立つ
（局所化は根基イデアルであることを保つからである）．$I' \subset R$ を
$I_P \subset R_P = R \otimes \mathbf{Q}$ に写る元全体の集合とする．すると $A \in U$ に対して
\begin{enumerate}
\item $I_A \subset I'_A = S_A^{-1}I'$，かつ
\item $M_A = I'_A/I_A$ は捩れ加群である．
\end{enumerate}
もちろん，すべての $A, B \in U$ に対して標準同一視
$S_{A \cup B}^{-1}M_A = S_{A \cup B}^{-1}M_B$ を得る．捩れ加群 $M_A$ をその
$p$-一次成分へ分解すると，$p$-冪捩れ $R$-加群 $M_p$ で
$$
M_A = \bigoplus\nolimits_{p \not \in A} M_p
$$
を満たし，しかも上の標準同一視と両立するものが存在することを，読者は容易に示せる．
$M = \bigoplus_{p \in P} M_p$ とおくと，上の標準同一視と両立する標準同型
$M_A = S_A^{-1}M$ が得られる．すると $R$-加群の標準写像
$$
I' \longrightarrow M
$$
で，すべての $A \in U$ に対して
% P12-ERR-0414: frozen I_A domain rather than I'_A; see ERRATA_CANDIDATES.jsonl.
写像 $I_A \to M_A$ を回収するものを得る．これは，上で述べたすべての両立性と，$M$ が
$\prod_{A \in U} M_A$ から $\prod_{A, B \in U} M_{A \cup B}$ への二写像の等化子であるという
先に証明した主張による．$I = \Ker(I' \to M)$ とおく．すると $I$ はイデアルであり，
$T = V(I)$ はすべての $A \in U$ に対して閉部分集合 $T_A$ を回収する閉部分集合である．
これが主張 II を証明する．

\begin{lemma}
\label{examples-lemma-Z-not-quasi-compact}
スキーム $\Spec(\mathbf{Z})$ は，スキームの圏の標準位相において準コンパクトでない．
\end{lemma}

\begin{proof}
上の記法の下で，射の族
$$
\mathcal{W} = \{\Spec(\mathbf{Z}_A) \to \Spec(\mathbf{Z})\}_{A \in U}
$$
を考える．Descent, Lemma \ref{descent-lemma-universal-effective-epimorphism} と上で証明した二つの主張により，
これは普遍的有効エピ射である．ファイバー積をもつ任意の圏において，普遍的有効エピ射は
$\mathcal{C}$ にサイトの構造を与え（容易に解決できる集合論上の問題を除く），標準位相を定める．
したがって $\mathcal{W}$ は標準位相の被覆である．一方，上で見たように，任意の有限部分族
$$
\{\Spec(\mathbf{Z}_{A_i}) \to \Spec(\mathbf{Z})\}_{i = 1, \ldots, n},\quad
n \in \mathbf{N}, A_1, \ldots, A_n \in U
$$
は，ある $p$ に対して $\Spec(\mathbf{Z}[1/p])$ を経由する．したがって，この有限族は普遍的有効
エピ射ではありえず，より一般に，普遍的有効エピ射
$\{g_j : T_j \to \Spec(\mathbf{Z})\}$ が
$\{\Spec(\mathbf{Z}_{A_i}) \to \Spec(\mathbf{Z})\}_{i = 1, \ldots, n}$ を細分することもない．
Sites, Definition \ref{sites-definition-quasi-compact} により，これは
$\Spec(\mathbf{Z})$ が標準位相において準コンパクトでないことを意味する．ここでの準コンパクト性の概念が
通常のトポス論的定義と一致することについては，Sites, Lemma \ref{sites-lemma-quasi-compact} を参照せよ．
\end{proof}
% END EXPANDED MODULE ch110_81.tex
